% =============================================================================
% Chapter 2 - Background
% =============================================================================
% Included from main.tex via % =============================================================================
% Chapter 2 - Background
% =============================================================================
% Included from main.tex via % =============================================================================
% Chapter 2 - Background
% =============================================================================
% Included from main.tex via % =============================================================================
% Chapter 2 - Background
% =============================================================================
% Included from main.tex via \include{contents/chapter2}
%
% In-chapter \autoref targets resolve once all nine sections are present.
% Cross-chapter references (Chapter 3, Chapter 4, Section 3.1.2, ...) are
% written as plain text until those chapters carry labels.
%
% Citation keys used in this chapter:
%   yan2014     [1]  Yan et al., CASME II
%   qu2016      [2]  Qu et al., CAS(ME)^2
%   liY2018     [3]  Li, Huang & Zhao, single apex frame
%   bai2021     [4]  Bai et al., magnification + pre-trained network
%   liong2019a  [5]  Liong et al., OFF-ApexNet
%   shreve2011  [6]  Shreve et al., spatio-temporal strain
%   zhaoS2021   [7]  Zhao et al., two-stage 3D CNN
%   yang2021    [8]  Yang et al., SimAM
%   xia2020a    [9]  Xia et al., STRCN
%   zhang2022  [10]  Zhang et al., SLSTT
%   see2019    [11]  See et al., MEGC 2019
% =============================================================================

\chapter{Background}
\label{ch:background}

This chapter sets out the mechanisms the rest of the thesis depends on, in the order the pipeline applies them: what a micro-expression is, what the recording supplies, how small motions are magnified and converted into a motion representation, what a neural network does, which blocks this system assembles from that foundation, how it is trained under a corpus that cannot be balanced, and how it is measured when that corpus is small.

It explains \textbf{what each mechanism is and how it works}. Chapter~3 surveys who has used each one, with what result, and what gap remains; where a fact is needed in both places it is established here and cross-referenced there. Scope is limited to what the implemented system actually does --- techniques a reader might expect in a background chapter but which this project does not use are omitted, and their absence is noted where it would otherwise be assumed.

Two constraints recur. The corpus is small enough that its size, rather than any modelling preference, forces most of the design. And several stages depart from the textbook form of the technique they implement; each departure is named where it occurs.


\section{The Phenomenon: Micro-Expressions as Involuntary, Brief, Low-Intensity Facial Movement}
\label{sec:phenomenon}

A micro-expression is a brief facial movement that leaks a felt emotion the person is actively trying to conceal~\cite{yan2014}. It is best understood by contrast with an ordinary, or macro-, expression, on three axes rather than one. Macro-expressions typically last upwards of half a second, and may run to several seconds, and can be produced at will; micro-expressions, by contrast, are ``characterized by short durations, involuntary generation and low intensity''~\cite{qu2016}. The third axis is the one most easily overlooked: a micro-expression cannot be performed on instruction. It is a symptom of concealment rather than a communicative act. That single property governs how such data can be gathered at all, and the consequences for corpus construction are taken up in Section~3.1.

\subsection{Duration: the first constraint}
\label{sec:duration}

The field's working definition fixes an upper bound on duration. Yan et al.~\cite{yan2014} state that ``the generally accepted upper limit of the duration is 1/2 s.'' Brevity is the first source of the difficulty in recognising the phenomenon automatically, and the mechanism is simple: whatever signal distinguishes one emotion from another must be carried by however many video frames fit inside that half-second window. At a modest frame rate that window may contain only a handful of frames, several of which are near-identical to their neighbours; the sequence available to a classifier is short by construction, not by any deficiency of the recording.

\subsection{Low intensity: the second constraint}
\label{sec:low-intensity}

Duration alone would still leave a tractable problem if the movement were large. It is not. Yan et al.~\cite{yan2014} observe that micro-expressions are ``usually low in intensity -- it might be so brief for the facial muscles to become fully-stretched with suppression'', with the consequence that ``because of the short duration and low intensity, it is usually imperceptible or neglected by the naked eyes''. In practice the movement of interest may change only a few grey levels over a small patch of skin --- around an eyebrow, or the corner of a mouth --- while the rest of the face stays static. This is a distinct obstacle from brevity: a longer recording does nothing to help it, because the difficulty is not how many frames are available but how far the signal in each frame sits above the noise floor of ordinary video. Duration limits how much evidence exists; intensity limits how visible that evidence is once it does.

\subsection{Onset, apex and offset}
\label{sec:onset-apex-offset}

Because a micro-expression is a movement rather than a static configuration, it is conventionally described by three temporal landmarks. The \textbf{onset} is the frame at which the face begins to depart from its neutral baseline; the \textbf{offset} is the frame at which it returns to neutral; and the \textbf{apex} is the frame at which the movement is judged most intense --- the point at which, as Ekman puts it, a ``snapshot taken at an [sic] point when the expression is at its apex can easily convey the emotion message'' (as reported by Li, Huang \& Zhao~\cite{liY2018}). All three landmarks are available in the corpus used here. This thesis, however, uses only the onset and offset frames to bound each clip passed into the pipeline; the apex frame is never read. This is a deliberate design choice rather than an oversight, and its consequences --- for what information the model can and cannot see, and for how its results relate to apex-frame methods in the literature --- are taken up in Chapters~3 and~4.

\subsection{Action units, and the label actually used}
\label{sec:action-units}

The Facial Action Coding System (FACS) describes any facial movement, however subtle, as a combination of discrete Action Units (AUs), ``an objective method for labeling facial movements in terms of component actions''~\cite{yan2014}. An AU is a description of movement, not of emotion, and the mapping between the two is not one-to-one; corpora therefore assign emotion categories using AUs together with other evidence, by a procedure that differs between datasets and is described for the corpus used here in Section~3.1.2.

It is worth being precise about how much of that machinery survives into this thesis: none of it. The work reported here uses only the resulting emotion label. The Action Units recorded for each clip are present in the label table as inherited metadata from the original annotation --- carried along because the released coding includes them --- but they are never read by any stage of the data pipeline and never presented to the model. Nothing in this thesis performs, or claims to perform, action-unit recognition; the sole target throughout is the emotion category.

How brief and how faint that movement is determines what the recording must capture, and \autoref{sec:recording-to-input} turns to the video this pipeline actually receives.


\section{From Recording to Pipeline Input}
\label{sec:recording-to-input}

Before any mechanism of interest --- magnification, flow, strain, or a trained network --- can act on a micro-expression, something has to capture it, and something has to hand it to the pipeline in a usable form. This section describes that handoff: what the recording had to provide given the constraints fixed in \autoref{sec:phenomenon}, what the corpus actually supplies, and what this project does and does not do to a clip before it enters processing.

\subsection{High-speed capture}
\label{sec:high-speed-capture}

\autoref{sec:duration} fixed a half-second upper bound on the duration of a micro-expression. That bound has an immediate consequence for recording: whatever number of frames fit inside the event is the entire evidence a classifier will ever see of it, so the frame rate at which the video was captured is not a matter of picture quality but a precondition for the signal existing in the data at all. A 300~ms movement recorded at 25 or 30~fps --- ordinary video rates --- yields only seven or eight frames. The same event recorded at 200~fps yields roughly sixty frames, enough for onset, apex and offset to be distinct, separately observable moments rather than a blur collapsed into a handful of frames.

CASME II was recorded at 200~fps for exactly this reason~\cite{yan2014}. The pipeline used in this thesis takes 200 as a configured constant (\texttt{fps: int = 200} in \texttt{CASMEIIConfig}, \texttt{Stage1\_DataPipeline/config.py}), rather than reading it from the video files themselves, since it is a fixed property of the corpus rather than something to be inferred per clip.

\subsection{What the corpus supplies, and what this project does not do}
\label{sec:corpus-supplies}

CASME II is distributed in two forms: the frames as they were recorded, and a normalised release in which each face has been cropped and registered. This pipeline reads the first. Its metadata-unification stage resolves each clip to a per-subject, per-clip folder of the original recorded frames (\texttt{Stage1\_DataPipeline/metadata\_unifier.py}), so what enters processing is the full recorded field of view at 640~$\times$~480 (Section~3.1.2), within which the face occupies only a part.

It is worth stating plainly what this means, because the alternative assumption is the natural one to make. This project performs no face detection, no landmark localisation, no geometric registration, no alignment and no cropping of its own, anywhere in the pipeline --- and it does not inherit those operations from the corpus either, because it does not read the release that carries them. The normalisation that produced that release --- 68-point Active Shape Model landmarking on each clip's first frame, registered to a neutral model face by a Local Weighted Mean transform, described in Chapter~3, Section~3.1.2 --- is the corpus authors' work, and Section~3.1.2 records that this thesis departs from it. The face reaching the network is therefore an unregistered one, at whatever position and scale the camera recorded it, and whatever tolerance to that variation the model acquires it must acquire from the data.

\subsection{Frame selection and preparation}
\label{sec:frame-preparation}

Each frame is read in greyscale and resized to 224~$\times$~224 by bilinear interpolation. The resize is applied unconditionally to the whole frame, so it both downsamples and, since 640~$\times$~480 is not square, does not preserve the recorded aspect ratio. From the frames spanning a clip's annotated onset to its annotated offset, a fixed number of frames is selected for further processing --- the scheme by which that selection is made is described in \autoref{sec:three-channel-tensor}, together with the tensor it produces.

\subsection{Which clips enter the pipeline}
\label{sec:clip-selection}

Not every row of the corpus's coding table reaches processing. Four conditions are applied to select the rows that do: the row must belong to the CASME II dataset, its coded expression type must be \texttt{micro-expression} rather than one of the other categories the table records, its frame folder must be present on disk and non-empty, and its sequence length --- the onset-to-offset frame count, counted inclusively --- must exceed two. The first two conditions restrict the pipeline to the corpus and the expression class this thesis addresses, and the third discards rows whose media the coding table lists but the filesystem does not supply. The fourth rules out clips too short to be usable: a sequence of two frames or fewer contains at most one frame-to-frame transition, which is not enough to construct a motion sequence of any length, whatever downstream representation is built from it.

Together, these properties of the recording and the corpus --- a 200~fps capture rate, an unregistered full recorded frame, a greyscale 224~$\times$~224 resize of it, and a filtered, onset-to-offset clip --- are what reaches the pipeline before any processing of interest begins. Eulerian motion magnification is the first such process applied to these frames, and \autoref{sec:motion-magnification} describes the mechanism by which it amplifies motion without ever tracking it.


\section{Motion magnification}
\label{sec:motion-magnification}

\subsection{The Eulerian idea}
\label{sec:eulerian-idea}

A moving object can be described in two ways. A \textit{Lagrangian} description follows a point on the object as it moves --- this is what optical flow and feature tracking do: locate a point in frame $t$, find its correspondence in frame $t+1$, and report the displacement. A \textit{Eulerian} description instead fixes attention on a location in the image and asks how the signal observed there changes over time. Eulerian video magnification (EVM) takes the second view: no correspondence between frames is computed. At every pixel, it treats the sequence of intensity values over time as a 1-D signal, isolates the part occurring in a chosen temporal frequency band, multiplies it by a factor $\alpha$, and adds the result back to the original sequence.

The justification for why this amplifies \textit{motion}, and not merely intensity, is a first-order argument. If a small one-dimensional image profile $f(x)$ translates by a time-varying displacement $\delta(t)$, the observed intensity at a fixed location is $I(x,t) = f(x+\delta(t))$. A first-order Taylor expansion around $x$ gives

\[
	I(x,t) \approx f(x) + \delta(t)\,\frac{\partial f(x)}{\partial x},
\]

so the temporal variation at a fixed pixel is approximately proportional to the true displacement $\delta(t)$, scaled by the local spatial gradient. Amplifying that temporal variation by $\alpha$ before adding it back is therefore, to first order, equivalent to synthesising a signal consistent with a displacement of $(1+\alpha)\delta(t)$ rather than $\delta(t)$: motion is amplified as a consequence of amplifying intensity change, with no tracking step anywhere in the process. This is the account given by Bai et al.~\cite{bai2021} (following Wu et al.).

\subsection{Spatial decomposition: the Laplacian pyramid}
\label{sec:laplacian-pyramid}

The Taylor approximation above only holds where the spatial gradient is well-behaved relative to the displacement, so in practice the frame is first decomposed into a set of spatial-frequency bands and the argument is applied band-by-band. The standard tool is the \textbf{Laplacian pyramid}. Each level is built by blurring the current image with a small low-pass kernel, downsampling by a factor of two, upsampling the result back to the original resolution, and subtracting it from the pre-downsampling image. The subtraction discards everything the blur-and-downsample step could reconstruct, leaving only the spatial detail specific to that scale --- the Laplacian band. The downsampled image feeds the next iteration, and the recursion terminates in a coarse, heavily blurred residual (a Gaussian, not a Laplacian, band). A pyramid of depth $L$ yields $L$ Laplacian bands, each isolating a different range of spatial frequencies, plus one coarsest residual that carries no band-specific detail and exists only to seed reconstruction.

\subsection{Temporal filtering}
\label{sec:temporal-filtering}

Each spatial band is treated as a stack of scalar time series, one per pixel, and passed through a \textbf{temporal band-pass filter} that keeps only the variation falling inside a chosen frequency range and discards the rest. The filtered signal at each band is multiplied by $\alpha$ and added back to that band's original values; the bands are then recombined by repeating the pyramid construction in reverse --- upsampling the coarsest residual and successively adding each (now-amplified) Laplacian band back in. The published amplitude-based method realises the band-pass step with a \textbf{Butterworth filter}, a smooth infinite-impulse-response design~\cite{bai2021}.

\subsection{The amplification trade-off}
\label{sec:amplification-tradeoff}

The magnification factor $\alpha$ trades sensitivity against artefact. A larger $\alpha$ makes a fainter in-band motion visible, but the same multiplication applies indiscriminately to everything the band-pass filter passes: sensor noise, compression artefacts, and any other real motion whose temporal frequency happens to fall inside the chosen band are amplified exactly as the motion of interest is. The method has no way to distinguish a wanted signal from an unwanted one occupying the same frequency range --- the filter only knows how fast a signal oscillates, not what it is. Raising $\alpha$ therefore does not monotonically improve the result; past some point amplified noise and induced displacement artefacts dominate whatever gain in visibility was sought. This tension is intrinsic to the method rather than a defect of any one implementation, and it recurs, with numbers, in Chapter~3.

\subsection{Departures from the textbook method in this implementation}
\label{sec:evm-departures}

The magnifier used in this thesis (\texttt{Stage1\_DataPipeline/evm\_magnifier.py}) implements the mechanism above with a genuine Laplacian pyramid --- a $[1,4,6,4,1]/16$ blur kernel, downsample-by-two, upsample, subtract --- followed by a coarsest Gaussian residual, matching \autoref{sec:laplacian-pyramid} exactly. It departs from amplitude-based magnification as originally formulated in four respects; the third is a property of that original method rather than of the summary given by Bai et al.~\cite{bai2021}.

First, the temporal band-pass is an \textbf{ideal filter}: a hard binary mask applied directly to the discrete Fourier transform of each pixel's time series (via \texttt{scipy.fftpack}), rather than a Butterworth (IIR) filter. Frequencies inside $[\omega_{\mathrm{low}}, \omega_{\mathrm{high}}]$ pass with a gain of exactly one and everything else is zeroed, with no smooth roll-off at the band edges.

Second, every Laplacian band is amplified by the \textbf{same scalar $\alpha$}. Bai et al.~\cite{bai2021} state the amplified reconstruction with a \textit{frequency-dependent} factor $\alpha_k$ --- one per band, reduced at higher spatial frequencies to limit noise amplification --- and no such damping is applied here. The coarsest residual is never filtered or amplified --- it only seeds reconstruction.

Third, the implementation works directly on the greyscale intensity tensor, so there is no chromatic attenuation step.

Fourth, and a property of pipeline ordering rather than of the magnifier itself: the calling code (\texttt{tensor\_pipeline\_manager.py}) resamples each clip to a fixed 33 frames \textit{before} invoking the magnifier, and passes the corpus's nominal recording rate (200~fps) as the filter's frame-rate parameter regardless of the resampled sequence's true frame spacing. The consequence --- a temporal filter whose frequency axis no longer matches the signal it is filtering --- is analysed in Section~3.5.7, not here.

The operating point this study runs at is given in Section~4.2.5; whatever the amplification factor, band and pyramid depth, the amplified output is clipped to $[0,255]$ and quantised back to 8-bit before being handed to the downstream flow-and-strain extractor.

What that extractor does with the amplified sequence --- and why motion, rather than appearance, is what the network is given --- is the subject of \autoref{sec:video-to-motion}.


\section{From Video to Motion: Optical Flow and Optical Strain}
\label{sec:video-to-motion}

\subsection{Why motion, not pixels}
\label{sec:why-motion}

A micro-expression is a movement, so the natural quantity to hand a classifier is a description of movement rather than of appearance. Raw pixel intensity carries the signal of interest --- the same faint change described in \autoref{sec:low-intensity} --- embedded in everything else the camera also recorded: who the subject is, how the scene is lit, the fixed geometry of a particular face. None of that is discarded by presenting a network with frames directly, and on a corpus this narrow in identity and illumination a model has every incentive to fit the wrong part of the image. A displacement field is different in kind: it records how each point moved, not what colour it is, so identity and illumination do not enter it at all. That principle motivates the rest of this section; the comparative case for it is made in Section~3.3.1.

\subsection{Optical flow: the brightness-constancy constraint}
\label{sec:optical-flow}

Optical flow estimates, for every pixel, the apparent two-dimensional displacement between two frames. The estimate rests on the brightness-constancy assumption: a point on a moving surface is taken to have the same intensity in the next frame as in the current one, so that any change in a pixel's grey level is attributed to something having moved past that location, never to the thing itself changing colour or shade~\cite{liong2019a}.

Writing $I_t(x,y)$ for the image intensity at position $(x,y)$ and time $t$, and supposing the point there moves to $(x+\delta x,\, y+\delta y)$ by time $t+1$, brightness constancy states

\[
	I_t(x,y) = I_{t+1}(x+\delta x,\, y+\delta y), \qquad \delta x = u_t\,\delta t,\;\; \delta y = v_t\,\delta t,
\]

where $u_t(x,y)$ and $v_t(x,y)$ are the horizontal and vertical components of the flow at that point. Expanding the right-hand side by a first-order Taylor series and substituting into the constancy equation, the intensity terms cancel and division by $\delta t$ leaves the \textbf{optical flow constraint equation}:

\[
	u_t(x,y)\,\frac{\partial I}{\partial x} + v_t(x,y)\,\frac{\partial I}{\partial y} + \frac{\partial I}{\partial t} = 0.
\]

This single scalar equation holds at every pixel but contains two unknowns, $u_t$ and $v_t$: the image gradients are measurable directly from the frames, while the displacement is not determined by them alone. This under-determinacy is the \textbf{aperture problem} --- a local patch of image constrains the component of motion perpendicular to an edge or gradient, but says nothing about motion parallel to it. One equation per pixel cannot fix two unknowns per pixel, so every practical estimator closes the system by adding a further assumption relating the flow at neighbouring pixels, typically that the field varies smoothly across the image. That added assumption is what distinguishes one flow algorithm from another; this thesis does not adjudicate between them (Section~3.3.3) and simply adopts one, described in \autoref{sec:three-channel-tensor}. The output, for a pair of frames, is the pair of scalar fields $u(x,y)$ and $v(x,y)$ --- the horizontal and vertical displacement at every pixel --- forming the first two channels of the motion representation used throughout this thesis.

\subsection{Optical strain: the finite strain tensor}
\label{sec:optical-strain}

Flow answers where a point moved; it does not by itself distinguish a patch of skin that moved together with its neighbours from one that stretched or sheared relative to them. Optical strain makes that distinction by taking the spatial derivative of the flow field rather than the field itself.

Writing the displacement at a point as $\mathbf{u} = [u, v]^{\mathsf T}$, the \textbf{finite strain tensor} is defined as the symmetric part of the displacement gradient~\cite{shreve2011}:

\[
	\varepsilon = \tfrac{1}{2}\left[\nabla\mathbf{u} + (\nabla\mathbf{u})^{\mathsf T}\right] = \begin{bmatrix} \varepsilon_{xx} & \varepsilon_{xy} \\ \varepsilon_{yx} & \varepsilon_{yy} \end{bmatrix}, \qquad \varepsilon_{xx} = \frac{\partial u}{\partial x},\;\; \varepsilon_{yy} = \frac{\partial v}{\partial y},\;\; \varepsilon_{xy} = \varepsilon_{yx} = \frac{1}{2}\left(\frac{\partial u}{\partial y} + \frac{\partial v}{\partial x}\right).
\]

$\varepsilon_{xx}$ and $\varepsilon_{yy}$ are the \textbf{normal strain} components, describing stretching or compression along each axis; $\varepsilon_{xy}$ is the \textbf{shear strain} component, describing the change of angle between two initially perpendicular directions on the surface. The tensor is reduced to a single scalar per pixel. Shreve et al.~\cite{shreve2011} sum the component values to obtain a strain magnitude; the reduction implemented here is instead a root sum of squares, with the shear term counted once,

\[
	\varepsilon_{\text{mag}} = \sqrt{\varepsilon_{xx}^{2} + \varepsilon_{yy}^{2} + \varepsilon_{xy}^{2}},
\]

a form that diverges from the convention used across the reviewed corpus. Section~3.4.3 records that divergence and what it costs.

The property that makes this useful, rather than a second copy of the flow field, follows from its definition as a \textit{gradient}. If a region of the face translates rigidly --- every pixel moving by the same $u$ and the same $v$, as a small head movement produces --- then $u$ and $v$ are locally constant there, and their spatial derivatives, hence every component of $\varepsilon$, are zero: a rigid translation contributes nothing to the strain field, however large it is. A localised muscle contraction, by contrast, moves neighbouring points of skin by different amounts, so the displacement field has a non-zero gradient exactly where the deformation occurs, and $\varepsilon_{\text{mag}}$ is large there. Strain is thus sensitive to non-rigid deformation specifically, insensitive by construction to whatever rigid motion the flow field also contains.

\subsection{Assembling the three-channel tensor}
\label{sec:three-channel-tensor}

For each clip, the frames spanning onset to offset are uniformly resampled to a fixed sequence of 33 frames, and dense optical flow is computed between every adjacent pair using OpenCV's implementation of Farneb\"ack's method~\cite{zhaoS2021} --- \texttt{cv2.calcOpticalFlowFarneback}, with pyramid scale 0.5, three pyramid levels, window size 15, three iterations, a polynomial neighbourhood of size 5 and a polynomial smoothing of 1.2, applied to 8-bit greyscale frames. This yields \textbf{32 flow fields} per clip, each a $(u, v)$ pair over the 224~$\times$~224 face region. Strain is computed from each field's spatial gradients using \texttt{np.gradient} (central differences at unit pixel spacing, over the two spatial axes only; no temporal gradient is used), giving $\varepsilon_{xx}$, $\varepsilon_{yy}$, $\varepsilon_{xy}$ and hence $\varepsilon_{\text{mag}}$ for each of the 32 pairs.

The three quantities --- $u$, $v$ and $\varepsilon_{\text{mag}}$ --- are stacked as three channels, giving a tensor of shape $(3, 32, 224, 224)$ in channel order $[u, v, \text{strain}]$. When Eulerian video magnification is applied (\autoref{sec:motion-magnification}), flow and strain are computed on the magnified frames rather than the originals; the construction is otherwise unchanged.

Normalisation is applied in two separate stages, and both remain active. At extraction time, each of the three channels is independently min--max normalised to $[0, 1]$ across the whole clip, so that a channel's own range --- which differs enormously between a flow component and a strain magnitude --- does not by itself determine its influence. Separately, at training load time (\texttt{Ablation\_Study/dataset.py}, with \texttt{normalize=True}), each channel is z-scored per sample over its $T \times H \times W$ extent. Neither stage substitutes for the other: the first is a per-clip rescaling done once, during data preparation; the second is a per-sample standardisation applied every time a clip is loaded, centring and scaling each channel to zero mean and unit variance.

How this tensor is built is now settled; what consumes it is not. \autoref{sec:nn-fundamentals} begins with what a neural network is, before \autoref{sec:building-blocks} reaches the blocks themselves.


\section{Neural Network Fundamentals}
\label{sec:nn-fundamentals}

Everything in this thesis described as ``trained'' is, underneath, a neural network: a function built from simple, repeated computations whose numerical parameters are set from data rather than specified by hand. This section builds that machinery from its smallest unit upward, only as far as is needed for the architecture-specific mechanisms of \autoref{sec:building-blocks}, the training regime of \autoref{sec:scarcity-skew} and the evaluation apparatus of \autoref{sec:evaluation} to be legible when they introduce anything that depends on it.

\subsection{Units, layers and the forward pass}
\label{sec:units-layers}

The basic computational unit takes a vector of inputs $x$, forms a weighted sum of them plus an offset term, and passes the result through a fixed non-linear function $\phi$:
\[
	a = \phi(w^\top x + b).
\]
The entries of $w$ are the unit's \textbf{weights}, one per input; $b$ is its \textbf{bias}. Weights and biases together are the unit's \textbf{parameters} --- the numbers a training procedure is free to adjust --- and $a$, the unit's output, is its \textbf{activation}. A \textbf{layer} is a bank of such units applied to the same input vector in parallel, each with its own weights and bias, producing a vector of activations rather than a single scalar. A network is a sequence of layers, each layer's output vector becoming the next layer's input; evaluating the network on an input by applying each layer in turn, from the first to the last, is the \textbf{forward pass}, and its final output is the network's prediction. The number of layers is the network's \textbf{depth}; the number of units in a given layer is that layer's \textbf{width}. Both are choices made when the architecture is specified, not something training decides.

\subsection{Non-linearity}
\label{sec:non-linearity}

If $\phi$ were the identity function, every layer would compute only an affine transformation of its input, and the composition of any number of affine transformations is itself a single affine transformation. A network built entirely this way would therefore compute no more than a network with one layer could, however many layers it stacked --- depth would buy nothing. The non-linear function $\phi$, applied after the weighted sum at every unit, is what stops layers from collapsing into one another; it is the reason depth is a meaningful design axis at all rather than a redundant one. A common choice, used throughout the components described in this thesis, is the rectified linear unit, $\text{ReLU}(x) = \max(0, x)$: it passes positive values through unchanged and clips negative ones to zero. Its appeal is mechanical rather than subtle --- it is cheap to compute, its derivative is either zero or one almost everywhere, and unlike activation functions that saturate on both sides, a unit with a positive input keeps gradient flowing through it at full strength.

\subsection{Learning by gradient descent}
\label{sec:gradient-descent}

A network's forward pass is fixed once its parameters are fixed, so ``training'' means searching for parameter values under which the forward pass produces good predictions. A \textbf{loss function} quantifies how far a prediction sits from its true label, returning a single number that is small when they agree and large when they do not; the specific loss used in this thesis, and why, is a matter for \autoref{sec:scarcity-skew}, but the role a loss plays is a property of every trained network and belongs here. Training proceeds by adjusting parameters to make the loss smaller, and the mechanism that makes that adjustment tractable is \textbf{backpropagation}: because the forward pass is a composition of layers, the chain rule from calculus lets the gradient of the loss with respect to every parameter in the network --- however deep --- be computed by working backward from the output layer to the input, reusing intermediate results rather than recomputing each parameter's gradient from scratch. Once every parameter's gradient is known, an \textbf{update rule} moves each parameter a small step against its gradient, since the gradient points in the direction the loss increases fastest and moving the opposite way decreases it. The size of that step is the \textbf{learning rate}. Parameters are not updated from the whole training set at once in a single step; instead the data is divided into \textbf{batches}, and one update is made per batch, while one full pass through every batch in the training set is one \textbf{epoch}, of which training runs many in succession. Throughout, a network has two distinct modes of use: in training mode its parameters are being adjusted from labelled examples by the procedure just described, while in evaluation mode its parameters are held fixed and it is only asked to produce predictions, with no update following from how correct they are.

\subsection{Convolution as a learned filter}
\label{sec:convolution}

A layer of the kind described in \autoref{sec:units-layers} connects every input to every unit, which is wasteful for an image: a pattern worth detecting --- an edge, a corner, a small patch of texture --- can appear anywhere in the frame, and a fully connected layer would need to learn a separate copy of the detector for every position it might appear at. A convolutional layer instead learns one small filter, a \textbf{kernel}, consisting of a grid of weights (a $3\times3$ or $5\times5$ window is typical) plus a single bias, and slides that same kernel across every position of the input, at each position computing a weighted sum of the pixels currently under it. Because the identical kernel is reused at every position, the layer is said to exhibit \textbf{weight sharing}: it looks for the same pattern everywhere in the image and costs only as many parameters as the kernel itself holds, regardless of how large the input is. The result of sliding one kernel across the whole input is a single \textbf{feature map}, one value per position, recording how strongly that kernel's pattern was present there. A convolutional layer ordinarily learns several kernels at once, each producing its own feature map; stacked together these form the layer's output \textbf{channels}, so a layer with $C_{\text{out}}$ kernels produces $C_{\text{out}}$ channels from whatever number of input channels, $C_{\text{in}}$, it was given (three for an ordinary colour image, one for a greyscale image, or more once earlier layers have produced feature maps of their own). Because each kernel spans every input channel, not just one, the number of learned parameters in such a layer, for a kernel of height $k_H$ and width $k_W$, is
\[
	C_{\text{out}} \times \bigl(C_{\text{in}} \times k_H \times k_W + 1\bigr),
\]
the $+1$ accounting for one bias per output channel --- a formula this thesis returns to repeatedly when comparing the cost of alternative designs. Two further settings control how the kernel is applied. \textbf{Stride} is the number of positions the kernel moves between one application and the next; a stride greater than one skips positions and shrinks the output accordingly. \textbf{Padding} adds a border, typically of zeros, around the input before the kernel is applied, which controls how much the output shrinks relative to the input and lets a layer preserve spatial size exactly if that is wanted. Because a single kernel only ever looks at a small window, one convolutional layer alone only detects small, local patterns; stacking convolutional layers lets a unit in a later layer depend, indirectly, on a much larger region of the original input, since each layer's output already summarises a neighbourhood of the layer before it. This growing \textbf{receptive field} is why early convolutional layers tend to respond to local structure --- edges, small textures --- while later ones respond to larger, more composite structure built from what the earlier layers detected.

\subsection{Pooling and spatial reduction}
\label{sec:pooling}

A \textbf{pooling} operation reduces the spatial size of a feature map without learning any weights of its own: it slides a small window across the map, exactly as a kernel does, but instead of a weighted sum it reports a fixed summary of the values it covers --- most commonly the maximum, giving max pooling. Pooling deliberately discards information: a $2\times2$ max-pool halves both height and width, keeping only the strongest response in each small neighbourhood and throwing the rest away. What is bought in exchange is a lower computational cost for every layer that follows, since they now operate on a smaller map, and a larger receptive field per unit of depth, since a later layer's kernel now covers a coarser, wider-reaching grid of the original input for the same kernel size. Pooling is used where the exact position of a detected feature matters less than its coarser location and the network can no longer afford to carry full spatial detail forward.

\subsection{Parameters, capacity and small data}
\label{sec:capacity}

The total count of a network's parameters --- every weight and bias, across every layer --- determines its \textbf{capacity}: loosely, how large and how intricate a set of functions it is able to represent. A higher-capacity network can fit more complicated relationships between input and label, but capacity is measured against the data available to constrain it, not in isolation. With a training corpus of only a few hundred or few thousand labelled examples, a network with far more parameters than examples has enough freedom to fit each training example almost exactly, including whatever is peculiar to that specific example --- sensor noise, an unrepresentative pose, an idiosyncrasy of one subject --- rather than the pattern that generalises across examples. This failure mode is \textbf{overfitting}: training loss keeps falling because the network is, in effect, memorising its training set, while performance on examples it did not train on stops improving or gets worse. The property actually wanted is \textbf{generalisation} --- accurate prediction on data the network never saw during training --- and it cannot be read off the training loss at all, since a memorising network can drive that loss arbitrarily low. Detecting whether generalisation is actually happening requires a \textbf{held-out set}: a portion of the labelled data set aside from the start, never used to compute a gradient or update a parameter, on which the network's predictions are checked only after training. A gap between good performance on the training data and poor performance on the held-out set is the signature of overfitting; parameter count, relative to the number of available training examples, is what makes that gap likely in the first place. This is precisely the constraint that governs the size and construction of every architecture described next, in \autoref{sec:building-blocks}.


\section{Network Building Blocks}
\label{sec:building-blocks}

Chapter~3 argues why each of the following components was chosen and what its measured contribution was. This section stays one level below that argument, setting out the mechanism each component computes.

\subsection{Three-dimensional convolution and kernel shape}
\label{sec:conv3d}

A \texttt{Conv3d} layer generalises image convolution to a five-dimensional tensor of shape $(B, C, D, H, W)$ --- batch, channel, and three spatial-like axes, here depth $D$ (time), height $H$ and width $W$. A kernel of shape $(k_D, k_H, k_W)$ extends the sliding, weight-sharing computation of \autoref{sec:convolution} across all three non-channel axes at once, with stride and padding again governing the output's size along each.

The detail that matters most for what follows is the kernel's temporal extent, $k_D$. If $k_D > 1$, the filter reaches across multiple frames at every application, so a single output value is a genuine function of more than one time step --- a spatio-temporal filter in the literal sense. If $k_D = 1$, the filter at time step $t$ only ever reads input at time step $t$; no output value depends on any other frame. A stack of layers with kernel shape $(1, k_H, k_W)$ is therefore a two-dimensional convolution applied identically and independently to every frame, merely expressed using \texttt{Conv3d} operations so the tensor never leaves its $(C, D, H, W)$ layout. This is the case the model in this thesis uses throughout its convolutional stem: every kernel is shaped $(1, 3, 3)$, with padding $(0, 1, 1)$ preserving the spatial extent while the temporal axis is neither padded nor touched. The temporal dimension is carried through the stem unchanged in length; only $H$ and $W$ are affected, first by the convolutions' own padding and then by a spatial-only max-pool of shape $(1, 2, 2)$ that halves both.

\subsection{Normalisation and regularisation}
\label{sec:normalisation}

\texttt{BatchNorm3d} normalises each channel independently: for channel $c$ it computes a mean and variance over every other axis at once --- the batch dimension and all spatial-temporal positions $(D, H, W)$ --- and rescales,
\[
	\hat{x} = \frac{x - \mu_c}{\sqrt{\sigma_c^2 + \epsilon}}, \qquad y = \gamma_c \hat{x} + \beta_c,
\]
where $\gamma_c$ and $\beta_c$ are learned per-channel scale and shift parameters. Training uses batch statistics; inference uses a running estimate accumulated during training, so the layer behaves exactly as its two-dimensional counterpart does, only pooled over one further axis.

\texttt{Dropout3d} differs from ordinary \texttt{Dropout} in what unit it zeroes. Ordinary dropout masks individual scalar activations independently at random, which is a weak regulariser on a convolutional feature map, since neighbouring positions in $H$, $W$ and $D$ are highly correlated and the network can read a missing value off an adjacent one. \texttt{Dropout3d} instead zeroes an entire channel --- every $(D, H, W)$ position of it, for a given sample --- with the given probability, removing a whole feature detector rather than a single correlated pixel.

\texttt{LayerNorm} normalises across the feature dimension of a single position for a single sample, independent of the batch, which is why it is the default choice inside transformer blocks rather than \texttt{BatchNorm}. Here it appears inside every transformer encoder layer and once more at the head of the classifier, which applies it before a dropout and the final linear projection to the output classes.

\subsection{Parameter-free attention: SimAM}
\label{sec:simam}

SimAM~\cite{yang2021} re-weights a feature map without learning any new parameters, by scoring each neuron on how distinctive it is from its surrounding context and using that score directly as a multiplicative gate. Yang et al. define an energy $e^{*}$ that is \textit{lower} for a more distinctive neuron, so it is the \textbf{reciprocal} of that energy that serves as the gate. For a neuron with activation $x$ in a channel whose mean and variance are $\mu$ and $v$ (estimated by pooling over that channel's own spatial extent), that reciprocal has the closed form
\[
	\frac{1}{e^{*}(x)} = \frac{(x-\mu)^2}{4(v+\lambda)} + 0.5,
\]
where $\lambda$ is the module's single hyperparameter. The refined output is then $x \cdot \sigma\!\left(1/e^{*}(x)\right)$, a monotonic sigmoid gate applied element-wise, the sigmoid present only to bound the gate rather than to reshape it. What makes this mechanism unusual among attention modules is exactly what the formula shows: every quantity it needs --- the mean, the variance, the resulting score --- is read off the feature map itself, at inference as much as at training time, so no weight matrix, bias, or reduction layer is introduced anywhere. The gate is a fixed function of the activations it is applied to, not a learned one.

The implementation in this thesis follows that formula exactly, with $\lambda = 10^{-4}$, but estimates $\mu$ and $v$ by pooling over the entire spatio-temporal volume $(D, H, W)$ of a stream's feature map rather than a single frame's spatial extent alone, so a neuron's distinctiveness is judged against the whole clip.

\subsection{Self-attention}
\label{sec:self-attention}

Self-attention relates every position in a sequence to every other position by learned linear projection rather than a fixed spatial neighbourhood. Each input vector is projected three ways, into a query $Q$, a key $K$ and a value $V$; attention weights are the scaled dot product of queries against keys,
\[
	\text{Attention}(Q,K,V) = \text{softmax}\!\left(\frac{QK^\top}{\sqrt{d_k}}\right)V,
\]
with the $\sqrt{d_k}$ scaling present to stop the dot products, and hence the softmax, from saturating as the key dimension $d_k$ grows. Multi-head attention runs several such projections in parallel, each into a smaller subspace, concatenating their outputs before a final linear mixing layer, so different heads can attend to different relationships within the same sequence.

Because this is a weighted sum over the whole sequence with weights derived purely from content, it is invariant to the order of the input positions: permuting the sequence permutes the output identically, with nothing in the mechanism encoding \textit{where} a token sits. A positional encoding is added to each input vector to break that symmetry. Two schemes are common. A sinusoidal encoding assigns each position a fixed vector built from sine and cosine functions at geometrically varying frequencies across the embedding dimensions --- deterministic, requiring no training, and evaluable at sequence lengths never seen during training. A learned encoding instead assigns each position index its own trainable embedding, which can fit the data more closely but cannot extrapolate beyond the longest position trained on. The encoder in this thesis uses the fixed sinusoidal form, held as a constant buffer rather than a parameter.

A second placement choice concerns normalisation relative to each sublayer. Post-norm --- the original arrangement --- applies \texttt{LayerNorm} after the residual addition that follows each sublayer; pre-norm applies it beforehand, to the sublayer's input, so the residual path carries an un-normalised signal straight through the stack, which is generally reported to train more stably since gradients reach earlier layers unimpeded. This is the placement used here. Beyond what the encoder layer itself contributes, no additional residual or skip connection is introduced anywhere else in the architecture.

\subsection{The fallbacks}
\label{sec:fallbacks}

Each learned component here can be switched off, and each switch has a specific, near parameter-free replacement rather than simply removing that stage. With the convolutional stem off, every frame's motion channels are instead reduced by adaptive average pooling to a fixed $4\times4$ spatial grid regardless of input resolution, flattened per frame and passed through a single linear layer into the model's working dimension. With the transformer off, the sequence of per-frame feature vectors is instead collapsed by a plain mean over the time axis, with no parameters and no representation of frame order at all.

Both fallbacks add as little capacity as possible: the pooling operations introduce no parameters, and the one linear layer needed to match dimensions in the no-CNN case is negligible next to what it replaces. That near-absence of added capacity is what makes each fallback a fair control, isolating what the learned component contributes without also changing how much the network is capable of fitting.

Isolating a component fairly also depends on how the network is trained, and \autoref{sec:scarcity-skew} sets out the training regime held constant across every configuration.


\section{Learning Under Scarcity and Skew}
\label{sec:scarcity-skew}

This section sets out the training machinery --- the loss function, the sampling scheme, regularisation, and the optimiser --- at the level of what each mechanism computes and how it is configured here. Chapter~3 (Section~3.9) reviews the literature motivating these choices, gives the corpus's class counts, and reports the failure mode that follows from combining two of them carelessly; that narrative is not repeated below. What follows is the mechanism and the configuration actually used.

\subsection{What skew does to cross-entropy}
\label{sec:skew-cross-entropy}

Standard cross-entropy sums a per-example term $-\log p_t$, where $p_t$ is the model's predicted probability for the true class, weighting every example equally. Under a skewed class distribution that equal weighting is not equal in effect: a class's total gradient contribution is its example count times the per-example gradient, so a majority class dominates the sum purely by outnumbering the rest. A model facing this loss has a cheap route to a lower value --- predict close to the prior on ambiguous cases --- because that is correct often enough on the majority class to outweigh being wrong on every minority example. Loss can fall steadily while the model learns comparatively little about classes it rarely sees. The remainder of this section responds to that problem.

\subsection{Focal loss: re-weighting by difficulty, not by frequency}
\label{sec:focal-loss}

Focal loss (Lin et al., 2017) inserts a modulating factor in front of the log term:

\[
	\text{FL}(p_t) = -\alpha_t (1-p_t)^{\gamma} \log p_t.
\]

Zhao et al.~\cite{zhaoS2021} apply this formulation to micro-expression recognition and state the mechanism plainly: $\gamma$ is ``the balance factor for loss'' and $\alpha_t$ ``the weight balance factor for samples''. The distinction matters for what follows. The factor $(1-p_t)^{\gamma}$ shrinks toward zero as the model becomes confident and correct on an example, re-weighting by how \textit{difficult} it currently finds that example, irrespective of its class --- an easy majority-class example contributes almost nothing once learned; a hard example, of any class, keeps contributing. Nothing in that factor looks at how many examples of a class exist. Class frequency enters only through the separate $\alpha_t$ term, an explicit per-class multiplier ordinarily set near inverse class frequency.

The implementation here (\texttt{Ablation\_Study/losses.py}) follows this exactly: it computes \texttt{log\_softmax}, gathers $\log p_t$ for the true class, exponentiates to recover $p_t$, and forms \texttt{focal\_weight = (1 - p\_t) ** gamma}. Label smoothing (\autoref{sec:label-smoothing}) is applied first, blending the negative log-likelihood term with a uniform term over classes, so the focal weight multiplies the \textit{smoothed} loss rather than the raw NLL. The optional $\alpha$ vector, where supplied, multiplies the result afterwards, indexed by each example's true class. The value of $\gamma$ used in this study follows Zhao et al.'s reported setting; Section~4.5.1 gives it, and the label-smoothing value, together with the rest of the training configuration.

\subsection{Balanced sampling}
\label{sec:balanced-sampling}

The second lever acts on the data a batch is drawn from rather than on the loss computed over it. A \texttt{WeightedRandomSampler} assigns each training example a weight equal to the inverse of its class's count, computed over the training split of the current fold only. Sampling then proceeds with replacement, drawing \texttt{num\_samples = len(train)} indices per epoch, so each class is presented with roughly equal probability regardless of its true rarity. This sampler is attached to the training loader alone; the validation loader uses no sampler and iterates with \texttt{shuffle=False}, so validation scores reflect the split's true distribution.

\subsection{Why combining both is a choice, not a default}
\label{sec:combining-both}

Focal loss's $\alpha$ term and the balanced sampler both correct for class frequency, at different points in the pipeline, so running both at once corrects for the same imbalance twice. A pipeline that applies both is therefore making a choice, not accepting a default, and the choice has to be resolved somewhere explicit. Section~4.5.2 states how this study resolves it and which of the two corrections is left active. Section~3.9.7 (cross-referenced rather than repeated here) explains why this rule exists and what happens when it is violated; the fact to carry forward is simply that the two mechanisms are not additive by default in this codebase.

\subsection{Label smoothing}
\label{sec:label-smoothing}

Label smoothing replaces the one-hot target with a softened one, mixing a fraction $\epsilon$ of uniform probability mass across all classes into the true-class target. Inside the focal-loss implementation this is a weighted combination of the ordinary negative log-likelihood term and a term averaging $-\log p_c$ over every class $c$; Section~4.5.1 gives the value of $\epsilon$ used here. The effect is to discourage the model from driving the true-class logit arbitrarily high --- regularisation against over-confidence, applied uniformly regardless of class, and distinct from a correction for imbalance. AdamW, mixed precision and gradient clipping have no dedicated treatment in this project's \texttt{docs/} corpus and are used here as standard practice.

\subsection{Augmentation under scarcity}
\label{sec:augmentation}

With few labelled clips per class, augmentation acts before the loss or the sampler is reached at all, and is applied to the training split only. Xia et al.~\cite{xia2020a} motivate this for micro-expression data directly: ``temporal data augmentation strategies as well as a balanced loss are jointly used for our deep network'' to address ``limited and imbalanced training samples''. Two transformations are implemented (\texttt{Ablation\_Study/dataset.py}), but only one of them acts. A random temporal crop of the input window (centred at evaluation) fires only where a clip's stored length exceeds the length the model expects; extraction stores exactly the number of frames the model is configured to consume (\autoref{sec:three-channel-tensor}), so that condition never holds and the crop is inert throughout this study. What remains is a horizontal flip applied with probability one half. Because the input is optical flow rather than raw pixels (\autoref{sec:video-to-motion}), the flip cannot be a plain mirror --- reversing the scene horizontally reverses the sign of horizontal displacement --- so the flip is paired with a negation of the flow tensor's $u$ (horizontal) channel. Flipping without negating $u$ would hand the network motion vectors pointing the wrong way for the geometry shown, a quiet corruption motion-tensor augmentation must specifically guard against.

\subsection{Optimisation}
\label{sec:optimisation}

Training uses AdamW, with decoupled weight decay applied uniformly to every parameter --- no separate no-decay group for biases or normalisation terms. The schedule combines a short linear warmup with cosine annealing: \texttt{LinearLR} ramps the rate from a fraction of target up to full over the warmup epochs, after which \texttt{CosineAnnealingLR} anneals it to a small floor over the remainder; the two are joined by \texttt{SequentialLR} and stepped once per epoch. Section~4.5.4 gives the learning rate, weight decay, warmup length and epoch count used here. Cosine annealing is adopted elsewhere in the micro-expression literature on similar small-data grounds --- Zhang et al.~\cite{zhang2022} use it in a spatio-temporal transformer for this same task.

Gradient norms are clipped after the AMP scaler's gradients are unscaled, via \texttt{clip\_grad\_norm\_}, guarding against the occasional large update a small, skewed dataset can produce. Training runs under mixed precision (\texttt{torch.amp.GradScaler}, \texttt{autocast}), enabled where CUDA is available and not disabled by flag. In the same unscaled branch, and therefore only while clipping is on, every parameter's gradient is checked for non-finite values before the optimiser step; if any is found, the step is skipped entirely, gradients are zeroed, and a warning logged, rather than letting a corrupted update reach the weights. Model selection keeps the checkpoint with the highest validation macro F1, not accuracy, consistent with the metric argument of Section~3.9.5 that accuracy is a poor guide under skew; ties favour the later epoch. There is no early stopping: every configuration trains its full epoch budget regardless of when its best epoch occurred.

One category of technique is absent by design: quantisation, pruning, and knowledge distillation are used nowhere in this project. Efficiency is measured rather than engineered here --- Section~3.6 reports the compute cost of the architectural choices actually made.

Measuring rather than engineering a property requires an evaluation apparatus equal to the task; \autoref{sec:evaluation} builds the one this corpus demands, from the confusion matrix upward.


\section{Evaluating on a Small Corpus}
\label{sec:evaluation}

Chapter~3 argues why leave-one-subject-out evaluation and pooled macro F1 suit this corpus, and what those choices cost here (Sections~3.1.6 and~3.1.8). This section gives the mechanism underneath: the confusion matrix and per-class precision, recall, F1; averaging scores versus pooling counts; the easily conflated distinction between averaging across folds and within them; what cross-validation buys and what leave-one-subject-out protects against; and, to close, an audit of what this thesis's own measurement apparatus does and does not supply.

\subsection{The confusion matrix, and precision, recall, F1 for one class}
\label{sec:confusion-matrix}

For a single class $c$ treated as positive against all others, every prediction falls into one of four counts: true positives $TP_c$ (correctly predicted $c$), false positives $FP_c$ (predicted $c$ but truly something else), false negatives $FN_c$ (truly $c$ but predicted otherwise), and true negatives. Arranging all classes this way at once gives the confusion matrix: entry $(i,j)$ counts clips whose true label is $i$ and predicted label is $j$, so the diagonal holds correct predictions and every off-diagonal cell names a specific error --- which true class is mistaken for which other.

From those counts, precision and recall isolate two failure modes:
\[
	\text{Precision}_c = \frac{TP_c}{TP_c + FP_c}, \qquad \text{Recall}_c = \frac{TP_c}{TP_c + FN_c}.
\]
Precision penalises false alarms: it falls whenever the classifier calls $c$ on a clip that is not $c$, regardless of how many true $c$ instances it also catches. Recall penalises misses: it falls whenever a genuine $c$ instance is assigned elsewhere, regardless of how many false alarms accompany the hits. Neither failure mode is visible from the other, which is why F1, the harmonic mean,
\[
	F1_c = \frac{2 \cdot \text{Precision}_c \cdot \text{Recall}_c}{\text{Precision}_c + \text{Recall}_c} = \frac{2\,TP_c}{2\,TP_c + FP_c + FN_c}
\]
is the usual single-number summary: being harmonic, it stays low whenever either input is low, so a model cannot buy a good F1 by trading recall for precision.

\subsection{Macro versus micro averaging}
\label{sec:macro-micro}

Combining per-class F1 scores into one number can be done two ways. Macro averaging computes $F1_c$ per class and takes an unweighted mean, $\text{Macro-F1} = \tfrac1C\sum_c F1_c$: every class counts equally regardless of size, so ignoring a rare class is penalised as heavily as ignoring a common one. Micro averaging instead pools the raw counts across classes first --- summing every class's $TP$, $FP$ and $FN$ --- and computes one F1 from those totals.

For single-label multi-class classification, where every clip gets exactly one predicted label, this pooling has a consequence worth stating plainly: every misclassification is simultaneously a false positive for the class it was wrongly assigned to and a false negative for its true class, so summed $FP$ and summed $FN$ equal each other and the total error count, and pooled micro-F1 reduces algebraically to $TP_{\text{total}}/N$ --- plain accuracy. Micro-F1 and accuracy are not two different metrics here; they are the same number computed two ways. The contrast that actually matters is therefore not ``macro versus micro'' but \textbf{macro-F1 versus accuracy}: an unweighted per-class average against a count a majority class can dominate.

\subsection{Pooled versus per-fold averaging}
\label{sec:pooled-vs-fold}

A further averaging choice sits underneath macro-F1 once evaluation is split into folds, and it governs how every result in this thesis is reported. One option accumulates the confusion counts --- $TP_c$, $FP_c$, $FN_c$ --- across every fold first, then computes one set of per-class F1 scores, and one macro-F1, from that pooled total. The other computes a complete macro-F1 inside each fold from that fold's own predictions, then averages the per-fold values. See et al.~\cite{see2019} define the Unweighted F1 (UF1) of the MEGC 2019 protocol the first way: the macro-averaged F1 obtained by accumulating true positives, false positives and false negatives over all folds of a leave-one-subject-out run before computing and averaging per-class F1.

These are estimators of two different quantities, not two ways of writing the same one: pooled macro-F1 describes this class's F1 across the whole evaluation, once every held-out prediction sits in one confusion matrix together, while per-fold-averaged macro-F1 describes the typical per-fold macro-F1 --- a quantity that depends on how classes happen to fall within each fold, not only on the classifier. When a fold's held-out set does not contain every class, its per-class F1 for the missing class is degenerate, and averaging such folds in with fully-populated ones changes what the number reflects for reasons unrelated to classifier quality; on a small, unevenly distributed corpus this is a live possibility, and fold composition can bound the per-fold estimator below what the pooled one would report. Sections~3.1.6 and~3.1.8 work through how the fold structure of this study does exactly that; which one a reported number is becomes apparent only once its computation is stated.

\subsection{Cross-validation and subject-disjointness}
\label{sec:cross-validation}

$k$-fold cross-validation partitions the data into $k$ disjoint subsets, trains $k$ models --- each on $k-1$ folds --- and evaluates each on the one fold it never saw, so every example is scored exactly once by a model that never trained on it. Leave-one-subject-out (LOSO) is the case where $k$ equals the number of subjects, and each fold's held-out set is precisely one subject's clips, every other subject's clips forming that fold's training set.

The property this buys is subject-disjointness: no fold's training and validation sets ever share a subject. Without it, a model evaluated on a clip from a subject it has already trained on can succeed by recognising that person's face or recording conditions rather than the expression class itself --- an identity leak inflating the score without the model having learned anything transferable to an unseen person. Enforcing subject-disjointness in every fold closes that channel.

\subsection{What an evaluation apparatus of this kind can fail to provide}
\label{sec:apparatus-gaps}

Three shortcomings belong to an evaluation apparatus rather than to any model it measures, and each matters more at this corpus size than it would at a larger one. They are set out here as properties to look for; Section~4.6 reports where this study's apparatus stands on each.

The first is \textbf{which aggregate a pipeline actually stores}. \autoref{sec:pooled-vs-fold} distinguishes pooling the confusion matrices before computing a macro average from averaging per-fold macro scores, and shows the two answer different questions. A pipeline may compute one, the other, or both under the same field name, and the quantity a reader wants is not necessarily the one written to disk.

The second is \textbf{whether model selection is blind to the data it is scored against}. A nested design keeps three sets apart: one to fit parameters, a second to choose between candidate models or epochs, and a third, untouched until the end, to score the chosen one. Collapsing the second and third --- selecting a checkpoint on the same data that then reports the result --- leaves the reported figure not blind to what it is checked against, and biases it optimistically by an amount that cannot be recovered after the fact.

The third is \textbf{whether enough is retained to support inference}. Aggregate metrics answer what a configuration scored; they cannot answer whether two configurations differ by more than chance. A paired test between configurations needs the per-clip predictions of each, so an apparatus that persists only aggregates forecloses that question whatever its scores say.


\section{Conclusion}
\label{sec:ch2-conclusion}

This chapter has set out the mechanisms on which the rest of the thesis depends, in the order the pipeline applies them. \autoref{sec:phenomenon} defined the micro-expression by the three properties that make it hard to recognise automatically --- it is brief, it is faint, and it cannot be produced on instruction --- and fixed the vocabulary of onset, apex and offset, noting that only the emotion label, not the action-unit coding, is used here. \autoref{sec:recording-to-input} established what actually reaches the pipeline: a 200~fps recording, the corpus's original frames rather than its registered release, and a greyscale frame at a fixed resolution, with no detection, registration or cropping performed by this project. \autoref{sec:motion-magnification} and \autoref{sec:video-to-motion} covered the two transformations applied to those frames --- Eulerian magnification of small motions, and the conversion of the magnified sequence into dense optical flow and optical strain.

\autoref{sec:nn-fundamentals} then built the machinery of a neural network from a single unit up to the convolutional layer, ending with the relationship between parameter count and corpus size that constrains every architectural choice this thesis makes. \autoref{sec:building-blocks} described the specific blocks assembled on that foundation, together with the parameter-free replacement each one falls back to when switched off. \autoref{sec:scarcity-skew} covered training under a class distribution that cannot be balanced by collecting more data, and \autoref{sec:evaluation} the evaluation apparatus needed when a corpus is small enough that the choice of averaging changes the conclusion.

Two threads run through all of it. The first is that the size of the corpus, rather than any modelling preference, forces most of the design: it dictates how few parameters the network can carry, why corrections for class skew are necessary, and why the evaluation protocol has to be chosen with care. The second is that several stages depart from the textbook or reference form of the technique they implement, and each departure has been named at the point it occurs rather than left for a reader to discover.

What this chapter has deliberately not done is say who else has used any of these techniques, with what result, or what remains unresolved. That is the work of Chapter~3, which reviews the literature for each component in turn and identifies the gap this thesis addresses.

%
% In-chapter \autoref targets resolve once all nine sections are present.
% Cross-chapter references (Chapter 3, Chapter 4, Section 3.1.2, ...) are
% written as plain text until those chapters carry labels.
%
% Citation keys used in this chapter:
%   yan2014     [1]  Yan et al., CASME II
%   qu2016      [2]  Qu et al., CAS(ME)^2
%   liY2018     [3]  Li, Huang & Zhao, single apex frame
%   bai2021     [4]  Bai et al., magnification + pre-trained network
%   liong2019a  [5]  Liong et al., OFF-ApexNet
%   shreve2011  [6]  Shreve et al., spatio-temporal strain
%   zhaoS2021   [7]  Zhao et al., two-stage 3D CNN
%   yang2021    [8]  Yang et al., SimAM
%   xia2020a    [9]  Xia et al., STRCN
%   zhang2022  [10]  Zhang et al., SLSTT
%   see2019    [11]  See et al., MEGC 2019
% =============================================================================

\chapter{Background}
\label{ch:background}

This chapter sets out the mechanisms the rest of the thesis depends on, in the order the pipeline applies them: what a micro-expression is, what the recording supplies, how small motions are magnified and converted into a motion representation, what a neural network does, which blocks this system assembles from that foundation, how it is trained under a corpus that cannot be balanced, and how it is measured when that corpus is small.

It explains \textbf{what each mechanism is and how it works}. Chapter~3 surveys who has used each one, with what result, and what gap remains; where a fact is needed in both places it is established here and cross-referenced there. Scope is limited to what the implemented system actually does --- techniques a reader might expect in a background chapter but which this project does not use are omitted, and their absence is noted where it would otherwise be assumed.

Two constraints recur. The corpus is small enough that its size, rather than any modelling preference, forces most of the design. And several stages depart from the textbook form of the technique they implement; each departure is named where it occurs.


\section{The Phenomenon: Micro-Expressions as Involuntary, Brief, Low-Intensity Facial Movement}
\label{sec:phenomenon}

A micro-expression is a brief facial movement that leaks a felt emotion the person is actively trying to conceal~\cite{yan2014}. It is best understood by contrast with an ordinary, or macro-, expression, on three axes rather than one. Macro-expressions typically last upwards of half a second, and may run to several seconds, and can be produced at will; micro-expressions, by contrast, are ``characterized by short durations, involuntary generation and low intensity''~\cite{qu2016}. The third axis is the one most easily overlooked: a micro-expression cannot be performed on instruction. It is a symptom of concealment rather than a communicative act. That single property governs how such data can be gathered at all, and the consequences for corpus construction are taken up in Section~3.1.

\subsection{Duration: the first constraint}
\label{sec:duration}

The field's working definition fixes an upper bound on duration. Yan et al.~\cite{yan2014} state that ``the generally accepted upper limit of the duration is 1/2 s.'' Brevity is the first source of the difficulty in recognising the phenomenon automatically, and the mechanism is simple: whatever signal distinguishes one emotion from another must be carried by however many video frames fit inside that half-second window. At a modest frame rate that window may contain only a handful of frames, several of which are near-identical to their neighbours; the sequence available to a classifier is short by construction, not by any deficiency of the recording.

\subsection{Low intensity: the second constraint}
\label{sec:low-intensity}

Duration alone would still leave a tractable problem if the movement were large. It is not. Yan et al.~\cite{yan2014} observe that micro-expressions are ``usually low in intensity -- it might be so brief for the facial muscles to become fully-stretched with suppression'', with the consequence that ``because of the short duration and low intensity, it is usually imperceptible or neglected by the naked eyes''. In practice the movement of interest may change only a few grey levels over a small patch of skin --- around an eyebrow, or the corner of a mouth --- while the rest of the face stays static. This is a distinct obstacle from brevity: a longer recording does nothing to help it, because the difficulty is not how many frames are available but how far the signal in each frame sits above the noise floor of ordinary video. Duration limits how much evidence exists; intensity limits how visible that evidence is once it does.

\subsection{Onset, apex and offset}
\label{sec:onset-apex-offset}

Because a micro-expression is a movement rather than a static configuration, it is conventionally described by three temporal landmarks. The \textbf{onset} is the frame at which the face begins to depart from its neutral baseline; the \textbf{offset} is the frame at which it returns to neutral; and the \textbf{apex} is the frame at which the movement is judged most intense --- the point at which, as Ekman puts it, a ``snapshot taken at an [sic] point when the expression is at its apex can easily convey the emotion message'' (as reported by Li, Huang \& Zhao~\cite{liY2018}). All three landmarks are available in the corpus used here. This thesis, however, uses only the onset and offset frames to bound each clip passed into the pipeline; the apex frame is never read. This is a deliberate design choice rather than an oversight, and its consequences --- for what information the model can and cannot see, and for how its results relate to apex-frame methods in the literature --- are taken up in Chapters~3 and~4.

\subsection{Action units, and the label actually used}
\label{sec:action-units}

The Facial Action Coding System (FACS) describes any facial movement, however subtle, as a combination of discrete Action Units (AUs), ``an objective method for labeling facial movements in terms of component actions''~\cite{yan2014}. An AU is a description of movement, not of emotion, and the mapping between the two is not one-to-one; corpora therefore assign emotion categories using AUs together with other evidence, by a procedure that differs between datasets and is described for the corpus used here in Section~3.1.2.

It is worth being precise about how much of that machinery survives into this thesis: none of it. The work reported here uses only the resulting emotion label. The Action Units recorded for each clip are present in the label table as inherited metadata from the original annotation --- carried along because the released coding includes them --- but they are never read by any stage of the data pipeline and never presented to the model. Nothing in this thesis performs, or claims to perform, action-unit recognition; the sole target throughout is the emotion category.

How brief and how faint that movement is determines what the recording must capture, and \autoref{sec:recording-to-input} turns to the video this pipeline actually receives.


\section{From Recording to Pipeline Input}
\label{sec:recording-to-input}

Before any mechanism of interest --- magnification, flow, strain, or a trained network --- can act on a micro-expression, something has to capture it, and something has to hand it to the pipeline in a usable form. This section describes that handoff: what the recording had to provide given the constraints fixed in \autoref{sec:phenomenon}, what the corpus actually supplies, and what this project does and does not do to a clip before it enters processing.

\subsection{High-speed capture}
\label{sec:high-speed-capture}

\autoref{sec:duration} fixed a half-second upper bound on the duration of a micro-expression. That bound has an immediate consequence for recording: whatever number of frames fit inside the event is the entire evidence a classifier will ever see of it, so the frame rate at which the video was captured is not a matter of picture quality but a precondition for the signal existing in the data at all. A 300~ms movement recorded at 25 or 30~fps --- ordinary video rates --- yields only seven or eight frames. The same event recorded at 200~fps yields roughly sixty frames, enough for onset, apex and offset to be distinct, separately observable moments rather than a blur collapsed into a handful of frames.

CASME II was recorded at 200~fps for exactly this reason~\cite{yan2014}. The pipeline used in this thesis takes 200 as a configured constant (\texttt{fps: int = 200} in \texttt{CASMEIIConfig}, \texttt{Stage1\_DataPipeline/config.py}), rather than reading it from the video files themselves, since it is a fixed property of the corpus rather than something to be inferred per clip.

\subsection{What the corpus supplies, and what this project does not do}
\label{sec:corpus-supplies}

CASME II is distributed in two forms: the frames as they were recorded, and a normalised release in which each face has been cropped and registered. This pipeline reads the first. Its metadata-unification stage resolves each clip to a per-subject, per-clip folder of the original recorded frames (\texttt{Stage1\_DataPipeline/metadata\_unifier.py}), so what enters processing is the full recorded field of view at 640~$\times$~480 (Section~3.1.2), within which the face occupies only a part.

It is worth stating plainly what this means, because the alternative assumption is the natural one to make. This project performs no face detection, no landmark localisation, no geometric registration, no alignment and no cropping of its own, anywhere in the pipeline --- and it does not inherit those operations from the corpus either, because it does not read the release that carries them. The normalisation that produced that release --- 68-point Active Shape Model landmarking on each clip's first frame, registered to a neutral model face by a Local Weighted Mean transform, described in Chapter~3, Section~3.1.2 --- is the corpus authors' work, and Section~3.1.2 records that this thesis departs from it. The face reaching the network is therefore an unregistered one, at whatever position and scale the camera recorded it, and whatever tolerance to that variation the model acquires it must acquire from the data.

\subsection{Frame selection and preparation}
\label{sec:frame-preparation}

Each frame is read in greyscale and resized to 224~$\times$~224 by bilinear interpolation. The resize is applied unconditionally to the whole frame, so it both downsamples and, since 640~$\times$~480 is not square, does not preserve the recorded aspect ratio. From the frames spanning a clip's annotated onset to its annotated offset, a fixed number of frames is selected for further processing --- the scheme by which that selection is made is described in \autoref{sec:three-channel-tensor}, together with the tensor it produces.

\subsection{Which clips enter the pipeline}
\label{sec:clip-selection}

Not every row of the corpus's coding table reaches processing. Four conditions are applied to select the rows that do: the row must belong to the CASME II dataset, its coded expression type must be \texttt{micro-expression} rather than one of the other categories the table records, its frame folder must be present on disk and non-empty, and its sequence length --- the onset-to-offset frame count, counted inclusively --- must exceed two. The first two conditions restrict the pipeline to the corpus and the expression class this thesis addresses, and the third discards rows whose media the coding table lists but the filesystem does not supply. The fourth rules out clips too short to be usable: a sequence of two frames or fewer contains at most one frame-to-frame transition, which is not enough to construct a motion sequence of any length, whatever downstream representation is built from it.

Together, these properties of the recording and the corpus --- a 200~fps capture rate, an unregistered full recorded frame, a greyscale 224~$\times$~224 resize of it, and a filtered, onset-to-offset clip --- are what reaches the pipeline before any processing of interest begins. Eulerian motion magnification is the first such process applied to these frames, and \autoref{sec:motion-magnification} describes the mechanism by which it amplifies motion without ever tracking it.


\section{Motion magnification}
\label{sec:motion-magnification}

\subsection{The Eulerian idea}
\label{sec:eulerian-idea}

A moving object can be described in two ways. A \textit{Lagrangian} description follows a point on the object as it moves --- this is what optical flow and feature tracking do: locate a point in frame $t$, find its correspondence in frame $t+1$, and report the displacement. A \textit{Eulerian} description instead fixes attention on a location in the image and asks how the signal observed there changes over time. Eulerian video magnification (EVM) takes the second view: no correspondence between frames is computed. At every pixel, it treats the sequence of intensity values over time as a 1-D signal, isolates the part occurring in a chosen temporal frequency band, multiplies it by a factor $\alpha$, and adds the result back to the original sequence.

The justification for why this amplifies \textit{motion}, and not merely intensity, is a first-order argument. If a small one-dimensional image profile $f(x)$ translates by a time-varying displacement $\delta(t)$, the observed intensity at a fixed location is $I(x,t) = f(x+\delta(t))$. A first-order Taylor expansion around $x$ gives

\[
	I(x,t) \approx f(x) + \delta(t)\,\frac{\partial f(x)}{\partial x},
\]

so the temporal variation at a fixed pixel is approximately proportional to the true displacement $\delta(t)$, scaled by the local spatial gradient. Amplifying that temporal variation by $\alpha$ before adding it back is therefore, to first order, equivalent to synthesising a signal consistent with a displacement of $(1+\alpha)\delta(t)$ rather than $\delta(t)$: motion is amplified as a consequence of amplifying intensity change, with no tracking step anywhere in the process. This is the account given by Bai et al.~\cite{bai2021} (following Wu et al.).

\subsection{Spatial decomposition: the Laplacian pyramid}
\label{sec:laplacian-pyramid}

The Taylor approximation above only holds where the spatial gradient is well-behaved relative to the displacement, so in practice the frame is first decomposed into a set of spatial-frequency bands and the argument is applied band-by-band. The standard tool is the \textbf{Laplacian pyramid}. Each level is built by blurring the current image with a small low-pass kernel, downsampling by a factor of two, upsampling the result back to the original resolution, and subtracting it from the pre-downsampling image. The subtraction discards everything the blur-and-downsample step could reconstruct, leaving only the spatial detail specific to that scale --- the Laplacian band. The downsampled image feeds the next iteration, and the recursion terminates in a coarse, heavily blurred residual (a Gaussian, not a Laplacian, band). A pyramid of depth $L$ yields $L$ Laplacian bands, each isolating a different range of spatial frequencies, plus one coarsest residual that carries no band-specific detail and exists only to seed reconstruction.

\subsection{Temporal filtering}
\label{sec:temporal-filtering}

Each spatial band is treated as a stack of scalar time series, one per pixel, and passed through a \textbf{temporal band-pass filter} that keeps only the variation falling inside a chosen frequency range and discards the rest. The filtered signal at each band is multiplied by $\alpha$ and added back to that band's original values; the bands are then recombined by repeating the pyramid construction in reverse --- upsampling the coarsest residual and successively adding each (now-amplified) Laplacian band back in. The published amplitude-based method realises the band-pass step with a \textbf{Butterworth filter}, a smooth infinite-impulse-response design~\cite{bai2021}.

\subsection{The amplification trade-off}
\label{sec:amplification-tradeoff}

The magnification factor $\alpha$ trades sensitivity against artefact. A larger $\alpha$ makes a fainter in-band motion visible, but the same multiplication applies indiscriminately to everything the band-pass filter passes: sensor noise, compression artefacts, and any other real motion whose temporal frequency happens to fall inside the chosen band are amplified exactly as the motion of interest is. The method has no way to distinguish a wanted signal from an unwanted one occupying the same frequency range --- the filter only knows how fast a signal oscillates, not what it is. Raising $\alpha$ therefore does not monotonically improve the result; past some point amplified noise and induced displacement artefacts dominate whatever gain in visibility was sought. This tension is intrinsic to the method rather than a defect of any one implementation, and it recurs, with numbers, in Chapter~3.

\subsection{Departures from the textbook method in this implementation}
\label{sec:evm-departures}

The magnifier used in this thesis (\texttt{Stage1\_DataPipeline/evm\_magnifier.py}) implements the mechanism above with a genuine Laplacian pyramid --- a $[1,4,6,4,1]/16$ blur kernel, downsample-by-two, upsample, subtract --- followed by a coarsest Gaussian residual, matching \autoref{sec:laplacian-pyramid} exactly. It departs from amplitude-based magnification as originally formulated in four respects; the third is a property of that original method rather than of the summary given by Bai et al.~\cite{bai2021}.

First, the temporal band-pass is an \textbf{ideal filter}: a hard binary mask applied directly to the discrete Fourier transform of each pixel's time series (via \texttt{scipy.fftpack}), rather than a Butterworth (IIR) filter. Frequencies inside $[\omega_{\mathrm{low}}, \omega_{\mathrm{high}}]$ pass with a gain of exactly one and everything else is zeroed, with no smooth roll-off at the band edges.

Second, every Laplacian band is amplified by the \textbf{same scalar $\alpha$}. Bai et al.~\cite{bai2021} state the amplified reconstruction with a \textit{frequency-dependent} factor $\alpha_k$ --- one per band, reduced at higher spatial frequencies to limit noise amplification --- and no such damping is applied here. The coarsest residual is never filtered or amplified --- it only seeds reconstruction.

Third, the implementation works directly on the greyscale intensity tensor, so there is no chromatic attenuation step.

Fourth, and a property of pipeline ordering rather than of the magnifier itself: the calling code (\texttt{tensor\_pipeline\_manager.py}) resamples each clip to a fixed 33 frames \textit{before} invoking the magnifier, and passes the corpus's nominal recording rate (200~fps) as the filter's frame-rate parameter regardless of the resampled sequence's true frame spacing. The consequence --- a temporal filter whose frequency axis no longer matches the signal it is filtering --- is analysed in Section~3.5.7, not here.

The operating point this study runs at is given in Section~4.2.5; whatever the amplification factor, band and pyramid depth, the amplified output is clipped to $[0,255]$ and quantised back to 8-bit before being handed to the downstream flow-and-strain extractor.

What that extractor does with the amplified sequence --- and why motion, rather than appearance, is what the network is given --- is the subject of \autoref{sec:video-to-motion}.


\section{From Video to Motion: Optical Flow and Optical Strain}
\label{sec:video-to-motion}

\subsection{Why motion, not pixels}
\label{sec:why-motion}

A micro-expression is a movement, so the natural quantity to hand a classifier is a description of movement rather than of appearance. Raw pixel intensity carries the signal of interest --- the same faint change described in \autoref{sec:low-intensity} --- embedded in everything else the camera also recorded: who the subject is, how the scene is lit, the fixed geometry of a particular face. None of that is discarded by presenting a network with frames directly, and on a corpus this narrow in identity and illumination a model has every incentive to fit the wrong part of the image. A displacement field is different in kind: it records how each point moved, not what colour it is, so identity and illumination do not enter it at all. That principle motivates the rest of this section; the comparative case for it is made in Section~3.3.1.

\subsection{Optical flow: the brightness-constancy constraint}
\label{sec:optical-flow}

Optical flow estimates, for every pixel, the apparent two-dimensional displacement between two frames. The estimate rests on the brightness-constancy assumption: a point on a moving surface is taken to have the same intensity in the next frame as in the current one, so that any change in a pixel's grey level is attributed to something having moved past that location, never to the thing itself changing colour or shade~\cite{liong2019a}.

Writing $I_t(x,y)$ for the image intensity at position $(x,y)$ and time $t$, and supposing the point there moves to $(x+\delta x,\, y+\delta y)$ by time $t+1$, brightness constancy states

\[
	I_t(x,y) = I_{t+1}(x+\delta x,\, y+\delta y), \qquad \delta x = u_t\,\delta t,\;\; \delta y = v_t\,\delta t,
\]

where $u_t(x,y)$ and $v_t(x,y)$ are the horizontal and vertical components of the flow at that point. Expanding the right-hand side by a first-order Taylor series and substituting into the constancy equation, the intensity terms cancel and division by $\delta t$ leaves the \textbf{optical flow constraint equation}:

\[
	u_t(x,y)\,\frac{\partial I}{\partial x} + v_t(x,y)\,\frac{\partial I}{\partial y} + \frac{\partial I}{\partial t} = 0.
\]

This single scalar equation holds at every pixel but contains two unknowns, $u_t$ and $v_t$: the image gradients are measurable directly from the frames, while the displacement is not determined by them alone. This under-determinacy is the \textbf{aperture problem} --- a local patch of image constrains the component of motion perpendicular to an edge or gradient, but says nothing about motion parallel to it. One equation per pixel cannot fix two unknowns per pixel, so every practical estimator closes the system by adding a further assumption relating the flow at neighbouring pixels, typically that the field varies smoothly across the image. That added assumption is what distinguishes one flow algorithm from another; this thesis does not adjudicate between them (Section~3.3.3) and simply adopts one, described in \autoref{sec:three-channel-tensor}. The output, for a pair of frames, is the pair of scalar fields $u(x,y)$ and $v(x,y)$ --- the horizontal and vertical displacement at every pixel --- forming the first two channels of the motion representation used throughout this thesis.

\subsection{Optical strain: the finite strain tensor}
\label{sec:optical-strain}

Flow answers where a point moved; it does not by itself distinguish a patch of skin that moved together with its neighbours from one that stretched or sheared relative to them. Optical strain makes that distinction by taking the spatial derivative of the flow field rather than the field itself.

Writing the displacement at a point as $\mathbf{u} = [u, v]^{\mathsf T}$, the \textbf{finite strain tensor} is defined as the symmetric part of the displacement gradient~\cite{shreve2011}:

\[
	\varepsilon = \tfrac{1}{2}\left[\nabla\mathbf{u} + (\nabla\mathbf{u})^{\mathsf T}\right] = \begin{bmatrix} \varepsilon_{xx} & \varepsilon_{xy} \\ \varepsilon_{yx} & \varepsilon_{yy} \end{bmatrix}, \qquad \varepsilon_{xx} = \frac{\partial u}{\partial x},\;\; \varepsilon_{yy} = \frac{\partial v}{\partial y},\;\; \varepsilon_{xy} = \varepsilon_{yx} = \frac{1}{2}\left(\frac{\partial u}{\partial y} + \frac{\partial v}{\partial x}\right).
\]

$\varepsilon_{xx}$ and $\varepsilon_{yy}$ are the \textbf{normal strain} components, describing stretching or compression along each axis; $\varepsilon_{xy}$ is the \textbf{shear strain} component, describing the change of angle between two initially perpendicular directions on the surface. The tensor is reduced to a single scalar per pixel. Shreve et al.~\cite{shreve2011} sum the component values to obtain a strain magnitude; the reduction implemented here is instead a root sum of squares, with the shear term counted once,

\[
	\varepsilon_{\text{mag}} = \sqrt{\varepsilon_{xx}^{2} + \varepsilon_{yy}^{2} + \varepsilon_{xy}^{2}},
\]

a form that diverges from the convention used across the reviewed corpus. Section~3.4.3 records that divergence and what it costs.

The property that makes this useful, rather than a second copy of the flow field, follows from its definition as a \textit{gradient}. If a region of the face translates rigidly --- every pixel moving by the same $u$ and the same $v$, as a small head movement produces --- then $u$ and $v$ are locally constant there, and their spatial derivatives, hence every component of $\varepsilon$, are zero: a rigid translation contributes nothing to the strain field, however large it is. A localised muscle contraction, by contrast, moves neighbouring points of skin by different amounts, so the displacement field has a non-zero gradient exactly where the deformation occurs, and $\varepsilon_{\text{mag}}$ is large there. Strain is thus sensitive to non-rigid deformation specifically, insensitive by construction to whatever rigid motion the flow field also contains.

\subsection{Assembling the three-channel tensor}
\label{sec:three-channel-tensor}

For each clip, the frames spanning onset to offset are uniformly resampled to a fixed sequence of 33 frames, and dense optical flow is computed between every adjacent pair using OpenCV's implementation of Farneb\"ack's method~\cite{zhaoS2021} --- \texttt{cv2.calcOpticalFlowFarneback}, with pyramid scale 0.5, three pyramid levels, window size 15, three iterations, a polynomial neighbourhood of size 5 and a polynomial smoothing of 1.2, applied to 8-bit greyscale frames. This yields \textbf{32 flow fields} per clip, each a $(u, v)$ pair over the 224~$\times$~224 face region. Strain is computed from each field's spatial gradients using \texttt{np.gradient} (central differences at unit pixel spacing, over the two spatial axes only; no temporal gradient is used), giving $\varepsilon_{xx}$, $\varepsilon_{yy}$, $\varepsilon_{xy}$ and hence $\varepsilon_{\text{mag}}$ for each of the 32 pairs.

The three quantities --- $u$, $v$ and $\varepsilon_{\text{mag}}$ --- are stacked as three channels, giving a tensor of shape $(3, 32, 224, 224)$ in channel order $[u, v, \text{strain}]$. When Eulerian video magnification is applied (\autoref{sec:motion-magnification}), flow and strain are computed on the magnified frames rather than the originals; the construction is otherwise unchanged.

Normalisation is applied in two separate stages, and both remain active. At extraction time, each of the three channels is independently min--max normalised to $[0, 1]$ across the whole clip, so that a channel's own range --- which differs enormously between a flow component and a strain magnitude --- does not by itself determine its influence. Separately, at training load time (\texttt{Ablation\_Study/dataset.py}, with \texttt{normalize=True}), each channel is z-scored per sample over its $T \times H \times W$ extent. Neither stage substitutes for the other: the first is a per-clip rescaling done once, during data preparation; the second is a per-sample standardisation applied every time a clip is loaded, centring and scaling each channel to zero mean and unit variance.

How this tensor is built is now settled; what consumes it is not. \autoref{sec:nn-fundamentals} begins with what a neural network is, before \autoref{sec:building-blocks} reaches the blocks themselves.


\section{Neural Network Fundamentals}
\label{sec:nn-fundamentals}

Everything in this thesis described as ``trained'' is, underneath, a neural network: a function built from simple, repeated computations whose numerical parameters are set from data rather than specified by hand. This section builds that machinery from its smallest unit upward, only as far as is needed for the architecture-specific mechanisms of \autoref{sec:building-blocks}, the training regime of \autoref{sec:scarcity-skew} and the evaluation apparatus of \autoref{sec:evaluation} to be legible when they introduce anything that depends on it.

\subsection{Units, layers and the forward pass}
\label{sec:units-layers}

The basic computational unit takes a vector of inputs $x$, forms a weighted sum of them plus an offset term, and passes the result through a fixed non-linear function $\phi$:
\[
	a = \phi(w^\top x + b).
\]
The entries of $w$ are the unit's \textbf{weights}, one per input; $b$ is its \textbf{bias}. Weights and biases together are the unit's \textbf{parameters} --- the numbers a training procedure is free to adjust --- and $a$, the unit's output, is its \textbf{activation}. A \textbf{layer} is a bank of such units applied to the same input vector in parallel, each with its own weights and bias, producing a vector of activations rather than a single scalar. A network is a sequence of layers, each layer's output vector becoming the next layer's input; evaluating the network on an input by applying each layer in turn, from the first to the last, is the \textbf{forward pass}, and its final output is the network's prediction. The number of layers is the network's \textbf{depth}; the number of units in a given layer is that layer's \textbf{width}. Both are choices made when the architecture is specified, not something training decides.

\subsection{Non-linearity}
\label{sec:non-linearity}

If $\phi$ were the identity function, every layer would compute only an affine transformation of its input, and the composition of any number of affine transformations is itself a single affine transformation. A network built entirely this way would therefore compute no more than a network with one layer could, however many layers it stacked --- depth would buy nothing. The non-linear function $\phi$, applied after the weighted sum at every unit, is what stops layers from collapsing into one another; it is the reason depth is a meaningful design axis at all rather than a redundant one. A common choice, used throughout the components described in this thesis, is the rectified linear unit, $\text{ReLU}(x) = \max(0, x)$: it passes positive values through unchanged and clips negative ones to zero. Its appeal is mechanical rather than subtle --- it is cheap to compute, its derivative is either zero or one almost everywhere, and unlike activation functions that saturate on both sides, a unit with a positive input keeps gradient flowing through it at full strength.

\subsection{Learning by gradient descent}
\label{sec:gradient-descent}

A network's forward pass is fixed once its parameters are fixed, so ``training'' means searching for parameter values under which the forward pass produces good predictions. A \textbf{loss function} quantifies how far a prediction sits from its true label, returning a single number that is small when they agree and large when they do not; the specific loss used in this thesis, and why, is a matter for \autoref{sec:scarcity-skew}, but the role a loss plays is a property of every trained network and belongs here. Training proceeds by adjusting parameters to make the loss smaller, and the mechanism that makes that adjustment tractable is \textbf{backpropagation}: because the forward pass is a composition of layers, the chain rule from calculus lets the gradient of the loss with respect to every parameter in the network --- however deep --- be computed by working backward from the output layer to the input, reusing intermediate results rather than recomputing each parameter's gradient from scratch. Once every parameter's gradient is known, an \textbf{update rule} moves each parameter a small step against its gradient, since the gradient points in the direction the loss increases fastest and moving the opposite way decreases it. The size of that step is the \textbf{learning rate}. Parameters are not updated from the whole training set at once in a single step; instead the data is divided into \textbf{batches}, and one update is made per batch, while one full pass through every batch in the training set is one \textbf{epoch}, of which training runs many in succession. Throughout, a network has two distinct modes of use: in training mode its parameters are being adjusted from labelled examples by the procedure just described, while in evaluation mode its parameters are held fixed and it is only asked to produce predictions, with no update following from how correct they are.

\subsection{Convolution as a learned filter}
\label{sec:convolution}

A layer of the kind described in \autoref{sec:units-layers} connects every input to every unit, which is wasteful for an image: a pattern worth detecting --- an edge, a corner, a small patch of texture --- can appear anywhere in the frame, and a fully connected layer would need to learn a separate copy of the detector for every position it might appear at. A convolutional layer instead learns one small filter, a \textbf{kernel}, consisting of a grid of weights (a $3\times3$ or $5\times5$ window is typical) plus a single bias, and slides that same kernel across every position of the input, at each position computing a weighted sum of the pixels currently under it. Because the identical kernel is reused at every position, the layer is said to exhibit \textbf{weight sharing}: it looks for the same pattern everywhere in the image and costs only as many parameters as the kernel itself holds, regardless of how large the input is. The result of sliding one kernel across the whole input is a single \textbf{feature map}, one value per position, recording how strongly that kernel's pattern was present there. A convolutional layer ordinarily learns several kernels at once, each producing its own feature map; stacked together these form the layer's output \textbf{channels}, so a layer with $C_{\text{out}}$ kernels produces $C_{\text{out}}$ channels from whatever number of input channels, $C_{\text{in}}$, it was given (three for an ordinary colour image, one for a greyscale image, or more once earlier layers have produced feature maps of their own). Because each kernel spans every input channel, not just one, the number of learned parameters in such a layer, for a kernel of height $k_H$ and width $k_W$, is
\[
	C_{\text{out}} \times \bigl(C_{\text{in}} \times k_H \times k_W + 1\bigr),
\]
the $+1$ accounting for one bias per output channel --- a formula this thesis returns to repeatedly when comparing the cost of alternative designs. Two further settings control how the kernel is applied. \textbf{Stride} is the number of positions the kernel moves between one application and the next; a stride greater than one skips positions and shrinks the output accordingly. \textbf{Padding} adds a border, typically of zeros, around the input before the kernel is applied, which controls how much the output shrinks relative to the input and lets a layer preserve spatial size exactly if that is wanted. Because a single kernel only ever looks at a small window, one convolutional layer alone only detects small, local patterns; stacking convolutional layers lets a unit in a later layer depend, indirectly, on a much larger region of the original input, since each layer's output already summarises a neighbourhood of the layer before it. This growing \textbf{receptive field} is why early convolutional layers tend to respond to local structure --- edges, small textures --- while later ones respond to larger, more composite structure built from what the earlier layers detected.

\subsection{Pooling and spatial reduction}
\label{sec:pooling}

A \textbf{pooling} operation reduces the spatial size of a feature map without learning any weights of its own: it slides a small window across the map, exactly as a kernel does, but instead of a weighted sum it reports a fixed summary of the values it covers --- most commonly the maximum, giving max pooling. Pooling deliberately discards information: a $2\times2$ max-pool halves both height and width, keeping only the strongest response in each small neighbourhood and throwing the rest away. What is bought in exchange is a lower computational cost for every layer that follows, since they now operate on a smaller map, and a larger receptive field per unit of depth, since a later layer's kernel now covers a coarser, wider-reaching grid of the original input for the same kernel size. Pooling is used where the exact position of a detected feature matters less than its coarser location and the network can no longer afford to carry full spatial detail forward.

\subsection{Parameters, capacity and small data}
\label{sec:capacity}

The total count of a network's parameters --- every weight and bias, across every layer --- determines its \textbf{capacity}: loosely, how large and how intricate a set of functions it is able to represent. A higher-capacity network can fit more complicated relationships between input and label, but capacity is measured against the data available to constrain it, not in isolation. With a training corpus of only a few hundred or few thousand labelled examples, a network with far more parameters than examples has enough freedom to fit each training example almost exactly, including whatever is peculiar to that specific example --- sensor noise, an unrepresentative pose, an idiosyncrasy of one subject --- rather than the pattern that generalises across examples. This failure mode is \textbf{overfitting}: training loss keeps falling because the network is, in effect, memorising its training set, while performance on examples it did not train on stops improving or gets worse. The property actually wanted is \textbf{generalisation} --- accurate prediction on data the network never saw during training --- and it cannot be read off the training loss at all, since a memorising network can drive that loss arbitrarily low. Detecting whether generalisation is actually happening requires a \textbf{held-out set}: a portion of the labelled data set aside from the start, never used to compute a gradient or update a parameter, on which the network's predictions are checked only after training. A gap between good performance on the training data and poor performance on the held-out set is the signature of overfitting; parameter count, relative to the number of available training examples, is what makes that gap likely in the first place. This is precisely the constraint that governs the size and construction of every architecture described next, in \autoref{sec:building-blocks}.


\section{Network Building Blocks}
\label{sec:building-blocks}

Chapter~3 argues why each of the following components was chosen and what its measured contribution was. This section stays one level below that argument, setting out the mechanism each component computes.

\subsection{Three-dimensional convolution and kernel shape}
\label{sec:conv3d}

A \texttt{Conv3d} layer generalises image convolution to a five-dimensional tensor of shape $(B, C, D, H, W)$ --- batch, channel, and three spatial-like axes, here depth $D$ (time), height $H$ and width $W$. A kernel of shape $(k_D, k_H, k_W)$ extends the sliding, weight-sharing computation of \autoref{sec:convolution} across all three non-channel axes at once, with stride and padding again governing the output's size along each.

The detail that matters most for what follows is the kernel's temporal extent, $k_D$. If $k_D > 1$, the filter reaches across multiple frames at every application, so a single output value is a genuine function of more than one time step --- a spatio-temporal filter in the literal sense. If $k_D = 1$, the filter at time step $t$ only ever reads input at time step $t$; no output value depends on any other frame. A stack of layers with kernel shape $(1, k_H, k_W)$ is therefore a two-dimensional convolution applied identically and independently to every frame, merely expressed using \texttt{Conv3d} operations so the tensor never leaves its $(C, D, H, W)$ layout. This is the case the model in this thesis uses throughout its convolutional stem: every kernel is shaped $(1, 3, 3)$, with padding $(0, 1, 1)$ preserving the spatial extent while the temporal axis is neither padded nor touched. The temporal dimension is carried through the stem unchanged in length; only $H$ and $W$ are affected, first by the convolutions' own padding and then by a spatial-only max-pool of shape $(1, 2, 2)$ that halves both.

\subsection{Normalisation and regularisation}
\label{sec:normalisation}

\texttt{BatchNorm3d} normalises each channel independently: for channel $c$ it computes a mean and variance over every other axis at once --- the batch dimension and all spatial-temporal positions $(D, H, W)$ --- and rescales,
\[
	\hat{x} = \frac{x - \mu_c}{\sqrt{\sigma_c^2 + \epsilon}}, \qquad y = \gamma_c \hat{x} + \beta_c,
\]
where $\gamma_c$ and $\beta_c$ are learned per-channel scale and shift parameters. Training uses batch statistics; inference uses a running estimate accumulated during training, so the layer behaves exactly as its two-dimensional counterpart does, only pooled over one further axis.

\texttt{Dropout3d} differs from ordinary \texttt{Dropout} in what unit it zeroes. Ordinary dropout masks individual scalar activations independently at random, which is a weak regulariser on a convolutional feature map, since neighbouring positions in $H$, $W$ and $D$ are highly correlated and the network can read a missing value off an adjacent one. \texttt{Dropout3d} instead zeroes an entire channel --- every $(D, H, W)$ position of it, for a given sample --- with the given probability, removing a whole feature detector rather than a single correlated pixel.

\texttt{LayerNorm} normalises across the feature dimension of a single position for a single sample, independent of the batch, which is why it is the default choice inside transformer blocks rather than \texttt{BatchNorm}. Here it appears inside every transformer encoder layer and once more at the head of the classifier, which applies it before a dropout and the final linear projection to the output classes.

\subsection{Parameter-free attention: SimAM}
\label{sec:simam}

SimAM~\cite{yang2021} re-weights a feature map without learning any new parameters, by scoring each neuron on how distinctive it is from its surrounding context and using that score directly as a multiplicative gate. Yang et al. define an energy $e^{*}$ that is \textit{lower} for a more distinctive neuron, so it is the \textbf{reciprocal} of that energy that serves as the gate. For a neuron with activation $x$ in a channel whose mean and variance are $\mu$ and $v$ (estimated by pooling over that channel's own spatial extent), that reciprocal has the closed form
\[
	\frac{1}{e^{*}(x)} = \frac{(x-\mu)^2}{4(v+\lambda)} + 0.5,
\]
where $\lambda$ is the module's single hyperparameter. The refined output is then $x \cdot \sigma\!\left(1/e^{*}(x)\right)$, a monotonic sigmoid gate applied element-wise, the sigmoid present only to bound the gate rather than to reshape it. What makes this mechanism unusual among attention modules is exactly what the formula shows: every quantity it needs --- the mean, the variance, the resulting score --- is read off the feature map itself, at inference as much as at training time, so no weight matrix, bias, or reduction layer is introduced anywhere. The gate is a fixed function of the activations it is applied to, not a learned one.

The implementation in this thesis follows that formula exactly, with $\lambda = 10^{-4}$, but estimates $\mu$ and $v$ by pooling over the entire spatio-temporal volume $(D, H, W)$ of a stream's feature map rather than a single frame's spatial extent alone, so a neuron's distinctiveness is judged against the whole clip.

\subsection{Self-attention}
\label{sec:self-attention}

Self-attention relates every position in a sequence to every other position by learned linear projection rather than a fixed spatial neighbourhood. Each input vector is projected three ways, into a query $Q$, a key $K$ and a value $V$; attention weights are the scaled dot product of queries against keys,
\[
	\text{Attention}(Q,K,V) = \text{softmax}\!\left(\frac{QK^\top}{\sqrt{d_k}}\right)V,
\]
with the $\sqrt{d_k}$ scaling present to stop the dot products, and hence the softmax, from saturating as the key dimension $d_k$ grows. Multi-head attention runs several such projections in parallel, each into a smaller subspace, concatenating their outputs before a final linear mixing layer, so different heads can attend to different relationships within the same sequence.

Because this is a weighted sum over the whole sequence with weights derived purely from content, it is invariant to the order of the input positions: permuting the sequence permutes the output identically, with nothing in the mechanism encoding \textit{where} a token sits. A positional encoding is added to each input vector to break that symmetry. Two schemes are common. A sinusoidal encoding assigns each position a fixed vector built from sine and cosine functions at geometrically varying frequencies across the embedding dimensions --- deterministic, requiring no training, and evaluable at sequence lengths never seen during training. A learned encoding instead assigns each position index its own trainable embedding, which can fit the data more closely but cannot extrapolate beyond the longest position trained on. The encoder in this thesis uses the fixed sinusoidal form, held as a constant buffer rather than a parameter.

A second placement choice concerns normalisation relative to each sublayer. Post-norm --- the original arrangement --- applies \texttt{LayerNorm} after the residual addition that follows each sublayer; pre-norm applies it beforehand, to the sublayer's input, so the residual path carries an un-normalised signal straight through the stack, which is generally reported to train more stably since gradients reach earlier layers unimpeded. This is the placement used here. Beyond what the encoder layer itself contributes, no additional residual or skip connection is introduced anywhere else in the architecture.

\subsection{The fallbacks}
\label{sec:fallbacks}

Each learned component here can be switched off, and each switch has a specific, near parameter-free replacement rather than simply removing that stage. With the convolutional stem off, every frame's motion channels are instead reduced by adaptive average pooling to a fixed $4\times4$ spatial grid regardless of input resolution, flattened per frame and passed through a single linear layer into the model's working dimension. With the transformer off, the sequence of per-frame feature vectors is instead collapsed by a plain mean over the time axis, with no parameters and no representation of frame order at all.

Both fallbacks add as little capacity as possible: the pooling operations introduce no parameters, and the one linear layer needed to match dimensions in the no-CNN case is negligible next to what it replaces. That near-absence of added capacity is what makes each fallback a fair control, isolating what the learned component contributes without also changing how much the network is capable of fitting.

Isolating a component fairly also depends on how the network is trained, and \autoref{sec:scarcity-skew} sets out the training regime held constant across every configuration.


\section{Learning Under Scarcity and Skew}
\label{sec:scarcity-skew}

This section sets out the training machinery --- the loss function, the sampling scheme, regularisation, and the optimiser --- at the level of what each mechanism computes and how it is configured here. Chapter~3 (Section~3.9) reviews the literature motivating these choices, gives the corpus's class counts, and reports the failure mode that follows from combining two of them carelessly; that narrative is not repeated below. What follows is the mechanism and the configuration actually used.

\subsection{What skew does to cross-entropy}
\label{sec:skew-cross-entropy}

Standard cross-entropy sums a per-example term $-\log p_t$, where $p_t$ is the model's predicted probability for the true class, weighting every example equally. Under a skewed class distribution that equal weighting is not equal in effect: a class's total gradient contribution is its example count times the per-example gradient, so a majority class dominates the sum purely by outnumbering the rest. A model facing this loss has a cheap route to a lower value --- predict close to the prior on ambiguous cases --- because that is correct often enough on the majority class to outweigh being wrong on every minority example. Loss can fall steadily while the model learns comparatively little about classes it rarely sees. The remainder of this section responds to that problem.

\subsection{Focal loss: re-weighting by difficulty, not by frequency}
\label{sec:focal-loss}

Focal loss (Lin et al., 2017) inserts a modulating factor in front of the log term:

\[
	\text{FL}(p_t) = -\alpha_t (1-p_t)^{\gamma} \log p_t.
\]

Zhao et al.~\cite{zhaoS2021} apply this formulation to micro-expression recognition and state the mechanism plainly: $\gamma$ is ``the balance factor for loss'' and $\alpha_t$ ``the weight balance factor for samples''. The distinction matters for what follows. The factor $(1-p_t)^{\gamma}$ shrinks toward zero as the model becomes confident and correct on an example, re-weighting by how \textit{difficult} it currently finds that example, irrespective of its class --- an easy majority-class example contributes almost nothing once learned; a hard example, of any class, keeps contributing. Nothing in that factor looks at how many examples of a class exist. Class frequency enters only through the separate $\alpha_t$ term, an explicit per-class multiplier ordinarily set near inverse class frequency.

The implementation here (\texttt{Ablation\_Study/losses.py}) follows this exactly: it computes \texttt{log\_softmax}, gathers $\log p_t$ for the true class, exponentiates to recover $p_t$, and forms \texttt{focal\_weight = (1 - p\_t) ** gamma}. Label smoothing (\autoref{sec:label-smoothing}) is applied first, blending the negative log-likelihood term with a uniform term over classes, so the focal weight multiplies the \textit{smoothed} loss rather than the raw NLL. The optional $\alpha$ vector, where supplied, multiplies the result afterwards, indexed by each example's true class. The value of $\gamma$ used in this study follows Zhao et al.'s reported setting; Section~4.5.1 gives it, and the label-smoothing value, together with the rest of the training configuration.

\subsection{Balanced sampling}
\label{sec:balanced-sampling}

The second lever acts on the data a batch is drawn from rather than on the loss computed over it. A \texttt{WeightedRandomSampler} assigns each training example a weight equal to the inverse of its class's count, computed over the training split of the current fold only. Sampling then proceeds with replacement, drawing \texttt{num\_samples = len(train)} indices per epoch, so each class is presented with roughly equal probability regardless of its true rarity. This sampler is attached to the training loader alone; the validation loader uses no sampler and iterates with \texttt{shuffle=False}, so validation scores reflect the split's true distribution.

\subsection{Why combining both is a choice, not a default}
\label{sec:combining-both}

Focal loss's $\alpha$ term and the balanced sampler both correct for class frequency, at different points in the pipeline, so running both at once corrects for the same imbalance twice. A pipeline that applies both is therefore making a choice, not accepting a default, and the choice has to be resolved somewhere explicit. Section~4.5.2 states how this study resolves it and which of the two corrections is left active. Section~3.9.7 (cross-referenced rather than repeated here) explains why this rule exists and what happens when it is violated; the fact to carry forward is simply that the two mechanisms are not additive by default in this codebase.

\subsection{Label smoothing}
\label{sec:label-smoothing}

Label smoothing replaces the one-hot target with a softened one, mixing a fraction $\epsilon$ of uniform probability mass across all classes into the true-class target. Inside the focal-loss implementation this is a weighted combination of the ordinary negative log-likelihood term and a term averaging $-\log p_c$ over every class $c$; Section~4.5.1 gives the value of $\epsilon$ used here. The effect is to discourage the model from driving the true-class logit arbitrarily high --- regularisation against over-confidence, applied uniformly regardless of class, and distinct from a correction for imbalance. AdamW, mixed precision and gradient clipping have no dedicated treatment in this project's \texttt{docs/} corpus and are used here as standard practice.

\subsection{Augmentation under scarcity}
\label{sec:augmentation}

With few labelled clips per class, augmentation acts before the loss or the sampler is reached at all, and is applied to the training split only. Xia et al.~\cite{xia2020a} motivate this for micro-expression data directly: ``temporal data augmentation strategies as well as a balanced loss are jointly used for our deep network'' to address ``limited and imbalanced training samples''. Two transformations are implemented (\texttt{Ablation\_Study/dataset.py}), but only one of them acts. A random temporal crop of the input window (centred at evaluation) fires only where a clip's stored length exceeds the length the model expects; extraction stores exactly the number of frames the model is configured to consume (\autoref{sec:three-channel-tensor}), so that condition never holds and the crop is inert throughout this study. What remains is a horizontal flip applied with probability one half. Because the input is optical flow rather than raw pixels (\autoref{sec:video-to-motion}), the flip cannot be a plain mirror --- reversing the scene horizontally reverses the sign of horizontal displacement --- so the flip is paired with a negation of the flow tensor's $u$ (horizontal) channel. Flipping without negating $u$ would hand the network motion vectors pointing the wrong way for the geometry shown, a quiet corruption motion-tensor augmentation must specifically guard against.

\subsection{Optimisation}
\label{sec:optimisation}

Training uses AdamW, with decoupled weight decay applied uniformly to every parameter --- no separate no-decay group for biases or normalisation terms. The schedule combines a short linear warmup with cosine annealing: \texttt{LinearLR} ramps the rate from a fraction of target up to full over the warmup epochs, after which \texttt{CosineAnnealingLR} anneals it to a small floor over the remainder; the two are joined by \texttt{SequentialLR} and stepped once per epoch. Section~4.5.4 gives the learning rate, weight decay, warmup length and epoch count used here. Cosine annealing is adopted elsewhere in the micro-expression literature on similar small-data grounds --- Zhang et al.~\cite{zhang2022} use it in a spatio-temporal transformer for this same task.

Gradient norms are clipped after the AMP scaler's gradients are unscaled, via \texttt{clip\_grad\_norm\_}, guarding against the occasional large update a small, skewed dataset can produce. Training runs under mixed precision (\texttt{torch.amp.GradScaler}, \texttt{autocast}), enabled where CUDA is available and not disabled by flag. In the same unscaled branch, and therefore only while clipping is on, every parameter's gradient is checked for non-finite values before the optimiser step; if any is found, the step is skipped entirely, gradients are zeroed, and a warning logged, rather than letting a corrupted update reach the weights. Model selection keeps the checkpoint with the highest validation macro F1, not accuracy, consistent with the metric argument of Section~3.9.5 that accuracy is a poor guide under skew; ties favour the later epoch. There is no early stopping: every configuration trains its full epoch budget regardless of when its best epoch occurred.

One category of technique is absent by design: quantisation, pruning, and knowledge distillation are used nowhere in this project. Efficiency is measured rather than engineered here --- Section~3.6 reports the compute cost of the architectural choices actually made.

Measuring rather than engineering a property requires an evaluation apparatus equal to the task; \autoref{sec:evaluation} builds the one this corpus demands, from the confusion matrix upward.


\section{Evaluating on a Small Corpus}
\label{sec:evaluation}

Chapter~3 argues why leave-one-subject-out evaluation and pooled macro F1 suit this corpus, and what those choices cost here (Sections~3.1.6 and~3.1.8). This section gives the mechanism underneath: the confusion matrix and per-class precision, recall, F1; averaging scores versus pooling counts; the easily conflated distinction between averaging across folds and within them; what cross-validation buys and what leave-one-subject-out protects against; and, to close, an audit of what this thesis's own measurement apparatus does and does not supply.

\subsection{The confusion matrix, and precision, recall, F1 for one class}
\label{sec:confusion-matrix}

For a single class $c$ treated as positive against all others, every prediction falls into one of four counts: true positives $TP_c$ (correctly predicted $c$), false positives $FP_c$ (predicted $c$ but truly something else), false negatives $FN_c$ (truly $c$ but predicted otherwise), and true negatives. Arranging all classes this way at once gives the confusion matrix: entry $(i,j)$ counts clips whose true label is $i$ and predicted label is $j$, so the diagonal holds correct predictions and every off-diagonal cell names a specific error --- which true class is mistaken for which other.

From those counts, precision and recall isolate two failure modes:
\[
	\text{Precision}_c = \frac{TP_c}{TP_c + FP_c}, \qquad \text{Recall}_c = \frac{TP_c}{TP_c + FN_c}.
\]
Precision penalises false alarms: it falls whenever the classifier calls $c$ on a clip that is not $c$, regardless of how many true $c$ instances it also catches. Recall penalises misses: it falls whenever a genuine $c$ instance is assigned elsewhere, regardless of how many false alarms accompany the hits. Neither failure mode is visible from the other, which is why F1, the harmonic mean,
\[
	F1_c = \frac{2 \cdot \text{Precision}_c \cdot \text{Recall}_c}{\text{Precision}_c + \text{Recall}_c} = \frac{2\,TP_c}{2\,TP_c + FP_c + FN_c}
\]
is the usual single-number summary: being harmonic, it stays low whenever either input is low, so a model cannot buy a good F1 by trading recall for precision.

\subsection{Macro versus micro averaging}
\label{sec:macro-micro}

Combining per-class F1 scores into one number can be done two ways. Macro averaging computes $F1_c$ per class and takes an unweighted mean, $\text{Macro-F1} = \tfrac1C\sum_c F1_c$: every class counts equally regardless of size, so ignoring a rare class is penalised as heavily as ignoring a common one. Micro averaging instead pools the raw counts across classes first --- summing every class's $TP$, $FP$ and $FN$ --- and computes one F1 from those totals.

For single-label multi-class classification, where every clip gets exactly one predicted label, this pooling has a consequence worth stating plainly: every misclassification is simultaneously a false positive for the class it was wrongly assigned to and a false negative for its true class, so summed $FP$ and summed $FN$ equal each other and the total error count, and pooled micro-F1 reduces algebraically to $TP_{\text{total}}/N$ --- plain accuracy. Micro-F1 and accuracy are not two different metrics here; they are the same number computed two ways. The contrast that actually matters is therefore not ``macro versus micro'' but \textbf{macro-F1 versus accuracy}: an unweighted per-class average against a count a majority class can dominate.

\subsection{Pooled versus per-fold averaging}
\label{sec:pooled-vs-fold}

A further averaging choice sits underneath macro-F1 once evaluation is split into folds, and it governs how every result in this thesis is reported. One option accumulates the confusion counts --- $TP_c$, $FP_c$, $FN_c$ --- across every fold first, then computes one set of per-class F1 scores, and one macro-F1, from that pooled total. The other computes a complete macro-F1 inside each fold from that fold's own predictions, then averages the per-fold values. See et al.~\cite{see2019} define the Unweighted F1 (UF1) of the MEGC 2019 protocol the first way: the macro-averaged F1 obtained by accumulating true positives, false positives and false negatives over all folds of a leave-one-subject-out run before computing and averaging per-class F1.

These are estimators of two different quantities, not two ways of writing the same one: pooled macro-F1 describes this class's F1 across the whole evaluation, once every held-out prediction sits in one confusion matrix together, while per-fold-averaged macro-F1 describes the typical per-fold macro-F1 --- a quantity that depends on how classes happen to fall within each fold, not only on the classifier. When a fold's held-out set does not contain every class, its per-class F1 for the missing class is degenerate, and averaging such folds in with fully-populated ones changes what the number reflects for reasons unrelated to classifier quality; on a small, unevenly distributed corpus this is a live possibility, and fold composition can bound the per-fold estimator below what the pooled one would report. Sections~3.1.6 and~3.1.8 work through how the fold structure of this study does exactly that; which one a reported number is becomes apparent only once its computation is stated.

\subsection{Cross-validation and subject-disjointness}
\label{sec:cross-validation}

$k$-fold cross-validation partitions the data into $k$ disjoint subsets, trains $k$ models --- each on $k-1$ folds --- and evaluates each on the one fold it never saw, so every example is scored exactly once by a model that never trained on it. Leave-one-subject-out (LOSO) is the case where $k$ equals the number of subjects, and each fold's held-out set is precisely one subject's clips, every other subject's clips forming that fold's training set.

The property this buys is subject-disjointness: no fold's training and validation sets ever share a subject. Without it, a model evaluated on a clip from a subject it has already trained on can succeed by recognising that person's face or recording conditions rather than the expression class itself --- an identity leak inflating the score without the model having learned anything transferable to an unseen person. Enforcing subject-disjointness in every fold closes that channel.

\subsection{What an evaluation apparatus of this kind can fail to provide}
\label{sec:apparatus-gaps}

Three shortcomings belong to an evaluation apparatus rather than to any model it measures, and each matters more at this corpus size than it would at a larger one. They are set out here as properties to look for; Section~4.6 reports where this study's apparatus stands on each.

The first is \textbf{which aggregate a pipeline actually stores}. \autoref{sec:pooled-vs-fold} distinguishes pooling the confusion matrices before computing a macro average from averaging per-fold macro scores, and shows the two answer different questions. A pipeline may compute one, the other, or both under the same field name, and the quantity a reader wants is not necessarily the one written to disk.

The second is \textbf{whether model selection is blind to the data it is scored against}. A nested design keeps three sets apart: one to fit parameters, a second to choose between candidate models or epochs, and a third, untouched until the end, to score the chosen one. Collapsing the second and third --- selecting a checkpoint on the same data that then reports the result --- leaves the reported figure not blind to what it is checked against, and biases it optimistically by an amount that cannot be recovered after the fact.

The third is \textbf{whether enough is retained to support inference}. Aggregate metrics answer what a configuration scored; they cannot answer whether two configurations differ by more than chance. A paired test between configurations needs the per-clip predictions of each, so an apparatus that persists only aggregates forecloses that question whatever its scores say.


\section{Conclusion}
\label{sec:ch2-conclusion}

This chapter has set out the mechanisms on which the rest of the thesis depends, in the order the pipeline applies them. \autoref{sec:phenomenon} defined the micro-expression by the three properties that make it hard to recognise automatically --- it is brief, it is faint, and it cannot be produced on instruction --- and fixed the vocabulary of onset, apex and offset, noting that only the emotion label, not the action-unit coding, is used here. \autoref{sec:recording-to-input} established what actually reaches the pipeline: a 200~fps recording, the corpus's original frames rather than its registered release, and a greyscale frame at a fixed resolution, with no detection, registration or cropping performed by this project. \autoref{sec:motion-magnification} and \autoref{sec:video-to-motion} covered the two transformations applied to those frames --- Eulerian magnification of small motions, and the conversion of the magnified sequence into dense optical flow and optical strain.

\autoref{sec:nn-fundamentals} then built the machinery of a neural network from a single unit up to the convolutional layer, ending with the relationship between parameter count and corpus size that constrains every architectural choice this thesis makes. \autoref{sec:building-blocks} described the specific blocks assembled on that foundation, together with the parameter-free replacement each one falls back to when switched off. \autoref{sec:scarcity-skew} covered training under a class distribution that cannot be balanced by collecting more data, and \autoref{sec:evaluation} the evaluation apparatus needed when a corpus is small enough that the choice of averaging changes the conclusion.

Two threads run through all of it. The first is that the size of the corpus, rather than any modelling preference, forces most of the design: it dictates how few parameters the network can carry, why corrections for class skew are necessary, and why the evaluation protocol has to be chosen with care. The second is that several stages depart from the textbook or reference form of the technique they implement, and each departure has been named at the point it occurs rather than left for a reader to discover.

What this chapter has deliberately not done is say who else has used any of these techniques, with what result, or what remains unresolved. That is the work of Chapter~3, which reviews the literature for each component in turn and identifies the gap this thesis addresses.

%
% In-chapter \autoref targets resolve once all nine sections are present.
% Cross-chapter references (Chapter 3, Chapter 4, Section 3.1.2, ...) are
% written as plain text until those chapters carry labels.
%
% Citation keys used in this chapter:
%   yan2014     [1]  Yan et al., CASME II
%   qu2016      [2]  Qu et al., CAS(ME)^2
%   liY2018     [3]  Li, Huang & Zhao, single apex frame
%   bai2021     [4]  Bai et al., magnification + pre-trained network
%   liong2019a  [5]  Liong et al., OFF-ApexNet
%   shreve2011  [6]  Shreve et al., spatio-temporal strain
%   zhaoS2021   [7]  Zhao et al., two-stage 3D CNN
%   yang2021    [8]  Yang et al., SimAM
%   xia2020a    [9]  Xia et al., STRCN
%   zhang2022  [10]  Zhang et al., SLSTT
%   see2019    [11]  See et al., MEGC 2019
% =============================================================================

\chapter{Background}
\label{ch:background}

This chapter sets out the mechanisms the rest of the thesis depends on, in the order the pipeline applies them: what a micro-expression is, what the recording supplies, how small motions are magnified and converted into a motion representation, what a neural network does, which blocks this system assembles from that foundation, how it is trained under a corpus that cannot be balanced, and how it is measured when that corpus is small.

It explains \textbf{what each mechanism is and how it works}. Chapter~3 surveys who has used each one, with what result, and what gap remains; where a fact is needed in both places it is established here and cross-referenced there. Scope is limited to what the implemented system actually does --- techniques a reader might expect in a background chapter but which this project does not use are omitted, and their absence is noted where it would otherwise be assumed.

Two constraints recur. The corpus is small enough that its size, rather than any modelling preference, forces most of the design. And several stages depart from the textbook form of the technique they implement; each departure is named where it occurs.


\section{The Phenomenon: Micro-Expressions as Involuntary, Brief, Low-Intensity Facial Movement}
\label{sec:phenomenon}

A micro-expression is a brief facial movement that leaks a felt emotion the person is actively trying to conceal~\cite{yan2014}. It is best understood by contrast with an ordinary, or macro-, expression, on three axes rather than one. Macro-expressions typically last upwards of half a second, and may run to several seconds, and can be produced at will; micro-expressions, by contrast, are ``characterized by short durations, involuntary generation and low intensity''~\cite{qu2016}. The third axis is the one most easily overlooked: a micro-expression cannot be performed on instruction. It is a symptom of concealment rather than a communicative act. That single property governs how such data can be gathered at all, and the consequences for corpus construction are taken up in Section~3.1.

\subsection{Duration: the first constraint}
\label{sec:duration}

The field's working definition fixes an upper bound on duration. Yan et al.~\cite{yan2014} state that ``the generally accepted upper limit of the duration is 1/2 s.'' Brevity is the first source of the difficulty in recognising the phenomenon automatically, and the mechanism is simple: whatever signal distinguishes one emotion from another must be carried by however many video frames fit inside that half-second window. At a modest frame rate that window may contain only a handful of frames, several of which are near-identical to their neighbours; the sequence available to a classifier is short by construction, not by any deficiency of the recording.

\subsection{Low intensity: the second constraint}
\label{sec:low-intensity}

Duration alone would still leave a tractable problem if the movement were large. It is not. Yan et al.~\cite{yan2014} observe that micro-expressions are ``usually low in intensity -- it might be so brief for the facial muscles to become fully-stretched with suppression'', with the consequence that ``because of the short duration and low intensity, it is usually imperceptible or neglected by the naked eyes''. In practice the movement of interest may change only a few grey levels over a small patch of skin --- around an eyebrow, or the corner of a mouth --- while the rest of the face stays static. This is a distinct obstacle from brevity: a longer recording does nothing to help it, because the difficulty is not how many frames are available but how far the signal in each frame sits above the noise floor of ordinary video. Duration limits how much evidence exists; intensity limits how visible that evidence is once it does.

\subsection{Onset, apex and offset}
\label{sec:onset-apex-offset}

Because a micro-expression is a movement rather than a static configuration, it is conventionally described by three temporal landmarks. The \textbf{onset} is the frame at which the face begins to depart from its neutral baseline; the \textbf{offset} is the frame at which it returns to neutral; and the \textbf{apex} is the frame at which the movement is judged most intense --- the point at which, as Ekman puts it, a ``snapshot taken at an [sic] point when the expression is at its apex can easily convey the emotion message'' (as reported by Li, Huang \& Zhao~\cite{liY2018}). All three landmarks are available in the corpus used here. This thesis, however, uses only the onset and offset frames to bound each clip passed into the pipeline; the apex frame is never read. This is a deliberate design choice rather than an oversight, and its consequences --- for what information the model can and cannot see, and for how its results relate to apex-frame methods in the literature --- are taken up in Chapters~3 and~4.

\subsection{Action units, and the label actually used}
\label{sec:action-units}

The Facial Action Coding System (FACS) describes any facial movement, however subtle, as a combination of discrete Action Units (AUs), ``an objective method for labeling facial movements in terms of component actions''~\cite{yan2014}. An AU is a description of movement, not of emotion, and the mapping between the two is not one-to-one; corpora therefore assign emotion categories using AUs together with other evidence, by a procedure that differs between datasets and is described for the corpus used here in Section~3.1.2.

It is worth being precise about how much of that machinery survives into this thesis: none of it. The work reported here uses only the resulting emotion label. The Action Units recorded for each clip are present in the label table as inherited metadata from the original annotation --- carried along because the released coding includes them --- but they are never read by any stage of the data pipeline and never presented to the model. Nothing in this thesis performs, or claims to perform, action-unit recognition; the sole target throughout is the emotion category.

How brief and how faint that movement is determines what the recording must capture, and \autoref{sec:recording-to-input} turns to the video this pipeline actually receives.


\section{From Recording to Pipeline Input}
\label{sec:recording-to-input}

Before any mechanism of interest --- magnification, flow, strain, or a trained network --- can act on a micro-expression, something has to capture it, and something has to hand it to the pipeline in a usable form. This section describes that handoff: what the recording had to provide given the constraints fixed in \autoref{sec:phenomenon}, what the corpus actually supplies, and what this project does and does not do to a clip before it enters processing.

\subsection{High-speed capture}
\label{sec:high-speed-capture}

\autoref{sec:duration} fixed a half-second upper bound on the duration of a micro-expression. That bound has an immediate consequence for recording: whatever number of frames fit inside the event is the entire evidence a classifier will ever see of it, so the frame rate at which the video was captured is not a matter of picture quality but a precondition for the signal existing in the data at all. A 300~ms movement recorded at 25 or 30~fps --- ordinary video rates --- yields only seven or eight frames. The same event recorded at 200~fps yields roughly sixty frames, enough for onset, apex and offset to be distinct, separately observable moments rather than a blur collapsed into a handful of frames.

CASME II was recorded at 200~fps for exactly this reason~\cite{yan2014}. The pipeline used in this thesis takes 200 as a configured constant (\texttt{fps: int = 200} in \texttt{CASMEIIConfig}, \texttt{Stage1\_DataPipeline/config.py}), rather than reading it from the video files themselves, since it is a fixed property of the corpus rather than something to be inferred per clip.

\subsection{What the corpus supplies, and what this project does not do}
\label{sec:corpus-supplies}

CASME II is distributed in two forms: the frames as they were recorded, and a normalised release in which each face has been cropped and registered. This pipeline reads the first. Its metadata-unification stage resolves each clip to a per-subject, per-clip folder of the original recorded frames (\texttt{Stage1\_DataPipeline/metadata\_unifier.py}), so what enters processing is the full recorded field of view at 640~$\times$~480 (Section~3.1.2), within which the face occupies only a part.

It is worth stating plainly what this means, because the alternative assumption is the natural one to make. This project performs no face detection, no landmark localisation, no geometric registration, no alignment and no cropping of its own, anywhere in the pipeline --- and it does not inherit those operations from the corpus either, because it does not read the release that carries them. The normalisation that produced that release --- 68-point Active Shape Model landmarking on each clip's first frame, registered to a neutral model face by a Local Weighted Mean transform, described in Chapter~3, Section~3.1.2 --- is the corpus authors' work, and Section~3.1.2 records that this thesis departs from it. The face reaching the network is therefore an unregistered one, at whatever position and scale the camera recorded it, and whatever tolerance to that variation the model acquires it must acquire from the data.

\subsection{Frame selection and preparation}
\label{sec:frame-preparation}

Each frame is read in greyscale and resized to 224~$\times$~224 by bilinear interpolation. The resize is applied unconditionally to the whole frame, so it both downsamples and, since 640~$\times$~480 is not square, does not preserve the recorded aspect ratio. From the frames spanning a clip's annotated onset to its annotated offset, a fixed number of frames is selected for further processing --- the scheme by which that selection is made is described in \autoref{sec:three-channel-tensor}, together with the tensor it produces.

\subsection{Which clips enter the pipeline}
\label{sec:clip-selection}

Not every row of the corpus's coding table reaches processing. Four conditions are applied to select the rows that do: the row must belong to the CASME II dataset, its coded expression type must be \texttt{micro-expression} rather than one of the other categories the table records, its frame folder must be present on disk and non-empty, and its sequence length --- the onset-to-offset frame count, counted inclusively --- must exceed two. The first two conditions restrict the pipeline to the corpus and the expression class this thesis addresses, and the third discards rows whose media the coding table lists but the filesystem does not supply. The fourth rules out clips too short to be usable: a sequence of two frames or fewer contains at most one frame-to-frame transition, which is not enough to construct a motion sequence of any length, whatever downstream representation is built from it.

Together, these properties of the recording and the corpus --- a 200~fps capture rate, an unregistered full recorded frame, a greyscale 224~$\times$~224 resize of it, and a filtered, onset-to-offset clip --- are what reaches the pipeline before any processing of interest begins. Eulerian motion magnification is the first such process applied to these frames, and \autoref{sec:motion-magnification} describes the mechanism by which it amplifies motion without ever tracking it.


\section{Motion magnification}
\label{sec:motion-magnification}

\subsection{The Eulerian idea}
\label{sec:eulerian-idea}

A moving object can be described in two ways. A \textit{Lagrangian} description follows a point on the object as it moves --- this is what optical flow and feature tracking do: locate a point in frame $t$, find its correspondence in frame $t+1$, and report the displacement. A \textit{Eulerian} description instead fixes attention on a location in the image and asks how the signal observed there changes over time. Eulerian video magnification (EVM) takes the second view: no correspondence between frames is computed. At every pixel, it treats the sequence of intensity values over time as a 1-D signal, isolates the part occurring in a chosen temporal frequency band, multiplies it by a factor $\alpha$, and adds the result back to the original sequence.

The justification for why this amplifies \textit{motion}, and not merely intensity, is a first-order argument. If a small one-dimensional image profile $f(x)$ translates by a time-varying displacement $\delta(t)$, the observed intensity at a fixed location is $I(x,t) = f(x+\delta(t))$. A first-order Taylor expansion around $x$ gives

\[
	I(x,t) \approx f(x) + \delta(t)\,\frac{\partial f(x)}{\partial x},
\]

so the temporal variation at a fixed pixel is approximately proportional to the true displacement $\delta(t)$, scaled by the local spatial gradient. Amplifying that temporal variation by $\alpha$ before adding it back is therefore, to first order, equivalent to synthesising a signal consistent with a displacement of $(1+\alpha)\delta(t)$ rather than $\delta(t)$: motion is amplified as a consequence of amplifying intensity change, with no tracking step anywhere in the process. This is the account given by Bai et al.~\cite{bai2021} (following Wu et al.).

\subsection{Spatial decomposition: the Laplacian pyramid}
\label{sec:laplacian-pyramid}

The Taylor approximation above only holds where the spatial gradient is well-behaved relative to the displacement, so in practice the frame is first decomposed into a set of spatial-frequency bands and the argument is applied band-by-band. The standard tool is the \textbf{Laplacian pyramid}. Each level is built by blurring the current image with a small low-pass kernel, downsampling by a factor of two, upsampling the result back to the original resolution, and subtracting it from the pre-downsampling image. The subtraction discards everything the blur-and-downsample step could reconstruct, leaving only the spatial detail specific to that scale --- the Laplacian band. The downsampled image feeds the next iteration, and the recursion terminates in a coarse, heavily blurred residual (a Gaussian, not a Laplacian, band). A pyramid of depth $L$ yields $L$ Laplacian bands, each isolating a different range of spatial frequencies, plus one coarsest residual that carries no band-specific detail and exists only to seed reconstruction.

\subsection{Temporal filtering}
\label{sec:temporal-filtering}

Each spatial band is treated as a stack of scalar time series, one per pixel, and passed through a \textbf{temporal band-pass filter} that keeps only the variation falling inside a chosen frequency range and discards the rest. The filtered signal at each band is multiplied by $\alpha$ and added back to that band's original values; the bands are then recombined by repeating the pyramid construction in reverse --- upsampling the coarsest residual and successively adding each (now-amplified) Laplacian band back in. The published amplitude-based method realises the band-pass step with a \textbf{Butterworth filter}, a smooth infinite-impulse-response design~\cite{bai2021}.

\subsection{The amplification trade-off}
\label{sec:amplification-tradeoff}

The magnification factor $\alpha$ trades sensitivity against artefact. A larger $\alpha$ makes a fainter in-band motion visible, but the same multiplication applies indiscriminately to everything the band-pass filter passes: sensor noise, compression artefacts, and any other real motion whose temporal frequency happens to fall inside the chosen band are amplified exactly as the motion of interest is. The method has no way to distinguish a wanted signal from an unwanted one occupying the same frequency range --- the filter only knows how fast a signal oscillates, not what it is. Raising $\alpha$ therefore does not monotonically improve the result; past some point amplified noise and induced displacement artefacts dominate whatever gain in visibility was sought. This tension is intrinsic to the method rather than a defect of any one implementation, and it recurs, with numbers, in Chapter~3.

\subsection{Departures from the textbook method in this implementation}
\label{sec:evm-departures}

The magnifier used in this thesis (\texttt{Stage1\_DataPipeline/evm\_magnifier.py}) implements the mechanism above with a genuine Laplacian pyramid --- a $[1,4,6,4,1]/16$ blur kernel, downsample-by-two, upsample, subtract --- followed by a coarsest Gaussian residual, matching \autoref{sec:laplacian-pyramid} exactly. It departs from amplitude-based magnification as originally formulated in four respects; the third is a property of that original method rather than of the summary given by Bai et al.~\cite{bai2021}.

First, the temporal band-pass is an \textbf{ideal filter}: a hard binary mask applied directly to the discrete Fourier transform of each pixel's time series (via \texttt{scipy.fftpack}), rather than a Butterworth (IIR) filter. Frequencies inside $[\omega_{\mathrm{low}}, \omega_{\mathrm{high}}]$ pass with a gain of exactly one and everything else is zeroed, with no smooth roll-off at the band edges.

Second, every Laplacian band is amplified by the \textbf{same scalar $\alpha$}. Bai et al.~\cite{bai2021} state the amplified reconstruction with a \textit{frequency-dependent} factor $\alpha_k$ --- one per band, reduced at higher spatial frequencies to limit noise amplification --- and no such damping is applied here. The coarsest residual is never filtered or amplified --- it only seeds reconstruction.

Third, the implementation works directly on the greyscale intensity tensor, so there is no chromatic attenuation step.

Fourth, and a property of pipeline ordering rather than of the magnifier itself: the calling code (\texttt{tensor\_pipeline\_manager.py}) resamples each clip to a fixed 33 frames \textit{before} invoking the magnifier, and passes the corpus's nominal recording rate (200~fps) as the filter's frame-rate parameter regardless of the resampled sequence's true frame spacing. The consequence --- a temporal filter whose frequency axis no longer matches the signal it is filtering --- is analysed in Section~3.5.7, not here.

The operating point this study runs at is given in Section~4.2.5; whatever the amplification factor, band and pyramid depth, the amplified output is clipped to $[0,255]$ and quantised back to 8-bit before being handed to the downstream flow-and-strain extractor.

What that extractor does with the amplified sequence --- and why motion, rather than appearance, is what the network is given --- is the subject of \autoref{sec:video-to-motion}.


\section{From Video to Motion: Optical Flow and Optical Strain}
\label{sec:video-to-motion}

\subsection{Why motion, not pixels}
\label{sec:why-motion}

A micro-expression is a movement, so the natural quantity to hand a classifier is a description of movement rather than of appearance. Raw pixel intensity carries the signal of interest --- the same faint change described in \autoref{sec:low-intensity} --- embedded in everything else the camera also recorded: who the subject is, how the scene is lit, the fixed geometry of a particular face. None of that is discarded by presenting a network with frames directly, and on a corpus this narrow in identity and illumination a model has every incentive to fit the wrong part of the image. A displacement field is different in kind: it records how each point moved, not what colour it is, so identity and illumination do not enter it at all. That principle motivates the rest of this section; the comparative case for it is made in Section~3.3.1.

\subsection{Optical flow: the brightness-constancy constraint}
\label{sec:optical-flow}

Optical flow estimates, for every pixel, the apparent two-dimensional displacement between two frames. The estimate rests on the brightness-constancy assumption: a point on a moving surface is taken to have the same intensity in the next frame as in the current one, so that any change in a pixel's grey level is attributed to something having moved past that location, never to the thing itself changing colour or shade~\cite{liong2019a}.

Writing $I_t(x,y)$ for the image intensity at position $(x,y)$ and time $t$, and supposing the point there moves to $(x+\delta x,\, y+\delta y)$ by time $t+1$, brightness constancy states

\[
	I_t(x,y) = I_{t+1}(x+\delta x,\, y+\delta y), \qquad \delta x = u_t\,\delta t,\;\; \delta y = v_t\,\delta t,
\]

where $u_t(x,y)$ and $v_t(x,y)$ are the horizontal and vertical components of the flow at that point. Expanding the right-hand side by a first-order Taylor series and substituting into the constancy equation, the intensity terms cancel and division by $\delta t$ leaves the \textbf{optical flow constraint equation}:

\[
	u_t(x,y)\,\frac{\partial I}{\partial x} + v_t(x,y)\,\frac{\partial I}{\partial y} + \frac{\partial I}{\partial t} = 0.
\]

This single scalar equation holds at every pixel but contains two unknowns, $u_t$ and $v_t$: the image gradients are measurable directly from the frames, while the displacement is not determined by them alone. This under-determinacy is the \textbf{aperture problem} --- a local patch of image constrains the component of motion perpendicular to an edge or gradient, but says nothing about motion parallel to it. One equation per pixel cannot fix two unknowns per pixel, so every practical estimator closes the system by adding a further assumption relating the flow at neighbouring pixels, typically that the field varies smoothly across the image. That added assumption is what distinguishes one flow algorithm from another; this thesis does not adjudicate between them (Section~3.3.3) and simply adopts one, described in \autoref{sec:three-channel-tensor}. The output, for a pair of frames, is the pair of scalar fields $u(x,y)$ and $v(x,y)$ --- the horizontal and vertical displacement at every pixel --- forming the first two channels of the motion representation used throughout this thesis.

\subsection{Optical strain: the finite strain tensor}
\label{sec:optical-strain}

Flow answers where a point moved; it does not by itself distinguish a patch of skin that moved together with its neighbours from one that stretched or sheared relative to them. Optical strain makes that distinction by taking the spatial derivative of the flow field rather than the field itself.

Writing the displacement at a point as $\mathbf{u} = [u, v]^{\mathsf T}$, the \textbf{finite strain tensor} is defined as the symmetric part of the displacement gradient~\cite{shreve2011}:

\[
	\varepsilon = \tfrac{1}{2}\left[\nabla\mathbf{u} + (\nabla\mathbf{u})^{\mathsf T}\right] = \begin{bmatrix} \varepsilon_{xx} & \varepsilon_{xy} \\ \varepsilon_{yx} & \varepsilon_{yy} \end{bmatrix}, \qquad \varepsilon_{xx} = \frac{\partial u}{\partial x},\;\; \varepsilon_{yy} = \frac{\partial v}{\partial y},\;\; \varepsilon_{xy} = \varepsilon_{yx} = \frac{1}{2}\left(\frac{\partial u}{\partial y} + \frac{\partial v}{\partial x}\right).
\]

$\varepsilon_{xx}$ and $\varepsilon_{yy}$ are the \textbf{normal strain} components, describing stretching or compression along each axis; $\varepsilon_{xy}$ is the \textbf{shear strain} component, describing the change of angle between two initially perpendicular directions on the surface. The tensor is reduced to a single scalar per pixel. Shreve et al.~\cite{shreve2011} sum the component values to obtain a strain magnitude; the reduction implemented here is instead a root sum of squares, with the shear term counted once,

\[
	\varepsilon_{\text{mag}} = \sqrt{\varepsilon_{xx}^{2} + \varepsilon_{yy}^{2} + \varepsilon_{xy}^{2}},
\]

a form that diverges from the convention used across the reviewed corpus. Section~3.4.3 records that divergence and what it costs.

The property that makes this useful, rather than a second copy of the flow field, follows from its definition as a \textit{gradient}. If a region of the face translates rigidly --- every pixel moving by the same $u$ and the same $v$, as a small head movement produces --- then $u$ and $v$ are locally constant there, and their spatial derivatives, hence every component of $\varepsilon$, are zero: a rigid translation contributes nothing to the strain field, however large it is. A localised muscle contraction, by contrast, moves neighbouring points of skin by different amounts, so the displacement field has a non-zero gradient exactly where the deformation occurs, and $\varepsilon_{\text{mag}}$ is large there. Strain is thus sensitive to non-rigid deformation specifically, insensitive by construction to whatever rigid motion the flow field also contains.

\subsection{Assembling the three-channel tensor}
\label{sec:three-channel-tensor}

For each clip, the frames spanning onset to offset are uniformly resampled to a fixed sequence of 33 frames, and dense optical flow is computed between every adjacent pair using OpenCV's implementation of Farneb\"ack's method~\cite{zhaoS2021} --- \texttt{cv2.calcOpticalFlowFarneback}, with pyramid scale 0.5, three pyramid levels, window size 15, three iterations, a polynomial neighbourhood of size 5 and a polynomial smoothing of 1.2, applied to 8-bit greyscale frames. This yields \textbf{32 flow fields} per clip, each a $(u, v)$ pair over the 224~$\times$~224 face region. Strain is computed from each field's spatial gradients using \texttt{np.gradient} (central differences at unit pixel spacing, over the two spatial axes only; no temporal gradient is used), giving $\varepsilon_{xx}$, $\varepsilon_{yy}$, $\varepsilon_{xy}$ and hence $\varepsilon_{\text{mag}}$ for each of the 32 pairs.

The three quantities --- $u$, $v$ and $\varepsilon_{\text{mag}}$ --- are stacked as three channels, giving a tensor of shape $(3, 32, 224, 224)$ in channel order $[u, v, \text{strain}]$. When Eulerian video magnification is applied (\autoref{sec:motion-magnification}), flow and strain are computed on the magnified frames rather than the originals; the construction is otherwise unchanged.

Normalisation is applied in two separate stages, and both remain active. At extraction time, each of the three channels is independently min--max normalised to $[0, 1]$ across the whole clip, so that a channel's own range --- which differs enormously between a flow component and a strain magnitude --- does not by itself determine its influence. Separately, at training load time (\texttt{Ablation\_Study/dataset.py}, with \texttt{normalize=True}), each channel is z-scored per sample over its $T \times H \times W$ extent. Neither stage substitutes for the other: the first is a per-clip rescaling done once, during data preparation; the second is a per-sample standardisation applied every time a clip is loaded, centring and scaling each channel to zero mean and unit variance.

How this tensor is built is now settled; what consumes it is not. \autoref{sec:nn-fundamentals} begins with what a neural network is, before \autoref{sec:building-blocks} reaches the blocks themselves.


\section{Neural Network Fundamentals}
\label{sec:nn-fundamentals}

Everything in this thesis described as ``trained'' is, underneath, a neural network: a function built from simple, repeated computations whose numerical parameters are set from data rather than specified by hand. This section builds that machinery from its smallest unit upward, only as far as is needed for the architecture-specific mechanisms of \autoref{sec:building-blocks}, the training regime of \autoref{sec:scarcity-skew} and the evaluation apparatus of \autoref{sec:evaluation} to be legible when they introduce anything that depends on it.

\subsection{Units, layers and the forward pass}
\label{sec:units-layers}

The basic computational unit takes a vector of inputs $x$, forms a weighted sum of them plus an offset term, and passes the result through a fixed non-linear function $\phi$:
\[
	a = \phi(w^\top x + b).
\]
The entries of $w$ are the unit's \textbf{weights}, one per input; $b$ is its \textbf{bias}. Weights and biases together are the unit's \textbf{parameters} --- the numbers a training procedure is free to adjust --- and $a$, the unit's output, is its \textbf{activation}. A \textbf{layer} is a bank of such units applied to the same input vector in parallel, each with its own weights and bias, producing a vector of activations rather than a single scalar. A network is a sequence of layers, each layer's output vector becoming the next layer's input; evaluating the network on an input by applying each layer in turn, from the first to the last, is the \textbf{forward pass}, and its final output is the network's prediction. The number of layers is the network's \textbf{depth}; the number of units in a given layer is that layer's \textbf{width}. Both are choices made when the architecture is specified, not something training decides.

\subsection{Non-linearity}
\label{sec:non-linearity}

If $\phi$ were the identity function, every layer would compute only an affine transformation of its input, and the composition of any number of affine transformations is itself a single affine transformation. A network built entirely this way would therefore compute no more than a network with one layer could, however many layers it stacked --- depth would buy nothing. The non-linear function $\phi$, applied after the weighted sum at every unit, is what stops layers from collapsing into one another; it is the reason depth is a meaningful design axis at all rather than a redundant one. A common choice, used throughout the components described in this thesis, is the rectified linear unit, $\text{ReLU}(x) = \max(0, x)$: it passes positive values through unchanged and clips negative ones to zero. Its appeal is mechanical rather than subtle --- it is cheap to compute, its derivative is either zero or one almost everywhere, and unlike activation functions that saturate on both sides, a unit with a positive input keeps gradient flowing through it at full strength.

\subsection{Learning by gradient descent}
\label{sec:gradient-descent}

A network's forward pass is fixed once its parameters are fixed, so ``training'' means searching for parameter values under which the forward pass produces good predictions. A \textbf{loss function} quantifies how far a prediction sits from its true label, returning a single number that is small when they agree and large when they do not; the specific loss used in this thesis, and why, is a matter for \autoref{sec:scarcity-skew}, but the role a loss plays is a property of every trained network and belongs here. Training proceeds by adjusting parameters to make the loss smaller, and the mechanism that makes that adjustment tractable is \textbf{backpropagation}: because the forward pass is a composition of layers, the chain rule from calculus lets the gradient of the loss with respect to every parameter in the network --- however deep --- be computed by working backward from the output layer to the input, reusing intermediate results rather than recomputing each parameter's gradient from scratch. Once every parameter's gradient is known, an \textbf{update rule} moves each parameter a small step against its gradient, since the gradient points in the direction the loss increases fastest and moving the opposite way decreases it. The size of that step is the \textbf{learning rate}. Parameters are not updated from the whole training set at once in a single step; instead the data is divided into \textbf{batches}, and one update is made per batch, while one full pass through every batch in the training set is one \textbf{epoch}, of which training runs many in succession. Throughout, a network has two distinct modes of use: in training mode its parameters are being adjusted from labelled examples by the procedure just described, while in evaluation mode its parameters are held fixed and it is only asked to produce predictions, with no update following from how correct they are.

\subsection{Convolution as a learned filter}
\label{sec:convolution}

A layer of the kind described in \autoref{sec:units-layers} connects every input to every unit, which is wasteful for an image: a pattern worth detecting --- an edge, a corner, a small patch of texture --- can appear anywhere in the frame, and a fully connected layer would need to learn a separate copy of the detector for every position it might appear at. A convolutional layer instead learns one small filter, a \textbf{kernel}, consisting of a grid of weights (a $3\times3$ or $5\times5$ window is typical) plus a single bias, and slides that same kernel across every position of the input, at each position computing a weighted sum of the pixels currently under it. Because the identical kernel is reused at every position, the layer is said to exhibit \textbf{weight sharing}: it looks for the same pattern everywhere in the image and costs only as many parameters as the kernel itself holds, regardless of how large the input is. The result of sliding one kernel across the whole input is a single \textbf{feature map}, one value per position, recording how strongly that kernel's pattern was present there. A convolutional layer ordinarily learns several kernels at once, each producing its own feature map; stacked together these form the layer's output \textbf{channels}, so a layer with $C_{\text{out}}$ kernels produces $C_{\text{out}}$ channels from whatever number of input channels, $C_{\text{in}}$, it was given (three for an ordinary colour image, one for a greyscale image, or more once earlier layers have produced feature maps of their own). Because each kernel spans every input channel, not just one, the number of learned parameters in such a layer, for a kernel of height $k_H$ and width $k_W$, is
\[
	C_{\text{out}} \times \bigl(C_{\text{in}} \times k_H \times k_W + 1\bigr),
\]
the $+1$ accounting for one bias per output channel --- a formula this thesis returns to repeatedly when comparing the cost of alternative designs. Two further settings control how the kernel is applied. \textbf{Stride} is the number of positions the kernel moves between one application and the next; a stride greater than one skips positions and shrinks the output accordingly. \textbf{Padding} adds a border, typically of zeros, around the input before the kernel is applied, which controls how much the output shrinks relative to the input and lets a layer preserve spatial size exactly if that is wanted. Because a single kernel only ever looks at a small window, one convolutional layer alone only detects small, local patterns; stacking convolutional layers lets a unit in a later layer depend, indirectly, on a much larger region of the original input, since each layer's output already summarises a neighbourhood of the layer before it. This growing \textbf{receptive field} is why early convolutional layers tend to respond to local structure --- edges, small textures --- while later ones respond to larger, more composite structure built from what the earlier layers detected.

\subsection{Pooling and spatial reduction}
\label{sec:pooling}

A \textbf{pooling} operation reduces the spatial size of a feature map without learning any weights of its own: it slides a small window across the map, exactly as a kernel does, but instead of a weighted sum it reports a fixed summary of the values it covers --- most commonly the maximum, giving max pooling. Pooling deliberately discards information: a $2\times2$ max-pool halves both height and width, keeping only the strongest response in each small neighbourhood and throwing the rest away. What is bought in exchange is a lower computational cost for every layer that follows, since they now operate on a smaller map, and a larger receptive field per unit of depth, since a later layer's kernel now covers a coarser, wider-reaching grid of the original input for the same kernel size. Pooling is used where the exact position of a detected feature matters less than its coarser location and the network can no longer afford to carry full spatial detail forward.

\subsection{Parameters, capacity and small data}
\label{sec:capacity}

The total count of a network's parameters --- every weight and bias, across every layer --- determines its \textbf{capacity}: loosely, how large and how intricate a set of functions it is able to represent. A higher-capacity network can fit more complicated relationships between input and label, but capacity is measured against the data available to constrain it, not in isolation. With a training corpus of only a few hundred or few thousand labelled examples, a network with far more parameters than examples has enough freedom to fit each training example almost exactly, including whatever is peculiar to that specific example --- sensor noise, an unrepresentative pose, an idiosyncrasy of one subject --- rather than the pattern that generalises across examples. This failure mode is \textbf{overfitting}: training loss keeps falling because the network is, in effect, memorising its training set, while performance on examples it did not train on stops improving or gets worse. The property actually wanted is \textbf{generalisation} --- accurate prediction on data the network never saw during training --- and it cannot be read off the training loss at all, since a memorising network can drive that loss arbitrarily low. Detecting whether generalisation is actually happening requires a \textbf{held-out set}: a portion of the labelled data set aside from the start, never used to compute a gradient or update a parameter, on which the network's predictions are checked only after training. A gap between good performance on the training data and poor performance on the held-out set is the signature of overfitting; parameter count, relative to the number of available training examples, is what makes that gap likely in the first place. This is precisely the constraint that governs the size and construction of every architecture described next, in \autoref{sec:building-blocks}.


\section{Network Building Blocks}
\label{sec:building-blocks}

Chapter~3 argues why each of the following components was chosen and what its measured contribution was. This section stays one level below that argument, setting out the mechanism each component computes.

\subsection{Three-dimensional convolution and kernel shape}
\label{sec:conv3d}

A \texttt{Conv3d} layer generalises image convolution to a five-dimensional tensor of shape $(B, C, D, H, W)$ --- batch, channel, and three spatial-like axes, here depth $D$ (time), height $H$ and width $W$. A kernel of shape $(k_D, k_H, k_W)$ extends the sliding, weight-sharing computation of \autoref{sec:convolution} across all three non-channel axes at once, with stride and padding again governing the output's size along each.

The detail that matters most for what follows is the kernel's temporal extent, $k_D$. If $k_D > 1$, the filter reaches across multiple frames at every application, so a single output value is a genuine function of more than one time step --- a spatio-temporal filter in the literal sense. If $k_D = 1$, the filter at time step $t$ only ever reads input at time step $t$; no output value depends on any other frame. A stack of layers with kernel shape $(1, k_H, k_W)$ is therefore a two-dimensional convolution applied identically and independently to every frame, merely expressed using \texttt{Conv3d} operations so the tensor never leaves its $(C, D, H, W)$ layout. This is the case the model in this thesis uses throughout its convolutional stem: every kernel is shaped $(1, 3, 3)$, with padding $(0, 1, 1)$ preserving the spatial extent while the temporal axis is neither padded nor touched. The temporal dimension is carried through the stem unchanged in length; only $H$ and $W$ are affected, first by the convolutions' own padding and then by a spatial-only max-pool of shape $(1, 2, 2)$ that halves both.

\subsection{Normalisation and regularisation}
\label{sec:normalisation}

\texttt{BatchNorm3d} normalises each channel independently: for channel $c$ it computes a mean and variance over every other axis at once --- the batch dimension and all spatial-temporal positions $(D, H, W)$ --- and rescales,
\[
	\hat{x} = \frac{x - \mu_c}{\sqrt{\sigma_c^2 + \epsilon}}, \qquad y = \gamma_c \hat{x} + \beta_c,
\]
where $\gamma_c$ and $\beta_c$ are learned per-channel scale and shift parameters. Training uses batch statistics; inference uses a running estimate accumulated during training, so the layer behaves exactly as its two-dimensional counterpart does, only pooled over one further axis.

\texttt{Dropout3d} differs from ordinary \texttt{Dropout} in what unit it zeroes. Ordinary dropout masks individual scalar activations independently at random, which is a weak regulariser on a convolutional feature map, since neighbouring positions in $H$, $W$ and $D$ are highly correlated and the network can read a missing value off an adjacent one. \texttt{Dropout3d} instead zeroes an entire channel --- every $(D, H, W)$ position of it, for a given sample --- with the given probability, removing a whole feature detector rather than a single correlated pixel.

\texttt{LayerNorm} normalises across the feature dimension of a single position for a single sample, independent of the batch, which is why it is the default choice inside transformer blocks rather than \texttt{BatchNorm}. Here it appears inside every transformer encoder layer and once more at the head of the classifier, which applies it before a dropout and the final linear projection to the output classes.

\subsection{Parameter-free attention: SimAM}
\label{sec:simam}

SimAM~\cite{yang2021} re-weights a feature map without learning any new parameters, by scoring each neuron on how distinctive it is from its surrounding context and using that score directly as a multiplicative gate. Yang et al. define an energy $e^{*}$ that is \textit{lower} for a more distinctive neuron, so it is the \textbf{reciprocal} of that energy that serves as the gate. For a neuron with activation $x$ in a channel whose mean and variance are $\mu$ and $v$ (estimated by pooling over that channel's own spatial extent), that reciprocal has the closed form
\[
	\frac{1}{e^{*}(x)} = \frac{(x-\mu)^2}{4(v+\lambda)} + 0.5,
\]
where $\lambda$ is the module's single hyperparameter. The refined output is then $x \cdot \sigma\!\left(1/e^{*}(x)\right)$, a monotonic sigmoid gate applied element-wise, the sigmoid present only to bound the gate rather than to reshape it. What makes this mechanism unusual among attention modules is exactly what the formula shows: every quantity it needs --- the mean, the variance, the resulting score --- is read off the feature map itself, at inference as much as at training time, so no weight matrix, bias, or reduction layer is introduced anywhere. The gate is a fixed function of the activations it is applied to, not a learned one.

The implementation in this thesis follows that formula exactly, with $\lambda = 10^{-4}$, but estimates $\mu$ and $v$ by pooling over the entire spatio-temporal volume $(D, H, W)$ of a stream's feature map rather than a single frame's spatial extent alone, so a neuron's distinctiveness is judged against the whole clip.

\subsection{Self-attention}
\label{sec:self-attention}

Self-attention relates every position in a sequence to every other position by learned linear projection rather than a fixed spatial neighbourhood. Each input vector is projected three ways, into a query $Q$, a key $K$ and a value $V$; attention weights are the scaled dot product of queries against keys,
\[
	\text{Attention}(Q,K,V) = \text{softmax}\!\left(\frac{QK^\top}{\sqrt{d_k}}\right)V,
\]
with the $\sqrt{d_k}$ scaling present to stop the dot products, and hence the softmax, from saturating as the key dimension $d_k$ grows. Multi-head attention runs several such projections in parallel, each into a smaller subspace, concatenating their outputs before a final linear mixing layer, so different heads can attend to different relationships within the same sequence.

Because this is a weighted sum over the whole sequence with weights derived purely from content, it is invariant to the order of the input positions: permuting the sequence permutes the output identically, with nothing in the mechanism encoding \textit{where} a token sits. A positional encoding is added to each input vector to break that symmetry. Two schemes are common. A sinusoidal encoding assigns each position a fixed vector built from sine and cosine functions at geometrically varying frequencies across the embedding dimensions --- deterministic, requiring no training, and evaluable at sequence lengths never seen during training. A learned encoding instead assigns each position index its own trainable embedding, which can fit the data more closely but cannot extrapolate beyond the longest position trained on. The encoder in this thesis uses the fixed sinusoidal form, held as a constant buffer rather than a parameter.

A second placement choice concerns normalisation relative to each sublayer. Post-norm --- the original arrangement --- applies \texttt{LayerNorm} after the residual addition that follows each sublayer; pre-norm applies it beforehand, to the sublayer's input, so the residual path carries an un-normalised signal straight through the stack, which is generally reported to train more stably since gradients reach earlier layers unimpeded. This is the placement used here. Beyond what the encoder layer itself contributes, no additional residual or skip connection is introduced anywhere else in the architecture.

\subsection{The fallbacks}
\label{sec:fallbacks}

Each learned component here can be switched off, and each switch has a specific, near parameter-free replacement rather than simply removing that stage. With the convolutional stem off, every frame's motion channels are instead reduced by adaptive average pooling to a fixed $4\times4$ spatial grid regardless of input resolution, flattened per frame and passed through a single linear layer into the model's working dimension. With the transformer off, the sequence of per-frame feature vectors is instead collapsed by a plain mean over the time axis, with no parameters and no representation of frame order at all.

Both fallbacks add as little capacity as possible: the pooling operations introduce no parameters, and the one linear layer needed to match dimensions in the no-CNN case is negligible next to what it replaces. That near-absence of added capacity is what makes each fallback a fair control, isolating what the learned component contributes without also changing how much the network is capable of fitting.

Isolating a component fairly also depends on how the network is trained, and \autoref{sec:scarcity-skew} sets out the training regime held constant across every configuration.


\section{Learning Under Scarcity and Skew}
\label{sec:scarcity-skew}

This section sets out the training machinery --- the loss function, the sampling scheme, regularisation, and the optimiser --- at the level of what each mechanism computes and how it is configured here. Chapter~3 (Section~3.9) reviews the literature motivating these choices, gives the corpus's class counts, and reports the failure mode that follows from combining two of them carelessly; that narrative is not repeated below. What follows is the mechanism and the configuration actually used.

\subsection{What skew does to cross-entropy}
\label{sec:skew-cross-entropy}

Standard cross-entropy sums a per-example term $-\log p_t$, where $p_t$ is the model's predicted probability for the true class, weighting every example equally. Under a skewed class distribution that equal weighting is not equal in effect: a class's total gradient contribution is its example count times the per-example gradient, so a majority class dominates the sum purely by outnumbering the rest. A model facing this loss has a cheap route to a lower value --- predict close to the prior on ambiguous cases --- because that is correct often enough on the majority class to outweigh being wrong on every minority example. Loss can fall steadily while the model learns comparatively little about classes it rarely sees. The remainder of this section responds to that problem.

\subsection{Focal loss: re-weighting by difficulty, not by frequency}
\label{sec:focal-loss}

Focal loss (Lin et al., 2017) inserts a modulating factor in front of the log term:

\[
	\text{FL}(p_t) = -\alpha_t (1-p_t)^{\gamma} \log p_t.
\]

Zhao et al.~\cite{zhaoS2021} apply this formulation to micro-expression recognition and state the mechanism plainly: $\gamma$ is ``the balance factor for loss'' and $\alpha_t$ ``the weight balance factor for samples''. The distinction matters for what follows. The factor $(1-p_t)^{\gamma}$ shrinks toward zero as the model becomes confident and correct on an example, re-weighting by how \textit{difficult} it currently finds that example, irrespective of its class --- an easy majority-class example contributes almost nothing once learned; a hard example, of any class, keeps contributing. Nothing in that factor looks at how many examples of a class exist. Class frequency enters only through the separate $\alpha_t$ term, an explicit per-class multiplier ordinarily set near inverse class frequency.

The implementation here (\texttt{Ablation\_Study/losses.py}) follows this exactly: it computes \texttt{log\_softmax}, gathers $\log p_t$ for the true class, exponentiates to recover $p_t$, and forms \texttt{focal\_weight = (1 - p\_t) ** gamma}. Label smoothing (\autoref{sec:label-smoothing}) is applied first, blending the negative log-likelihood term with a uniform term over classes, so the focal weight multiplies the \textit{smoothed} loss rather than the raw NLL. The optional $\alpha$ vector, where supplied, multiplies the result afterwards, indexed by each example's true class. The value of $\gamma$ used in this study follows Zhao et al.'s reported setting; Section~4.5.1 gives it, and the label-smoothing value, together with the rest of the training configuration.

\subsection{Balanced sampling}
\label{sec:balanced-sampling}

The second lever acts on the data a batch is drawn from rather than on the loss computed over it. A \texttt{WeightedRandomSampler} assigns each training example a weight equal to the inverse of its class's count, computed over the training split of the current fold only. Sampling then proceeds with replacement, drawing \texttt{num\_samples = len(train)} indices per epoch, so each class is presented with roughly equal probability regardless of its true rarity. This sampler is attached to the training loader alone; the validation loader uses no sampler and iterates with \texttt{shuffle=False}, so validation scores reflect the split's true distribution.

\subsection{Why combining both is a choice, not a default}
\label{sec:combining-both}

Focal loss's $\alpha$ term and the balanced sampler both correct for class frequency, at different points in the pipeline, so running both at once corrects for the same imbalance twice. A pipeline that applies both is therefore making a choice, not accepting a default, and the choice has to be resolved somewhere explicit. Section~4.5.2 states how this study resolves it and which of the two corrections is left active. Section~3.9.7 (cross-referenced rather than repeated here) explains why this rule exists and what happens when it is violated; the fact to carry forward is simply that the two mechanisms are not additive by default in this codebase.

\subsection{Label smoothing}
\label{sec:label-smoothing}

Label smoothing replaces the one-hot target with a softened one, mixing a fraction $\epsilon$ of uniform probability mass across all classes into the true-class target. Inside the focal-loss implementation this is a weighted combination of the ordinary negative log-likelihood term and a term averaging $-\log p_c$ over every class $c$; Section~4.5.1 gives the value of $\epsilon$ used here. The effect is to discourage the model from driving the true-class logit arbitrarily high --- regularisation against over-confidence, applied uniformly regardless of class, and distinct from a correction for imbalance. AdamW, mixed precision and gradient clipping have no dedicated treatment in this project's \texttt{docs/} corpus and are used here as standard practice.

\subsection{Augmentation under scarcity}
\label{sec:augmentation}

With few labelled clips per class, augmentation acts before the loss or the sampler is reached at all, and is applied to the training split only. Xia et al.~\cite{xia2020a} motivate this for micro-expression data directly: ``temporal data augmentation strategies as well as a balanced loss are jointly used for our deep network'' to address ``limited and imbalanced training samples''. Two transformations are implemented (\texttt{Ablation\_Study/dataset.py}), but only one of them acts. A random temporal crop of the input window (centred at evaluation) fires only where a clip's stored length exceeds the length the model expects; extraction stores exactly the number of frames the model is configured to consume (\autoref{sec:three-channel-tensor}), so that condition never holds and the crop is inert throughout this study. What remains is a horizontal flip applied with probability one half. Because the input is optical flow rather than raw pixels (\autoref{sec:video-to-motion}), the flip cannot be a plain mirror --- reversing the scene horizontally reverses the sign of horizontal displacement --- so the flip is paired with a negation of the flow tensor's $u$ (horizontal) channel. Flipping without negating $u$ would hand the network motion vectors pointing the wrong way for the geometry shown, a quiet corruption motion-tensor augmentation must specifically guard against.

\subsection{Optimisation}
\label{sec:optimisation}

Training uses AdamW, with decoupled weight decay applied uniformly to every parameter --- no separate no-decay group for biases or normalisation terms. The schedule combines a short linear warmup with cosine annealing: \texttt{LinearLR} ramps the rate from a fraction of target up to full over the warmup epochs, after which \texttt{CosineAnnealingLR} anneals it to a small floor over the remainder; the two are joined by \texttt{SequentialLR} and stepped once per epoch. Section~4.5.4 gives the learning rate, weight decay, warmup length and epoch count used here. Cosine annealing is adopted elsewhere in the micro-expression literature on similar small-data grounds --- Zhang et al.~\cite{zhang2022} use it in a spatio-temporal transformer for this same task.

Gradient norms are clipped after the AMP scaler's gradients are unscaled, via \texttt{clip\_grad\_norm\_}, guarding against the occasional large update a small, skewed dataset can produce. Training runs under mixed precision (\texttt{torch.amp.GradScaler}, \texttt{autocast}), enabled where CUDA is available and not disabled by flag. In the same unscaled branch, and therefore only while clipping is on, every parameter's gradient is checked for non-finite values before the optimiser step; if any is found, the step is skipped entirely, gradients are zeroed, and a warning logged, rather than letting a corrupted update reach the weights. Model selection keeps the checkpoint with the highest validation macro F1, not accuracy, consistent with the metric argument of Section~3.9.5 that accuracy is a poor guide under skew; ties favour the later epoch. There is no early stopping: every configuration trains its full epoch budget regardless of when its best epoch occurred.

One category of technique is absent by design: quantisation, pruning, and knowledge distillation are used nowhere in this project. Efficiency is measured rather than engineered here --- Section~3.6 reports the compute cost of the architectural choices actually made.

Measuring rather than engineering a property requires an evaluation apparatus equal to the task; \autoref{sec:evaluation} builds the one this corpus demands, from the confusion matrix upward.


\section{Evaluating on a Small Corpus}
\label{sec:evaluation}

Chapter~3 argues why leave-one-subject-out evaluation and pooled macro F1 suit this corpus, and what those choices cost here (Sections~3.1.6 and~3.1.8). This section gives the mechanism underneath: the confusion matrix and per-class precision, recall, F1; averaging scores versus pooling counts; the easily conflated distinction between averaging across folds and within them; what cross-validation buys and what leave-one-subject-out protects against; and, to close, an audit of what this thesis's own measurement apparatus does and does not supply.

\subsection{The confusion matrix, and precision, recall, F1 for one class}
\label{sec:confusion-matrix}

For a single class $c$ treated as positive against all others, every prediction falls into one of four counts: true positives $TP_c$ (correctly predicted $c$), false positives $FP_c$ (predicted $c$ but truly something else), false negatives $FN_c$ (truly $c$ but predicted otherwise), and true negatives. Arranging all classes this way at once gives the confusion matrix: entry $(i,j)$ counts clips whose true label is $i$ and predicted label is $j$, so the diagonal holds correct predictions and every off-diagonal cell names a specific error --- which true class is mistaken for which other.

From those counts, precision and recall isolate two failure modes:
\[
	\text{Precision}_c = \frac{TP_c}{TP_c + FP_c}, \qquad \text{Recall}_c = \frac{TP_c}{TP_c + FN_c}.
\]
Precision penalises false alarms: it falls whenever the classifier calls $c$ on a clip that is not $c$, regardless of how many true $c$ instances it also catches. Recall penalises misses: it falls whenever a genuine $c$ instance is assigned elsewhere, regardless of how many false alarms accompany the hits. Neither failure mode is visible from the other, which is why F1, the harmonic mean,
\[
	F1_c = \frac{2 \cdot \text{Precision}_c \cdot \text{Recall}_c}{\text{Precision}_c + \text{Recall}_c} = \frac{2\,TP_c}{2\,TP_c + FP_c + FN_c}
\]
is the usual single-number summary: being harmonic, it stays low whenever either input is low, so a model cannot buy a good F1 by trading recall for precision.

\subsection{Macro versus micro averaging}
\label{sec:macro-micro}

Combining per-class F1 scores into one number can be done two ways. Macro averaging computes $F1_c$ per class and takes an unweighted mean, $\text{Macro-F1} = \tfrac1C\sum_c F1_c$: every class counts equally regardless of size, so ignoring a rare class is penalised as heavily as ignoring a common one. Micro averaging instead pools the raw counts across classes first --- summing every class's $TP$, $FP$ and $FN$ --- and computes one F1 from those totals.

For single-label multi-class classification, where every clip gets exactly one predicted label, this pooling has a consequence worth stating plainly: every misclassification is simultaneously a false positive for the class it was wrongly assigned to and a false negative for its true class, so summed $FP$ and summed $FN$ equal each other and the total error count, and pooled micro-F1 reduces algebraically to $TP_{\text{total}}/N$ --- plain accuracy. Micro-F1 and accuracy are not two different metrics here; they are the same number computed two ways. The contrast that actually matters is therefore not ``macro versus micro'' but \textbf{macro-F1 versus accuracy}: an unweighted per-class average against a count a majority class can dominate.

\subsection{Pooled versus per-fold averaging}
\label{sec:pooled-vs-fold}

A further averaging choice sits underneath macro-F1 once evaluation is split into folds, and it governs how every result in this thesis is reported. One option accumulates the confusion counts --- $TP_c$, $FP_c$, $FN_c$ --- across every fold first, then computes one set of per-class F1 scores, and one macro-F1, from that pooled total. The other computes a complete macro-F1 inside each fold from that fold's own predictions, then averages the per-fold values. See et al.~\cite{see2019} define the Unweighted F1 (UF1) of the MEGC 2019 protocol the first way: the macro-averaged F1 obtained by accumulating true positives, false positives and false negatives over all folds of a leave-one-subject-out run before computing and averaging per-class F1.

These are estimators of two different quantities, not two ways of writing the same one: pooled macro-F1 describes this class's F1 across the whole evaluation, once every held-out prediction sits in one confusion matrix together, while per-fold-averaged macro-F1 describes the typical per-fold macro-F1 --- a quantity that depends on how classes happen to fall within each fold, not only on the classifier. When a fold's held-out set does not contain every class, its per-class F1 for the missing class is degenerate, and averaging such folds in with fully-populated ones changes what the number reflects for reasons unrelated to classifier quality; on a small, unevenly distributed corpus this is a live possibility, and fold composition can bound the per-fold estimator below what the pooled one would report. Sections~3.1.6 and~3.1.8 work through how the fold structure of this study does exactly that; which one a reported number is becomes apparent only once its computation is stated.

\subsection{Cross-validation and subject-disjointness}
\label{sec:cross-validation}

$k$-fold cross-validation partitions the data into $k$ disjoint subsets, trains $k$ models --- each on $k-1$ folds --- and evaluates each on the one fold it never saw, so every example is scored exactly once by a model that never trained on it. Leave-one-subject-out (LOSO) is the case where $k$ equals the number of subjects, and each fold's held-out set is precisely one subject's clips, every other subject's clips forming that fold's training set.

The property this buys is subject-disjointness: no fold's training and validation sets ever share a subject. Without it, a model evaluated on a clip from a subject it has already trained on can succeed by recognising that person's face or recording conditions rather than the expression class itself --- an identity leak inflating the score without the model having learned anything transferable to an unseen person. Enforcing subject-disjointness in every fold closes that channel.

\subsection{What an evaluation apparatus of this kind can fail to provide}
\label{sec:apparatus-gaps}

Three shortcomings belong to an evaluation apparatus rather than to any model it measures, and each matters more at this corpus size than it would at a larger one. They are set out here as properties to look for; Section~4.6 reports where this study's apparatus stands on each.

The first is \textbf{which aggregate a pipeline actually stores}. \autoref{sec:pooled-vs-fold} distinguishes pooling the confusion matrices before computing a macro average from averaging per-fold macro scores, and shows the two answer different questions. A pipeline may compute one, the other, or both under the same field name, and the quantity a reader wants is not necessarily the one written to disk.

The second is \textbf{whether model selection is blind to the data it is scored against}. A nested design keeps three sets apart: one to fit parameters, a second to choose between candidate models or epochs, and a third, untouched until the end, to score the chosen one. Collapsing the second and third --- selecting a checkpoint on the same data that then reports the result --- leaves the reported figure not blind to what it is checked against, and biases it optimistically by an amount that cannot be recovered after the fact.

The third is \textbf{whether enough is retained to support inference}. Aggregate metrics answer what a configuration scored; they cannot answer whether two configurations differ by more than chance. A paired test between configurations needs the per-clip predictions of each, so an apparatus that persists only aggregates forecloses that question whatever its scores say.


\section{Conclusion}
\label{sec:ch2-conclusion}

This chapter has set out the mechanisms on which the rest of the thesis depends, in the order the pipeline applies them. \autoref{sec:phenomenon} defined the micro-expression by the three properties that make it hard to recognise automatically --- it is brief, it is faint, and it cannot be produced on instruction --- and fixed the vocabulary of onset, apex and offset, noting that only the emotion label, not the action-unit coding, is used here. \autoref{sec:recording-to-input} established what actually reaches the pipeline: a 200~fps recording, the corpus's original frames rather than its registered release, and a greyscale frame at a fixed resolution, with no detection, registration or cropping performed by this project. \autoref{sec:motion-magnification} and \autoref{sec:video-to-motion} covered the two transformations applied to those frames --- Eulerian magnification of small motions, and the conversion of the magnified sequence into dense optical flow and optical strain.

\autoref{sec:nn-fundamentals} then built the machinery of a neural network from a single unit up to the convolutional layer, ending with the relationship between parameter count and corpus size that constrains every architectural choice this thesis makes. \autoref{sec:building-blocks} described the specific blocks assembled on that foundation, together with the parameter-free replacement each one falls back to when switched off. \autoref{sec:scarcity-skew} covered training under a class distribution that cannot be balanced by collecting more data, and \autoref{sec:evaluation} the evaluation apparatus needed when a corpus is small enough that the choice of averaging changes the conclusion.

Two threads run through all of it. The first is that the size of the corpus, rather than any modelling preference, forces most of the design: it dictates how few parameters the network can carry, why corrections for class skew are necessary, and why the evaluation protocol has to be chosen with care. The second is that several stages depart from the textbook or reference form of the technique they implement, and each departure has been named at the point it occurs rather than left for a reader to discover.

What this chapter has deliberately not done is say who else has used any of these techniques, with what result, or what remains unresolved. That is the work of Chapter~3, which reviews the literature for each component in turn and identifies the gap this thesis addresses.

%
% In-chapter \autoref targets resolve once all nine sections are present.
% Cross-chapter references (Chapter 3, Chapter 4, Section 3.1.2, ...) are
% written as plain text until those chapters carry labels.
%
% Citation keys used in this chapter:
%   yan2014     [1]  Yan et al., CASME II
%   qu2016      [2]  Qu et al., CAS(ME)^2
%   liY2018     [3]  Li, Huang & Zhao, single apex frame
%   bai2021     [4]  Bai et al., magnification + pre-trained network
%   liong2019a  [5]  Liong et al., OFF-ApexNet
%   shreve2011  [6]  Shreve et al., spatio-temporal strain
%   zhaoS2021   [7]  Zhao et al., two-stage 3D CNN
%   yang2021    [8]  Yang et al., SimAM
%   xia2020a    [9]  Xia et al., STRCN
%   zhang2022  [10]  Zhang et al., SLSTT
%   see2019    [11]  See et al., MEGC 2019
% =============================================================================

\chapter{Background}
\label{ch:background}

This chapter sets out the mechanisms the rest of the thesis depends on, in the order the pipeline applies them: what a micro-expression is, what the recording supplies, how small motions are magnified and converted into a motion representation, what a neural network does, which blocks this system assembles from that foundation, how it is trained under a corpus that cannot be balanced, and how it is measured when that corpus is small.

It explains \textbf{what each mechanism is and how it works}. Chapter~3 surveys who has used each one, with what result, and what gap remains; where a fact is needed in both places it is established here and cross-referenced there. Scope is limited to what the implemented system actually does --- techniques a reader might expect in a background chapter but which this project does not use are omitted, and their absence is noted where it would otherwise be assumed.

Two constraints recur. The corpus is small enough that its size, rather than any modelling preference, forces most of the design. And several stages depart from the textbook form of the technique they implement; each departure is named where it occurs.


\section{The Phenomenon: Micro-Expressions as Involuntary, Brief, Low-Intensity Facial Movement}
\label{sec:phenomenon}

A micro-expression is a brief facial movement that leaks a felt emotion the person is actively trying to conceal~\cite{yan2014}. It is best understood by contrast with an ordinary, or macro-, expression, on three axes rather than one. Macro-expressions typically last upwards of half a second, and may run to several seconds, and can be produced at will; micro-expressions, by contrast, are ``characterized by short durations, involuntary generation and low intensity''~\cite{qu2016}. The third axis is the one most easily overlooked: a micro-expression cannot be performed on instruction. It is a symptom of concealment rather than a communicative act. That single property governs how such data can be gathered at all, and the consequences for corpus construction are taken up in Section~3.1.

\subsection{Duration: the first constraint}
\label{sec:duration}

The field's working definition fixes an upper bound on duration. Yan et al.~\cite{yan2014} state that ``the generally accepted upper limit of the duration is 1/2 s.'' Brevity is the first source of the difficulty in recognising the phenomenon automatically, and the mechanism is simple: whatever signal distinguishes one emotion from another must be carried by however many video frames fit inside that half-second window. At a modest frame rate that window may contain only a handful of frames, several of which are near-identical to their neighbours; the sequence available to a classifier is short by construction, not by any deficiency of the recording.

\subsection{Low intensity: the second constraint}
\label{sec:low-intensity}

Duration alone would still leave a tractable problem if the movement were large. It is not. Yan et al.~\cite{yan2014} observe that micro-expressions are ``usually low in intensity -- it might be so brief for the facial muscles to become fully-stretched with suppression'', with the consequence that ``because of the short duration and low intensity, it is usually imperceptible or neglected by the naked eyes''. In practice the movement of interest may change only a few grey levels over a small patch of skin --- around an eyebrow, or the corner of a mouth --- while the rest of the face stays static. This is a distinct obstacle from brevity: a longer recording does nothing to help it, because the difficulty is not how many frames are available but how far the signal in each frame sits above the noise floor of ordinary video. Duration limits how much evidence exists; intensity limits how visible that evidence is once it does.

\subsection{Onset, apex and offset}
\label{sec:onset-apex-offset}

Because a micro-expression is a movement rather than a static configuration, it is conventionally described by three temporal landmarks. The \textbf{onset} is the frame at which the face begins to depart from its neutral baseline; the \textbf{offset} is the frame at which it returns to neutral; and the \textbf{apex} is the frame at which the movement is judged most intense --- the point at which, as Ekman puts it, a ``snapshot taken at an [sic] point when the expression is at its apex can easily convey the emotion message'' (as reported by Li, Huang \& Zhao~\cite{liY2018}). All three landmarks are available in the corpus used here. This thesis, however, uses only the onset and offset frames to bound each clip passed into the pipeline; the apex frame is never read. This is a deliberate design choice rather than an oversight, and its consequences --- for what information the model can and cannot see, and for how its results relate to apex-frame methods in the literature --- are taken up in Chapters~3 and~4.

\subsection{Action units, and the label actually used}
\label{sec:action-units}

The Facial Action Coding System (FACS) describes any facial movement, however subtle, as a combination of discrete Action Units (AUs), ``an objective method for labeling facial movements in terms of component actions''~\cite{yan2014}. An AU is a description of movement, not of emotion, and the mapping between the two is not one-to-one; corpora therefore assign emotion categories using AUs together with other evidence, by a procedure that differs between datasets and is described for the corpus used here in Section~3.1.2.

It is worth being precise about how much of that machinery survives into this thesis: none of it. The work reported here uses only the resulting emotion label. The Action Units recorded for each clip are present in the label table as inherited metadata from the original annotation --- carried along because the released coding includes them --- but they are never read by any stage of the data pipeline and never presented to the model. Nothing in this thesis performs, or claims to perform, action-unit recognition; the sole target throughout is the emotion category.

How brief and how faint that movement is determines what the recording must capture, and \autoref{sec:recording-to-input} turns to the video this pipeline actually receives.


\section{From Recording to Pipeline Input}
\label{sec:recording-to-input}

Before any mechanism of interest --- magnification, flow, strain, or a trained network --- can act on a micro-expression, something has to capture it, and something has to hand it to the pipeline in a usable form. This section describes that handoff: what the recording had to provide given the constraints fixed in \autoref{sec:phenomenon}, what the corpus actually supplies, and what this project does and does not do to a clip before it enters processing.

\subsection{High-speed capture}
\label{sec:high-speed-capture}

\autoref{sec:duration} fixed a half-second upper bound on the duration of a micro-expression. That bound has an immediate consequence for recording: whatever number of frames fit inside the event is the entire evidence a classifier will ever see of it, so the frame rate at which the video was captured is not a matter of picture quality but a precondition for the signal existing in the data at all. A 300~ms movement recorded at 25 or 30~fps --- ordinary video rates --- yields only seven or eight frames. The same event recorded at 200~fps yields roughly sixty frames, enough for onset, apex and offset to be distinct, separately observable moments rather than a blur collapsed into a handful of frames.

CASME II was recorded at 200~fps for exactly this reason~\cite{yan2014}. The pipeline used in this thesis takes 200 as a configured constant (\texttt{fps: int = 200} in \texttt{CASMEIIConfig}, \texttt{Stage1\_DataPipeline/config.py}), rather than reading it from the video files themselves, since it is a fixed property of the corpus rather than something to be inferred per clip.

\subsection{What the corpus supplies, and what this project does not do}
\label{sec:corpus-supplies}

CASME II is distributed in two forms: the frames as they were recorded, and a normalised release in which each face has been cropped and registered. This pipeline reads the first. Its metadata-unification stage resolves each clip to a per-subject, per-clip folder of the original recorded frames (\texttt{Stage1\_DataPipeline/metadata\_unifier.py}), so what enters processing is the full recorded field of view at 640~$\times$~480 (Section~3.1.2), within which the face occupies only a part.

It is worth stating plainly what this means, because the alternative assumption is the natural one to make. This project performs no face detection, no landmark localisation, no geometric registration, no alignment and no cropping of its own, anywhere in the pipeline --- and it does not inherit those operations from the corpus either, because it does not read the release that carries them. The normalisation that produced that release --- 68-point Active Shape Model landmarking on each clip's first frame, registered to a neutral model face by a Local Weighted Mean transform, described in Chapter~3, Section~3.1.2 --- is the corpus authors' work, and Section~3.1.2 records that this thesis departs from it. The face reaching the network is therefore an unregistered one, at whatever position and scale the camera recorded it, and whatever tolerance to that variation the model acquires it must acquire from the data.

\subsection{Frame selection and preparation}
\label{sec:frame-preparation}

Each frame is read in greyscale and resized to 224~$\times$~224 by bilinear interpolation. The resize is applied unconditionally to the whole frame, so it both downsamples and, since 640~$\times$~480 is not square, does not preserve the recorded aspect ratio. From the frames spanning a clip's annotated onset to its annotated offset, a fixed number of frames is selected for further processing --- the scheme by which that selection is made is described in \autoref{sec:three-channel-tensor}, together with the tensor it produces.

\subsection{Which clips enter the pipeline}
\label{sec:clip-selection}

Not every row of the corpus's coding table reaches processing. Four conditions are applied to select the rows that do: the row must belong to the CASME II dataset, its coded expression type must be \texttt{micro-expression} rather than one of the other categories the table records, its frame folder must be present on disk and non-empty, and its sequence length --- the onset-to-offset frame count, counted inclusively --- must exceed two. The first two conditions restrict the pipeline to the corpus and the expression class this thesis addresses, and the third discards rows whose media the coding table lists but the filesystem does not supply. The fourth rules out clips too short to be usable: a sequence of two frames or fewer contains at most one frame-to-frame transition, which is not enough to construct a motion sequence of any length, whatever downstream representation is built from it.

Together, these properties of the recording and the corpus --- a 200~fps capture rate, an unregistered full recorded frame, a greyscale 224~$\times$~224 resize of it, and a filtered, onset-to-offset clip --- are what reaches the pipeline before any processing of interest begins. Eulerian motion magnification is the first such process applied to these frames, and \autoref{sec:motion-magnification} describes the mechanism by which it amplifies motion without ever tracking it.


\section{Motion magnification}
\label{sec:motion-magnification}

\subsection{The Eulerian idea}
\label{sec:eulerian-idea}

A moving object can be described in two ways. A \textit{Lagrangian} description follows a point on the object as it moves --- this is what optical flow and feature tracking do: locate a point in frame $t$, find its correspondence in frame $t+1$, and report the displacement. A \textit{Eulerian} description instead fixes attention on a location in the image and asks how the signal observed there changes over time. Eulerian video magnification (EVM) takes the second view: no correspondence between frames is computed. At every pixel, it treats the sequence of intensity values over time as a 1-D signal, isolates the part occurring in a chosen temporal frequency band, multiplies it by a factor $\alpha$, and adds the result back to the original sequence.

The justification for why this amplifies \textit{motion}, and not merely intensity, is a first-order argument. If a small one-dimensional image profile $f(x)$ translates by a time-varying displacement $\delta(t)$, the observed intensity at a fixed location is $I(x,t) = f(x+\delta(t))$. A first-order Taylor expansion around $x$ gives

\[
	I(x,t) \approx f(x) + \delta(t)\,\frac{\partial f(x)}{\partial x},
\]

so the temporal variation at a fixed pixel is approximately proportional to the true displacement $\delta(t)$, scaled by the local spatial gradient. Amplifying that temporal variation by $\alpha$ before adding it back is therefore, to first order, equivalent to synthesising a signal consistent with a displacement of $(1+\alpha)\delta(t)$ rather than $\delta(t)$: motion is amplified as a consequence of amplifying intensity change, with no tracking step anywhere in the process. This is the account given by Bai et al.~\cite{bai2021} (following Wu et al.).

\subsection{Spatial decomposition: the Laplacian pyramid}
\label{sec:laplacian-pyramid}

The Taylor approximation above only holds where the spatial gradient is well-behaved relative to the displacement, so in practice the frame is first decomposed into a set of spatial-frequency bands and the argument is applied band-by-band. The standard tool is the \textbf{Laplacian pyramid}. Each level is built by blurring the current image with a small low-pass kernel, downsampling by a factor of two, upsampling the result back to the original resolution, and subtracting it from the pre-downsampling image. The subtraction discards everything the blur-and-downsample step could reconstruct, leaving only the spatial detail specific to that scale --- the Laplacian band. The downsampled image feeds the next iteration, and the recursion terminates in a coarse, heavily blurred residual (a Gaussian, not a Laplacian, band). A pyramid of depth $L$ yields $L$ Laplacian bands, each isolating a different range of spatial frequencies, plus one coarsest residual that carries no band-specific detail and exists only to seed reconstruction.

\subsection{Temporal filtering}
\label{sec:temporal-filtering}

Each spatial band is treated as a stack of scalar time series, one per pixel, and passed through a \textbf{temporal band-pass filter} that keeps only the variation falling inside a chosen frequency range and discards the rest. The filtered signal at each band is multiplied by $\alpha$ and added back to that band's original values; the bands are then recombined by repeating the pyramid construction in reverse --- upsampling the coarsest residual and successively adding each (now-amplified) Laplacian band back in. The published amplitude-based method realises the band-pass step with a \textbf{Butterworth filter}, a smooth infinite-impulse-response design~\cite{bai2021}.

\subsection{The amplification trade-off}
\label{sec:amplification-tradeoff}

The magnification factor $\alpha$ trades sensitivity against artefact. A larger $\alpha$ makes a fainter in-band motion visible, but the same multiplication applies indiscriminately to everything the band-pass filter passes: sensor noise, compression artefacts, and any other real motion whose temporal frequency happens to fall inside the chosen band are amplified exactly as the motion of interest is. The method has no way to distinguish a wanted signal from an unwanted one occupying the same frequency range --- the filter only knows how fast a signal oscillates, not what it is. Raising $\alpha$ therefore does not monotonically improve the result; past some point amplified noise and induced displacement artefacts dominate whatever gain in visibility was sought. This tension is intrinsic to the method rather than a defect of any one implementation, and it recurs, with numbers, in Chapter~3.

\subsection{Departures from the textbook method in this implementation}
\label{sec:evm-departures}

The magnifier used in this thesis (\texttt{Stage1\_DataPipeline/evm\_magnifier.py}) implements the mechanism above with a genuine Laplacian pyramid --- a $[1,4,6,4,1]/16$ blur kernel, downsample-by-two, upsample, subtract --- followed by a coarsest Gaussian residual, matching \autoref{sec:laplacian-pyramid} exactly. It departs from amplitude-based magnification as originally formulated in four respects; the third is a property of that original method rather than of the summary given by Bai et al.~\cite{bai2021}.

First, the temporal band-pass is an \textbf{ideal filter}: a hard binary mask applied directly to the discrete Fourier transform of each pixel's time series (via \texttt{scipy.fftpack}), rather than a Butterworth (IIR) filter. Frequencies inside $[\omega_{\mathrm{low}}, \omega_{\mathrm{high}}]$ pass with a gain of exactly one and everything else is zeroed, with no smooth roll-off at the band edges.

Second, every Laplacian band is amplified by the \textbf{same scalar $\alpha$}. Bai et al.~\cite{bai2021} state the amplified reconstruction with a \textit{frequency-dependent} factor $\alpha_k$ --- one per band, reduced at higher spatial frequencies to limit noise amplification --- and no such damping is applied here. The coarsest residual is never filtered or amplified --- it only seeds reconstruction.

Third, the implementation works directly on the greyscale intensity tensor, so there is no chromatic attenuation step.

Fourth, and a property of pipeline ordering rather than of the magnifier itself: the calling code (\texttt{tensor\_pipeline\_manager.py}) resamples each clip to a fixed 33 frames \textit{before} invoking the magnifier, and passes the corpus's nominal recording rate (200~fps) as the filter's frame-rate parameter regardless of the resampled sequence's true frame spacing. The consequence --- a temporal filter whose frequency axis no longer matches the signal it is filtering --- is analysed in Section~3.5.7, not here.

The operating point this study runs at is given in Section~4.2.5; whatever the amplification factor, band and pyramid depth, the amplified output is clipped to $[0,255]$ and quantised back to 8-bit before being handed to the downstream flow-and-strain extractor.

What that extractor does with the amplified sequence --- and why motion, rather than appearance, is what the network is given --- is the subject of \autoref{sec:video-to-motion}.


\section{From Video to Motion: Optical Flow and Optical Strain}
\label{sec:video-to-motion}

\subsection{Why motion, not pixels}
\label{sec:why-motion}

A micro-expression is a movement, so the natural quantity to hand a classifier is a description of movement rather than of appearance. Raw pixel intensity carries the signal of interest --- the same faint change described in \autoref{sec:low-intensity} --- embedded in everything else the camera also recorded: who the subject is, how the scene is lit, the fixed geometry of a particular face. None of that is discarded by presenting a network with frames directly, and on a corpus this narrow in identity and illumination a model has every incentive to fit the wrong part of the image. A displacement field is different in kind: it records how each point moved, not what colour it is, so identity and illumination do not enter it at all. That principle motivates the rest of this section; the comparative case for it is made in Section~3.3.1.

\subsection{Optical flow: the brightness-constancy constraint}
\label{sec:optical-flow}

Optical flow estimates, for every pixel, the apparent two-dimensional displacement between two frames. The estimate rests on the brightness-constancy assumption: a point on a moving surface is taken to have the same intensity in the next frame as in the current one, so that any change in a pixel's grey level is attributed to something having moved past that location, never to the thing itself changing colour or shade~\cite{liong2019a}.

Writing $I_t(x,y)$ for the image intensity at position $(x,y)$ and time $t$, and supposing the point there moves to $(x+\delta x,\, y+\delta y)$ by time $t+1$, brightness constancy states

\[
	I_t(x,y) = I_{t+1}(x+\delta x,\, y+\delta y), \qquad \delta x = u_t\,\delta t,\;\; \delta y = v_t\,\delta t,
\]

where $u_t(x,y)$ and $v_t(x,y)$ are the horizontal and vertical components of the flow at that point. Expanding the right-hand side by a first-order Taylor series and substituting into the constancy equation, the intensity terms cancel and division by $\delta t$ leaves the \textbf{optical flow constraint equation}:

\[
	u_t(x,y)\,\frac{\partial I}{\partial x} + v_t(x,y)\,\frac{\partial I}{\partial y} + \frac{\partial I}{\partial t} = 0.
\]

This single scalar equation holds at every pixel but contains two unknowns, $u_t$ and $v_t$: the image gradients are measurable directly from the frames, while the displacement is not determined by them alone. This under-determinacy is the \textbf{aperture problem} --- a local patch of image constrains the component of motion perpendicular to an edge or gradient, but says nothing about motion parallel to it. One equation per pixel cannot fix two unknowns per pixel, so every practical estimator closes the system by adding a further assumption relating the flow at neighbouring pixels, typically that the field varies smoothly across the image. That added assumption is what distinguishes one flow algorithm from another; this thesis does not adjudicate between them (Section~3.3.3) and simply adopts one, described in \autoref{sec:three-channel-tensor}. The output, for a pair of frames, is the pair of scalar fields $u(x,y)$ and $v(x,y)$ --- the horizontal and vertical displacement at every pixel --- forming the first two channels of the motion representation used throughout this thesis.

\subsection{Optical strain: the finite strain tensor}
\label{sec:optical-strain}

Flow answers where a point moved; it does not by itself distinguish a patch of skin that moved together with its neighbours from one that stretched or sheared relative to them. Optical strain makes that distinction by taking the spatial derivative of the flow field rather than the field itself.

Writing the displacement at a point as $\mathbf{u} = [u, v]^{\mathsf T}$, the \textbf{finite strain tensor} is defined as the symmetric part of the displacement gradient~\cite{shreve2011}:

\[
	\varepsilon = \tfrac{1}{2}\left[\nabla\mathbf{u} + (\nabla\mathbf{u})^{\mathsf T}\right] = \begin{bmatrix} \varepsilon_{xx} & \varepsilon_{xy} \\ \varepsilon_{yx} & \varepsilon_{yy} \end{bmatrix}, \qquad \varepsilon_{xx} = \frac{\partial u}{\partial x},\;\; \varepsilon_{yy} = \frac{\partial v}{\partial y},\;\; \varepsilon_{xy} = \varepsilon_{yx} = \frac{1}{2}\left(\frac{\partial u}{\partial y} + \frac{\partial v}{\partial x}\right).
\]

$\varepsilon_{xx}$ and $\varepsilon_{yy}$ are the \textbf{normal strain} components, describing stretching or compression along each axis; $\varepsilon_{xy}$ is the \textbf{shear strain} component, describing the change of angle between two initially perpendicular directions on the surface. The tensor is reduced to a single scalar per pixel. Shreve et al.~\cite{shreve2011} sum the component values to obtain a strain magnitude; the reduction implemented here is instead a root sum of squares, with the shear term counted once,

\[
	\varepsilon_{\text{mag}} = \sqrt{\varepsilon_{xx}^{2} + \varepsilon_{yy}^{2} + \varepsilon_{xy}^{2}},
\]

a form that diverges from the convention used across the reviewed corpus. Section~3.4.3 records that divergence and what it costs.

The property that makes this useful, rather than a second copy of the flow field, follows from its definition as a \textit{gradient}. If a region of the face translates rigidly --- every pixel moving by the same $u$ and the same $v$, as a small head movement produces --- then $u$ and $v$ are locally constant there, and their spatial derivatives, hence every component of $\varepsilon$, are zero: a rigid translation contributes nothing to the strain field, however large it is. A localised muscle contraction, by contrast, moves neighbouring points of skin by different amounts, so the displacement field has a non-zero gradient exactly where the deformation occurs, and $\varepsilon_{\text{mag}}$ is large there. Strain is thus sensitive to non-rigid deformation specifically, insensitive by construction to whatever rigid motion the flow field also contains.

\subsection{Assembling the three-channel tensor}
\label{sec:three-channel-tensor}

For each clip, the frames spanning onset to offset are uniformly resampled to a fixed sequence of 33 frames, and dense optical flow is computed between every adjacent pair using OpenCV's implementation of Farneb\"ack's method~\cite{zhaoS2021} --- \texttt{cv2.calcOpticalFlowFarneback}, with pyramid scale 0.5, three pyramid levels, window size 15, three iterations, a polynomial neighbourhood of size 5 and a polynomial smoothing of 1.2, applied to 8-bit greyscale frames. This yields \textbf{32 flow fields} per clip, each a $(u, v)$ pair over the 224~$\times$~224 face region. Strain is computed from each field's spatial gradients using \texttt{np.gradient} (central differences at unit pixel spacing, over the two spatial axes only; no temporal gradient is used), giving $\varepsilon_{xx}$, $\varepsilon_{yy}$, $\varepsilon_{xy}$ and hence $\varepsilon_{\text{mag}}$ for each of the 32 pairs.

The three quantities --- $u$, $v$ and $\varepsilon_{\text{mag}}$ --- are stacked as three channels, giving a tensor of shape $(3, 32, 224, 224)$ in channel order $[u, v, \text{strain}]$. When Eulerian video magnification is applied (\autoref{sec:motion-magnification}), flow and strain are computed on the magnified frames rather than the originals; the construction is otherwise unchanged.

Normalisation is applied in two separate stages, and both remain active. At extraction time, each of the three channels is independently min--max normalised to $[0, 1]$ across the whole clip, so that a channel's own range --- which differs enormously between a flow component and a strain magnitude --- does not by itself determine its influence. Separately, at training load time (\texttt{Ablation\_Study/dataset.py}, with \texttt{normalize=True}), each channel is z-scored per sample over its $T \times H \times W$ extent. Neither stage substitutes for the other: the first is a per-clip rescaling done once, during data preparation; the second is a per-sample standardisation applied every time a clip is loaded, centring and scaling each channel to zero mean and unit variance.

How this tensor is built is now settled; what consumes it is not. \autoref{sec:nn-fundamentals} begins with what a neural network is, before \autoref{sec:building-blocks} reaches the blocks themselves.


\section{Neural Network Fundamentals}
\label{sec:nn-fundamentals}

Everything in this thesis described as ``trained'' is, underneath, a neural network: a function built from simple, repeated computations whose numerical parameters are set from data rather than specified by hand. This section builds that machinery from its smallest unit upward, only as far as is needed for the architecture-specific mechanisms of \autoref{sec:building-blocks}, the training regime of \autoref{sec:scarcity-skew} and the evaluation apparatus of \autoref{sec:evaluation} to be legible when they introduce anything that depends on it.

\subsection{Units, layers and the forward pass}
\label{sec:units-layers}

The basic computational unit takes a vector of inputs $x$, forms a weighted sum of them plus an offset term, and passes the result through a fixed non-linear function $\phi$:
\[
	a = \phi(w^\top x + b).
\]
The entries of $w$ are the unit's \textbf{weights}, one per input; $b$ is its \textbf{bias}. Weights and biases together are the unit's \textbf{parameters} --- the numbers a training procedure is free to adjust --- and $a$, the unit's output, is its \textbf{activation}. A \textbf{layer} is a bank of such units applied to the same input vector in parallel, each with its own weights and bias, producing a vector of activations rather than a single scalar. A network is a sequence of layers, each layer's output vector becoming the next layer's input; evaluating the network on an input by applying each layer in turn, from the first to the last, is the \textbf{forward pass}, and its final output is the network's prediction. The number of layers is the network's \textbf{depth}; the number of units in a given layer is that layer's \textbf{width}. Both are choices made when the architecture is specified, not something training decides.

\subsection{Non-linearity}
\label{sec:non-linearity}

If $\phi$ were the identity function, every layer would compute only an affine transformation of its input, and the composition of any number of affine transformations is itself a single affine transformation. A network built entirely this way would therefore compute no more than a network with one layer could, however many layers it stacked --- depth would buy nothing. The non-linear function $\phi$, applied after the weighted sum at every unit, is what stops layers from collapsing into one another; it is the reason depth is a meaningful design axis at all rather than a redundant one. A common choice, used throughout the components described in this thesis, is the rectified linear unit, $\text{ReLU}(x) = \max(0, x)$: it passes positive values through unchanged and clips negative ones to zero. Its appeal is mechanical rather than subtle --- it is cheap to compute, its derivative is either zero or one almost everywhere, and unlike activation functions that saturate on both sides, a unit with a positive input keeps gradient flowing through it at full strength.

\subsection{Learning by gradient descent}
\label{sec:gradient-descent}

A network's forward pass is fixed once its parameters are fixed, so ``training'' means searching for parameter values under which the forward pass produces good predictions. A \textbf{loss function} quantifies how far a prediction sits from its true label, returning a single number that is small when they agree and large when they do not; the specific loss used in this thesis, and why, is a matter for \autoref{sec:scarcity-skew}, but the role a loss plays is a property of every trained network and belongs here. Training proceeds by adjusting parameters to make the loss smaller, and the mechanism that makes that adjustment tractable is \textbf{backpropagation}: because the forward pass is a composition of layers, the chain rule from calculus lets the gradient of the loss with respect to every parameter in the network --- however deep --- be computed by working backward from the output layer to the input, reusing intermediate results rather than recomputing each parameter's gradient from scratch. Once every parameter's gradient is known, an \textbf{update rule} moves each parameter a small step against its gradient, since the gradient points in the direction the loss increases fastest and moving the opposite way decreases it. The size of that step is the \textbf{learning rate}. Parameters are not updated from the whole training set at once in a single step; instead the data is divided into \textbf{batches}, and one update is made per batch, while one full pass through every batch in the training set is one \textbf{epoch}, of which training runs many in succession. Throughout, a network has two distinct modes of use: in training mode its parameters are being adjusted from labelled examples by the procedure just described, while in evaluation mode its parameters are held fixed and it is only asked to produce predictions, with no update following from how correct they are.

\subsection{Convolution as a learned filter}
\label{sec:convolution}

A layer of the kind described in \autoref{sec:units-layers} connects every input to every unit, which is wasteful for an image: a pattern worth detecting --- an edge, a corner, a small patch of texture --- can appear anywhere in the frame, and a fully connected layer would need to learn a separate copy of the detector for every position it might appear at. A convolutional layer instead learns one small filter, a \textbf{kernel}, consisting of a grid of weights (a $3\times3$ or $5\times5$ window is typical) plus a single bias, and slides that same kernel across every position of the input, at each position computing a weighted sum of the pixels currently under it. Because the identical kernel is reused at every position, the layer is said to exhibit \textbf{weight sharing}: it looks for the same pattern everywhere in the image and costs only as many parameters as the kernel itself holds, regardless of how large the input is. The result of sliding one kernel across the whole input is a single \textbf{feature map}, one value per position, recording how strongly that kernel's pattern was present there. A convolutional layer ordinarily learns several kernels at once, each producing its own feature map; stacked together these form the layer's output \textbf{channels}, so a layer with $C_{\text{out}}$ kernels produces $C_{\text{out}}$ channels from whatever number of input channels, $C_{\text{in}}$, it was given (three for an ordinary colour image, one for a greyscale image, or more once earlier layers have produced feature maps of their own). Because each kernel spans every input channel, not just one, the number of learned parameters in such a layer, for a kernel of height $k_H$ and width $k_W$, is
\[
	C_{\text{out}} \times \bigl(C_{\text{in}} \times k_H \times k_W + 1\bigr),
\]
the $+1$ accounting for one bias per output channel --- a formula this thesis returns to repeatedly when comparing the cost of alternative designs. Two further settings control how the kernel is applied. \textbf{Stride} is the number of positions the kernel moves between one application and the next; a stride greater than one skips positions and shrinks the output accordingly. \textbf{Padding} adds a border, typically of zeros, around the input before the kernel is applied, which controls how much the output shrinks relative to the input and lets a layer preserve spatial size exactly if that is wanted. Because a single kernel only ever looks at a small window, one convolutional layer alone only detects small, local patterns; stacking convolutional layers lets a unit in a later layer depend, indirectly, on a much larger region of the original input, since each layer's output already summarises a neighbourhood of the layer before it. This growing \textbf{receptive field} is why early convolutional layers tend to respond to local structure --- edges, small textures --- while later ones respond to larger, more composite structure built from what the earlier layers detected.

\subsection{Pooling and spatial reduction}
\label{sec:pooling}

A \textbf{pooling} operation reduces the spatial size of a feature map without learning any weights of its own: it slides a small window across the map, exactly as a kernel does, but instead of a weighted sum it reports a fixed summary of the values it covers --- most commonly the maximum, giving max pooling. Pooling deliberately discards information: a $2\times2$ max-pool halves both height and width, keeping only the strongest response in each small neighbourhood and throwing the rest away. What is bought in exchange is a lower computational cost for every layer that follows, since they now operate on a smaller map, and a larger receptive field per unit of depth, since a later layer's kernel now covers a coarser, wider-reaching grid of the original input for the same kernel size. Pooling is used where the exact position of a detected feature matters less than its coarser location and the network can no longer afford to carry full spatial detail forward.

\subsection{Parameters, capacity and small data}
\label{sec:capacity}

The total count of a network's parameters --- every weight and bias, across every layer --- determines its \textbf{capacity}: loosely, how large and how intricate a set of functions it is able to represent. A higher-capacity network can fit more complicated relationships between input and label, but capacity is measured against the data available to constrain it, not in isolation. With a training corpus of only a few hundred or few thousand labelled examples, a network with far more parameters than examples has enough freedom to fit each training example almost exactly, including whatever is peculiar to that specific example --- sensor noise, an unrepresentative pose, an idiosyncrasy of one subject --- rather than the pattern that generalises across examples. This failure mode is \textbf{overfitting}: training loss keeps falling because the network is, in effect, memorising its training set, while performance on examples it did not train on stops improving or gets worse. The property actually wanted is \textbf{generalisation} --- accurate prediction on data the network never saw during training --- and it cannot be read off the training loss at all, since a memorising network can drive that loss arbitrarily low. Detecting whether generalisation is actually happening requires a \textbf{held-out set}: a portion of the labelled data set aside from the start, never used to compute a gradient or update a parameter, on which the network's predictions are checked only after training. A gap between good performance on the training data and poor performance on the held-out set is the signature of overfitting; parameter count, relative to the number of available training examples, is what makes that gap likely in the first place. This is precisely the constraint that governs the size and construction of every architecture described next, in \autoref{sec:building-blocks}.


\section{Network Building Blocks}
\label{sec:building-blocks}

Chapter~3 argues why each of the following components was chosen and what its measured contribution was. This section stays one level below that argument, setting out the mechanism each component computes.

\subsection{Three-dimensional convolution and kernel shape}
\label{sec:conv3d}

A \texttt{Conv3d} layer generalises image convolution to a five-dimensional tensor of shape $(B, C, D, H, W)$ --- batch, channel, and three spatial-like axes, here depth $D$ (time), height $H$ and width $W$. A kernel of shape $(k_D, k_H, k_W)$ extends the sliding, weight-sharing computation of \autoref{sec:convolution} across all three non-channel axes at once, with stride and padding again governing the output's size along each.

The detail that matters most for what follows is the kernel's temporal extent, $k_D$. If $k_D > 1$, the filter reaches across multiple frames at every application, so a single output value is a genuine function of more than one time step --- a spatio-temporal filter in the literal sense. If $k_D = 1$, the filter at time step $t$ only ever reads input at time step $t$; no output value depends on any other frame. A stack of layers with kernel shape $(1, k_H, k_W)$ is therefore a two-dimensional convolution applied identically and independently to every frame, merely expressed using \texttt{Conv3d} operations so the tensor never leaves its $(C, D, H, W)$ layout. This is the case the model in this thesis uses throughout its convolutional stem: every kernel is shaped $(1, 3, 3)$, with padding $(0, 1, 1)$ preserving the spatial extent while the temporal axis is neither padded nor touched. The temporal dimension is carried through the stem unchanged in length; only $H$ and $W$ are affected, first by the convolutions' own padding and then by a spatial-only max-pool of shape $(1, 2, 2)$ that halves both.

\subsection{Normalisation and regularisation}
\label{sec:normalisation}

\texttt{BatchNorm3d} normalises each channel independently: for channel $c$ it computes a mean and variance over every other axis at once --- the batch dimension and all spatial-temporal positions $(D, H, W)$ --- and rescales,
\[
	\hat{x} = \frac{x - \mu_c}{\sqrt{\sigma_c^2 + \epsilon}}, \qquad y = \gamma_c \hat{x} + \beta_c,
\]
where $\gamma_c$ and $\beta_c$ are learned per-channel scale and shift parameters. Training uses batch statistics; inference uses a running estimate accumulated during training, so the layer behaves exactly as its two-dimensional counterpart does, only pooled over one further axis.

\texttt{Dropout3d} differs from ordinary \texttt{Dropout} in what unit it zeroes. Ordinary dropout masks individual scalar activations independently at random, which is a weak regulariser on a convolutional feature map, since neighbouring positions in $H$, $W$ and $D$ are highly correlated and the network can read a missing value off an adjacent one. \texttt{Dropout3d} instead zeroes an entire channel --- every $(D, H, W)$ position of it, for a given sample --- with the given probability, removing a whole feature detector rather than a single correlated pixel.

\texttt{LayerNorm} normalises across the feature dimension of a single position for a single sample, independent of the batch, which is why it is the default choice inside transformer blocks rather than \texttt{BatchNorm}. Here it appears inside every transformer encoder layer and once more at the head of the classifier, which applies it before a dropout and the final linear projection to the output classes.

\subsection{Parameter-free attention: SimAM}
\label{sec:simam}

SimAM~\cite{yang2021} re-weights a feature map without learning any new parameters, by scoring each neuron on how distinctive it is from its surrounding context and using that score directly as a multiplicative gate. Yang et al. define an energy $e^{*}$ that is \textit{lower} for a more distinctive neuron, so it is the \textbf{reciprocal} of that energy that serves as the gate. For a neuron with activation $x$ in a channel whose mean and variance are $\mu$ and $v$ (estimated by pooling over that channel's own spatial extent), that reciprocal has the closed form
\[
	\frac{1}{e^{*}(x)} = \frac{(x-\mu)^2}{4(v+\lambda)} + 0.5,
\]
where $\lambda$ is the module's single hyperparameter. The refined output is then $x \cdot \sigma\!\left(1/e^{*}(x)\right)$, a monotonic sigmoid gate applied element-wise, the sigmoid present only to bound the gate rather than to reshape it. What makes this mechanism unusual among attention modules is exactly what the formula shows: every quantity it needs --- the mean, the variance, the resulting score --- is read off the feature map itself, at inference as much as at training time, so no weight matrix, bias, or reduction layer is introduced anywhere. The gate is a fixed function of the activations it is applied to, not a learned one.

The implementation in this thesis follows that formula exactly, with $\lambda = 10^{-4}$, but estimates $\mu$ and $v$ by pooling over the entire spatio-temporal volume $(D, H, W)$ of a stream's feature map rather than a single frame's spatial extent alone, so a neuron's distinctiveness is judged against the whole clip.

\subsection{Self-attention}
\label{sec:self-attention}

Self-attention relates every position in a sequence to every other position by learned linear projection rather than a fixed spatial neighbourhood. Each input vector is projected three ways, into a query $Q$, a key $K$ and a value $V$; attention weights are the scaled dot product of queries against keys,
\[
	\text{Attention}(Q,K,V) = \text{softmax}\!\left(\frac{QK^\top}{\sqrt{d_k}}\right)V,
\]
with the $\sqrt{d_k}$ scaling present to stop the dot products, and hence the softmax, from saturating as the key dimension $d_k$ grows. Multi-head attention runs several such projections in parallel, each into a smaller subspace, concatenating their outputs before a final linear mixing layer, so different heads can attend to different relationships within the same sequence.

Because this is a weighted sum over the whole sequence with weights derived purely from content, it is invariant to the order of the input positions: permuting the sequence permutes the output identically, with nothing in the mechanism encoding \textit{where} a token sits. A positional encoding is added to each input vector to break that symmetry. Two schemes are common. A sinusoidal encoding assigns each position a fixed vector built from sine and cosine functions at geometrically varying frequencies across the embedding dimensions --- deterministic, requiring no training, and evaluable at sequence lengths never seen during training. A learned encoding instead assigns each position index its own trainable embedding, which can fit the data more closely but cannot extrapolate beyond the longest position trained on. The encoder in this thesis uses the fixed sinusoidal form, held as a constant buffer rather than a parameter.

A second placement choice concerns normalisation relative to each sublayer. Post-norm --- the original arrangement --- applies \texttt{LayerNorm} after the residual addition that follows each sublayer; pre-norm applies it beforehand, to the sublayer's input, so the residual path carries an un-normalised signal straight through the stack, which is generally reported to train more stably since gradients reach earlier layers unimpeded. This is the placement used here. Beyond what the encoder layer itself contributes, no additional residual or skip connection is introduced anywhere else in the architecture.

\subsection{The fallbacks}
\label{sec:fallbacks}

Each learned component here can be switched off, and each switch has a specific, near parameter-free replacement rather than simply removing that stage. With the convolutional stem off, every frame's motion channels are instead reduced by adaptive average pooling to a fixed $4\times4$ spatial grid regardless of input resolution, flattened per frame and passed through a single linear layer into the model's working dimension. With the transformer off, the sequence of per-frame feature vectors is instead collapsed by a plain mean over the time axis, with no parameters and no representation of frame order at all.

Both fallbacks add as little capacity as possible: the pooling operations introduce no parameters, and the one linear layer needed to match dimensions in the no-CNN case is negligible next to what it replaces. That near-absence of added capacity is what makes each fallback a fair control, isolating what the learned component contributes without also changing how much the network is capable of fitting.

Isolating a component fairly also depends on how the network is trained, and \autoref{sec:scarcity-skew} sets out the training regime held constant across every configuration.


\section{Learning Under Scarcity and Skew}
\label{sec:scarcity-skew}

This section sets out the training machinery --- the loss function, the sampling scheme, regularisation, and the optimiser --- at the level of what each mechanism computes and how it is configured here. Chapter~3 (Section~3.9) reviews the literature motivating these choices, gives the corpus's class counts, and reports the failure mode that follows from combining two of them carelessly; that narrative is not repeated below. What follows is the mechanism and the configuration actually used.

\subsection{What skew does to cross-entropy}
\label{sec:skew-cross-entropy}

Standard cross-entropy sums a per-example term $-\log p_t$, where $p_t$ is the model's predicted probability for the true class, weighting every example equally. Under a skewed class distribution that equal weighting is not equal in effect: a class's total gradient contribution is its example count times the per-example gradient, so a majority class dominates the sum purely by outnumbering the rest. A model facing this loss has a cheap route to a lower value --- predict close to the prior on ambiguous cases --- because that is correct often enough on the majority class to outweigh being wrong on every minority example. Loss can fall steadily while the model learns comparatively little about classes it rarely sees. The remainder of this section responds to that problem.

\subsection{Focal loss: re-weighting by difficulty, not by frequency}
\label{sec:focal-loss}

Focal loss (Lin et al., 2017) inserts a modulating factor in front of the log term:

\[
	\text{FL}(p_t) = -\alpha_t (1-p_t)^{\gamma} \log p_t.
\]

Zhao et al.~\cite{zhaoS2021} apply this formulation to micro-expression recognition and state the mechanism plainly: $\gamma$ is ``the balance factor for loss'' and $\alpha_t$ ``the weight balance factor for samples''. The distinction matters for what follows. The factor $(1-p_t)^{\gamma}$ shrinks toward zero as the model becomes confident and correct on an example, re-weighting by how \textit{difficult} it currently finds that example, irrespective of its class --- an easy majority-class example contributes almost nothing once learned; a hard example, of any class, keeps contributing. Nothing in that factor looks at how many examples of a class exist. Class frequency enters only through the separate $\alpha_t$ term, an explicit per-class multiplier ordinarily set near inverse class frequency.

The implementation here (\texttt{Ablation\_Study/losses.py}) follows this exactly: it computes \texttt{log\_softmax}, gathers $\log p_t$ for the true class, exponentiates to recover $p_t$, and forms \texttt{focal\_weight = (1 - p\_t) ** gamma}. Label smoothing (\autoref{sec:label-smoothing}) is applied first, blending the negative log-likelihood term with a uniform term over classes, so the focal weight multiplies the \textit{smoothed} loss rather than the raw NLL. The optional $\alpha$ vector, where supplied, multiplies the result afterwards, indexed by each example's true class. The value of $\gamma$ used in this study follows Zhao et al.'s reported setting; Section~4.5.1 gives it, and the label-smoothing value, together with the rest of the training configuration.

\subsection{Balanced sampling}
\label{sec:balanced-sampling}

The second lever acts on the data a batch is drawn from rather than on the loss computed over it. A \texttt{WeightedRandomSampler} assigns each training example a weight equal to the inverse of its class's count, computed over the training split of the current fold only. Sampling then proceeds with replacement, drawing \texttt{num\_samples = len(train)} indices per epoch, so each class is presented with roughly equal probability regardless of its true rarity. This sampler is attached to the training loader alone; the validation loader uses no sampler and iterates with \texttt{shuffle=False}, so validation scores reflect the split's true distribution.

\subsection{Why combining both is a choice, not a default}
\label{sec:combining-both}

Focal loss's $\alpha$ term and the balanced sampler both correct for class frequency, at different points in the pipeline, so running both at once corrects for the same imbalance twice. A pipeline that applies both is therefore making a choice, not accepting a default, and the choice has to be resolved somewhere explicit. Section~4.5.2 states how this study resolves it and which of the two corrections is left active. Section~3.9.7 (cross-referenced rather than repeated here) explains why this rule exists and what happens when it is violated; the fact to carry forward is simply that the two mechanisms are not additive by default in this codebase.

\subsection{Label smoothing}
\label{sec:label-smoothing}

Label smoothing replaces the one-hot target with a softened one, mixing a fraction $\epsilon$ of uniform probability mass across all classes into the true-class target. Inside the focal-loss implementation this is a weighted combination of the ordinary negative log-likelihood term and a term averaging $-\log p_c$ over every class $c$; Section~4.5.1 gives the value of $\epsilon$ used here. The effect is to discourage the model from driving the true-class logit arbitrarily high --- regularisation against over-confidence, applied uniformly regardless of class, and distinct from a correction for imbalance. AdamW, mixed precision and gradient clipping have no dedicated treatment in this project's \texttt{docs/} corpus and are used here as standard practice.

\subsection{Augmentation under scarcity}
\label{sec:augmentation}

With few labelled clips per class, augmentation acts before the loss or the sampler is reached at all, and is applied to the training split only. Xia et al.~\cite{xia2020a} motivate this for micro-expression data directly: ``temporal data augmentation strategies as well as a balanced loss are jointly used for our deep network'' to address ``limited and imbalanced training samples''. Two transformations are implemented (\texttt{Ablation\_Study/dataset.py}), but only one of them acts. A random temporal crop of the input window (centred at evaluation) fires only where a clip's stored length exceeds the length the model expects; extraction stores exactly the number of frames the model is configured to consume (\autoref{sec:three-channel-tensor}), so that condition never holds and the crop is inert throughout this study. What remains is a horizontal flip applied with probability one half. Because the input is optical flow rather than raw pixels (\autoref{sec:video-to-motion}), the flip cannot be a plain mirror --- reversing the scene horizontally reverses the sign of horizontal displacement --- so the flip is paired with a negation of the flow tensor's $u$ (horizontal) channel. Flipping without negating $u$ would hand the network motion vectors pointing the wrong way for the geometry shown, a quiet corruption motion-tensor augmentation must specifically guard against.

\subsection{Optimisation}
\label{sec:optimisation}

Training uses AdamW, with decoupled weight decay applied uniformly to every parameter --- no separate no-decay group for biases or normalisation terms. The schedule combines a short linear warmup with cosine annealing: \texttt{LinearLR} ramps the rate from a fraction of target up to full over the warmup epochs, after which \texttt{CosineAnnealingLR} anneals it to a small floor over the remainder; the two are joined by \texttt{SequentialLR} and stepped once per epoch. Section~4.5.4 gives the learning rate, weight decay, warmup length and epoch count used here. Cosine annealing is adopted elsewhere in the micro-expression literature on similar small-data grounds --- Zhang et al.~\cite{zhang2022} use it in a spatio-temporal transformer for this same task.

Gradient norms are clipped after the AMP scaler's gradients are unscaled, via \texttt{clip\_grad\_norm\_}, guarding against the occasional large update a small, skewed dataset can produce. Training runs under mixed precision (\texttt{torch.amp.GradScaler}, \texttt{autocast}), enabled where CUDA is available and not disabled by flag. In the same unscaled branch, and therefore only while clipping is on, every parameter's gradient is checked for non-finite values before the optimiser step; if any is found, the step is skipped entirely, gradients are zeroed, and a warning logged, rather than letting a corrupted update reach the weights. Model selection keeps the checkpoint with the highest validation macro F1, not accuracy, consistent with the metric argument of Section~3.9.5 that accuracy is a poor guide under skew; ties favour the later epoch. There is no early stopping: every configuration trains its full epoch budget regardless of when its best epoch occurred.

One category of technique is absent by design: quantisation, pruning, and knowledge distillation are used nowhere in this project. Efficiency is measured rather than engineered here --- Section~3.6 reports the compute cost of the architectural choices actually made.

Measuring rather than engineering a property requires an evaluation apparatus equal to the task; \autoref{sec:evaluation} builds the one this corpus demands, from the confusion matrix upward.


\section{Evaluating on a Small Corpus}
\label{sec:evaluation}

Chapter~3 argues why leave-one-subject-out evaluation and pooled macro F1 suit this corpus, and what those choices cost here (Sections~3.1.6 and~3.1.8). This section gives the mechanism underneath: the confusion matrix and per-class precision, recall, F1; averaging scores versus pooling counts; the easily conflated distinction between averaging across folds and within them; what cross-validation buys and what leave-one-subject-out protects against; and, to close, an audit of what this thesis's own measurement apparatus does and does not supply.

\subsection{The confusion matrix, and precision, recall, F1 for one class}
\label{sec:confusion-matrix}

For a single class $c$ treated as positive against all others, every prediction falls into one of four counts: true positives $TP_c$ (correctly predicted $c$), false positives $FP_c$ (predicted $c$ but truly something else), false negatives $FN_c$ (truly $c$ but predicted otherwise), and true negatives. Arranging all classes this way at once gives the confusion matrix: entry $(i,j)$ counts clips whose true label is $i$ and predicted label is $j$, so the diagonal holds correct predictions and every off-diagonal cell names a specific error --- which true class is mistaken for which other.

From those counts, precision and recall isolate two failure modes:
\[
	\text{Precision}_c = \frac{TP_c}{TP_c + FP_c}, \qquad \text{Recall}_c = \frac{TP_c}{TP_c + FN_c}.
\]
Precision penalises false alarms: it falls whenever the classifier calls $c$ on a clip that is not $c$, regardless of how many true $c$ instances it also catches. Recall penalises misses: it falls whenever a genuine $c$ instance is assigned elsewhere, regardless of how many false alarms accompany the hits. Neither failure mode is visible from the other, which is why F1, the harmonic mean,
\[
	F1_c = \frac{2 \cdot \text{Precision}_c \cdot \text{Recall}_c}{\text{Precision}_c + \text{Recall}_c} = \frac{2\,TP_c}{2\,TP_c + FP_c + FN_c}
\]
is the usual single-number summary: being harmonic, it stays low whenever either input is low, so a model cannot buy a good F1 by trading recall for precision.

\subsection{Macro versus micro averaging}
\label{sec:macro-micro}

Combining per-class F1 scores into one number can be done two ways. Macro averaging computes $F1_c$ per class and takes an unweighted mean, $\text{Macro-F1} = \tfrac1C\sum_c F1_c$: every class counts equally regardless of size, so ignoring a rare class is penalised as heavily as ignoring a common one. Micro averaging instead pools the raw counts across classes first --- summing every class's $TP$, $FP$ and $FN$ --- and computes one F1 from those totals.

For single-label multi-class classification, where every clip gets exactly one predicted label, this pooling has a consequence worth stating plainly: every misclassification is simultaneously a false positive for the class it was wrongly assigned to and a false negative for its true class, so summed $FP$ and summed $FN$ equal each other and the total error count, and pooled micro-F1 reduces algebraically to $TP_{\text{total}}/N$ --- plain accuracy. Micro-F1 and accuracy are not two different metrics here; they are the same number computed two ways. The contrast that actually matters is therefore not ``macro versus micro'' but \textbf{macro-F1 versus accuracy}: an unweighted per-class average against a count a majority class can dominate.

\subsection{Pooled versus per-fold averaging}
\label{sec:pooled-vs-fold}

A further averaging choice sits underneath macro-F1 once evaluation is split into folds, and it governs how every result in this thesis is reported. One option accumulates the confusion counts --- $TP_c$, $FP_c$, $FN_c$ --- across every fold first, then computes one set of per-class F1 scores, and one macro-F1, from that pooled total. The other computes a complete macro-F1 inside each fold from that fold's own predictions, then averages the per-fold values. See et al.~\cite{see2019} define the Unweighted F1 (UF1) of the MEGC 2019 protocol the first way: the macro-averaged F1 obtained by accumulating true positives, false positives and false negatives over all folds of a leave-one-subject-out run before computing and averaging per-class F1.

These are estimators of two different quantities, not two ways of writing the same one: pooled macro-F1 describes this class's F1 across the whole evaluation, once every held-out prediction sits in one confusion matrix together, while per-fold-averaged macro-F1 describes the typical per-fold macro-F1 --- a quantity that depends on how classes happen to fall within each fold, not only on the classifier. When a fold's held-out set does not contain every class, its per-class F1 for the missing class is degenerate, and averaging such folds in with fully-populated ones changes what the number reflects for reasons unrelated to classifier quality; on a small, unevenly distributed corpus this is a live possibility, and fold composition can bound the per-fold estimator below what the pooled one would report. Sections~3.1.6 and~3.1.8 work through how the fold structure of this study does exactly that; which one a reported number is becomes apparent only once its computation is stated.

\subsection{Cross-validation and subject-disjointness}
\label{sec:cross-validation}

$k$-fold cross-validation partitions the data into $k$ disjoint subsets, trains $k$ models --- each on $k-1$ folds --- and evaluates each on the one fold it never saw, so every example is scored exactly once by a model that never trained on it. Leave-one-subject-out (LOSO) is the case where $k$ equals the number of subjects, and each fold's held-out set is precisely one subject's clips, every other subject's clips forming that fold's training set.

The property this buys is subject-disjointness: no fold's training and validation sets ever share a subject. Without it, a model evaluated on a clip from a subject it has already trained on can succeed by recognising that person's face or recording conditions rather than the expression class itself --- an identity leak inflating the score without the model having learned anything transferable to an unseen person. Enforcing subject-disjointness in every fold closes that channel.

\subsection{What an evaluation apparatus of this kind can fail to provide}
\label{sec:apparatus-gaps}

Three shortcomings belong to an evaluation apparatus rather than to any model it measures, and each matters more at this corpus size than it would at a larger one. They are set out here as properties to look for; Section~4.6 reports where this study's apparatus stands on each.

The first is \textbf{which aggregate a pipeline actually stores}. \autoref{sec:pooled-vs-fold} distinguishes pooling the confusion matrices before computing a macro average from averaging per-fold macro scores, and shows the two answer different questions. A pipeline may compute one, the other, or both under the same field name, and the quantity a reader wants is not necessarily the one written to disk.

The second is \textbf{whether model selection is blind to the data it is scored against}. A nested design keeps three sets apart: one to fit parameters, a second to choose between candidate models or epochs, and a third, untouched until the end, to score the chosen one. Collapsing the second and third --- selecting a checkpoint on the same data that then reports the result --- leaves the reported figure not blind to what it is checked against, and biases it optimistically by an amount that cannot be recovered after the fact.

The third is \textbf{whether enough is retained to support inference}. Aggregate metrics answer what a configuration scored; they cannot answer whether two configurations differ by more than chance. A paired test between configurations needs the per-clip predictions of each, so an apparatus that persists only aggregates forecloses that question whatever its scores say.


\section{Conclusion}
\label{sec:ch2-conclusion}

This chapter has set out the mechanisms on which the rest of the thesis depends, in the order the pipeline applies them. \autoref{sec:phenomenon} defined the micro-expression by the three properties that make it hard to recognise automatically --- it is brief, it is faint, and it cannot be produced on instruction --- and fixed the vocabulary of onset, apex and offset, noting that only the emotion label, not the action-unit coding, is used here. \autoref{sec:recording-to-input} established what actually reaches the pipeline: a 200~fps recording, the corpus's original frames rather than its registered release, and a greyscale frame at a fixed resolution, with no detection, registration or cropping performed by this project. \autoref{sec:motion-magnification} and \autoref{sec:video-to-motion} covered the two transformations applied to those frames --- Eulerian magnification of small motions, and the conversion of the magnified sequence into dense optical flow and optical strain.

\autoref{sec:nn-fundamentals} then built the machinery of a neural network from a single unit up to the convolutional layer, ending with the relationship between parameter count and corpus size that constrains every architectural choice this thesis makes. \autoref{sec:building-blocks} described the specific blocks assembled on that foundation, together with the parameter-free replacement each one falls back to when switched off. \autoref{sec:scarcity-skew} covered training under a class distribution that cannot be balanced by collecting more data, and \autoref{sec:evaluation} the evaluation apparatus needed when a corpus is small enough that the choice of averaging changes the conclusion.

Two threads run through all of it. The first is that the size of the corpus, rather than any modelling preference, forces most of the design: it dictates how few parameters the network can carry, why corrections for class skew are necessary, and why the evaluation protocol has to be chosen with care. The second is that several stages depart from the textbook or reference form of the technique they implement, and each departure has been named at the point it occurs rather than left for a reader to discover.

What this chapter has deliberately not done is say who else has used any of these techniques, with what result, or what remains unresolved. That is the work of Chapter~3, which reviews the literature for each component in turn and identifies the gap this thesis addresses.
